% Le Cameroun dans la coopération bilatérale — ECM Tle Tech — Cours signé OnBuch+
% Programme officiel MINESEC. Compiler avec : tectonic N04-cameroun-cooperation-bilaterale.tex
\newcommand{\DOCMATIERE}{ECM}
\newcommand{\DOCNIVEAU}{Tle Tech}
\newcommand{\DOCMODULE}{ECM – classe de Terminale technique}
\newcommand{\DOCLECON}{N04}
\newcommand{\DOCTITRE}{Le Cameroun dans la coopération bilatérale}
% OnBuch+ — préambule « cours signés » ENRICHI (compilé avec tectonic / XeLaTeX)
% Système visuel « Pop » d'OnBuch+ : fond crème (jamais blanc), bordures encre
% épaisses, ombres « dures » décalées, titres Archivo Black en capitales.
% Version MATHÉMATIQUES : graphes (pgfplots/tikz), boîtes pédagogiques
% (définition, propriété, méthode, attention, le savais-tu, exercice résolu,
% exercice, TP, bilan), page de garde et signature OnBuch+ ancrée en bas.
%
% Macros attendues dans main.tex AVANT \input{preamble} :
%   \DOCMATIERE  \DOCNIVEAU  \DOCMODULE  \DOCLECON (numéro)  \DOCTITRE
\documentclass[11pt]{article}

\usepackage{fontspec}
\usepackage[french]{babel}
\usepackage[a4paper,margin=2cm,top=3.4cm,bottom=2.8cm,headheight=1.4cm,footskip=1.3cm]{geometry}
\usepackage{amsmath,amssymb,amsfonts}
\usepackage{mathtools} % \xmapsto, \coloneqq... souvent utilisés par les modèles
\usepackage{cancel} % simplifications barrées (bilans de forces)
% \abs{z} (module, valeur absolue) et \norm{u} (norme), utilisés par les modèles
\providecommand{\abs}{}\let\abs\relax\DeclarePairedDelimiter{\abs}{\lvert}{\rvert}
\providecommand{\norm}{}\let\norm\relax\DeclarePairedDelimiter{\norm}{\lVert}{\rVert}
\usepackage[table]{xcolor} % [table] : fournit aussi \rowcolors (alternance de lignes)
\usepackage{pagecolor}
\usepackage{tikz}
\usetikzlibrary{arrows.meta,positioning,calc,shapes.geometric,shapes.misc,shapes.arrows,decorations.pathmorphing,decorations.pathreplacing,decorations.markings,patterns,shadows,mindmap,trees,fit,backgrounds,matrix,intersections,angles,quotes}
\usepackage{pgfplots}
\usepackage[version=4]{mhchem} % équations nucléaires \ce{^{235}_{92}U}
\usepackage[european,siunitx]{circuitikz} % schémas électriques
\pgfplotsset{compat=1.18}
\usepgfplotslibrary{fillbetween,groupplots} % aires entre courbes (calcul intégral)
\usepackage{enumitem}
\usepackage{fancyhdr}
% [most] charge toutes les bibliothèques tcolorbox (listings, minted, raster,
% documentation...) et gonfle inutilement la mémoire TeX sur un document long
% (~300 boîtes) : on ne charge que ce qui est réellement utilisé.
\usepackage[skins,breakable]{tcolorbox}
\usepackage{graphicx}
\usepackage{colortbl}
\usepackage{booktabs}
\usepackage{tabularx}
\usepackage{multirow}
\usepackage{diagbox} % case d'en-tête diagonale des tableaux à double entrée
% Colonnes extensibles centrée / gauche / droite (C, L, R), souvent utilisées par les modèles
\newcolumntype{C}{>{\centering\arraybackslash}X}
\newcolumntype{L}{>{\raggedright\arraybackslash}X}
\newcolumntype{R}{>{\raggedleft\arraybackslash}X}
\usepackage{array}
\usepackage{multicol}
\usepackage{siunitx}
% parse-numbers=false : accepte aussi \SI{2,5 \times 10^{-3}}{...}
\sisetup{locale=FR,output-decimal-marker={,},group-minimum-digits=4,parse-numbers=false}
\DeclareSIUnit\ohm{\ensuremath{\symup{\Omega}}} % U+2126 absent de la police
\DeclareSIUnit\division{div} % réglages d'oscilloscope (ms/div, V/div)
\DeclareSIUnit\tr{tr} % tours (vitesse de rotation, ex. \SI{1500}{\tr\per\minute})
\DeclareSIUnit\molar{M} % concentration molaire (ex. \SI{3}{\molar}, \SI{1}{\micro\molar})
\DeclareSIUnit\an{an} % année (le modèle confond parfois avec \year, un compteur TeX) — ex. \SI{5}{\kilo\meter\per\an}
\DeclareSIUnit\FCFA{FCFA} % franc CFA (ex. \SI{500}{\FCFA\per\kilo\gram})
\DeclareSIUnit\degreeBrix{\textdegree Brix} % teneur en sucre d'un sirop/jus (réfractométrie)
\DeclareSIUnit\month{mois} % le modèle utilise parfois \month, un compteur TeX, comme unité de durée
\DeclareSIUnit\microliter{\micro\liter} % le modèle écrit parfois \microliter en un seul token
\DeclareSIUnit\million{millions} % dénombrement (ex. \SI{15}{\million\per\mL} spermatozoïdes)
\usepackage{qrcode}
% \degree : commande fréquemment utilisée par les modèles pour un angle ou une
% température en dehors de \SI{}{} (non fournie par les paquets chargés).
\providecommand{\degree}{^{\circ}}
% \qed (fin de démonstration), utilisé par les modèles sans amsthm
\providecommand{\qed}{\hfill$\square$}
% \barre : double barre des valeurs interdites dans un tableau de variations (tkz-tab)
\providecommand{\barre}{\ensuremath{\|}}
% environnement proof (amsthm non chargé) utilisé par les modèles
\ifdefined\proof\else\newenvironment{proof}[1][Démonstration]{\par\noindent\textit{#1.}\ }{\qed\par}\fi
\usepackage{lastpage}
\usepackage{hyperref}
\hypersetup{hidelinks}

% ============================================================
% Palette « Pop » OnBuch+ — mêmes hex que la classe P (plus_ui.dart)
% ============================================================
\definecolor{poporange}{HTML}{F59321}
\definecolor{popdark}{HTML}{E07A0C}
\definecolor{popcream}{HTML}{FFF6E8}
\definecolor{popink}{HTML}{1C1714}
\definecolor{poppurple}{HTML}{7A5AE0}
\definecolor{popgreen}{HTML}{2E9E6A}
\definecolor{poppink}{HTML}{E0457B}
\definecolor{popblue}{HTML}{2F7DE1}
\definecolor{popgold}{HTML}{C98A12}
\definecolor{popmuted}{HTML}{8A7F76}
\definecolor{popline}{HTML}{EADFD2}
\definecolor{popgray}{HTML}{8A7F76}
\colorlet{popgrey}{popgray}
% Alias tolérants (le modèle nomme parfois les couleurs différemment)
\colorlet{popwhite}{white}
\colorlet{popblack}{popink}
\colorlet{popred}{poppink}
\colorlet{popyellow}{popgold}
% Teintes claires (fonds de tableaux/encadrés) — le modèle nomme parfois
% les couleurs avec un suffixe L (ex. popblueL) sans les avoir définies.
\colorlet{poporangeL}{poporange!20}
\colorlet{popdarkL}{popdark!20}
\colorlet{poppurpleL}{poppurple!20}
\colorlet{popgreenL}{popgreen!20}
\colorlet{poppinkL}{poppink!20}
\colorlet{popblueL}{popblue!20}
\colorlet{popgoldL}{popgold!20}
\colorlet{popredL}{popred!20}
\colorlet{popgrayL}{popgray!20}
% Teintes claires pour fonds de tableaux / zones
\colorlet{popblueL}{popblue!10!white}
\colorlet{poporangeL}{poporange!14!white}
\colorlet{popgreenL}{popgreen!12!white}
\colorlet{poppurpleL}{poppurple!12!white}
\colorlet{poppinkL}{poppink!12!white}
\colorlet{popgoldL}{popgold!14!white}

\pagecolor{popcream}

% ============================================================
% Typographie
% ============================================================
\defaultfontfeatures{Ligatures=TeX}
\setmainfont{PlusJakartaSans-Medium.ttf}[Path=../../../fonts/,BoldFont=PlusJakartaSans-Bold.ttf]
% Maths en sans-serif (Fira Math) avec chiffres et lettres droites de Plus
% Jakarta, pour que les formules se fondent dans le texte.
\usepackage[math-style=ISO,bold-style=ISO]{unicode-math}
\setmathfont{FiraMath-Regular.otf}[Path=../../../fonts/]
\setmathfont{PlusJakartaSans-Medium.ttf}[Path=../../../fonts/,range={up/num,up/{Latin,latin},\mathup}]
\usepackage{icomma} % virgule décimale sans espace en mode math
% Caractères mathématiques Unicode absents des polices de texte (Plus Jakarta,
% Archivo Black) : rendus par leur équivalent mathématique, en texte comme en
% formule — sinon ils s'affichent en carré vide (ex. « Arithmétique dans ℤ »).
\usepackage{newunicodechar}
\newunicodechar{ℝ}{\ensuremath{\mathbb{R}}}
\newunicodechar{ℤ}{\ensuremath{\mathbb{Z}}}
\newunicodechar{ℂ}{\ensuremath{\mathbb{C}}}
\newunicodechar{ℕ}{\ensuremath{\mathbb{N}}}
\newunicodechar{ℚ}{\ensuremath{\mathbb{Q}}}
\newunicodechar{𝔻}{\ensuremath{\mathbb{D}}}
\newunicodechar{α}{\ensuremath{\alpha}}
\newunicodechar{β}{\ensuremath{\beta}}
\newunicodechar{γ}{\ensuremath{\gamma}}
\newunicodechar{δ}{\ensuremath{\delta}}
\newunicodechar{ε}{\ensuremath{\varepsilon}}
\newunicodechar{θ}{\ensuremath{\theta}}
\newunicodechar{λ}{\ensuremath{\lambda}}
\newunicodechar{μ}{\ensuremath{\mu}}
\newunicodechar{σ}{\ensuremath{\sigma}}
\newunicodechar{τ}{\ensuremath{\tau}}
\newunicodechar{φ}{\ensuremath{\varphi}}
\newunicodechar{ψ}{\ensuremath{\psi}}
\newunicodechar{ω}{\ensuremath{\omega}}
\newunicodechar{Σ}{\ensuremath{\Sigma}}
% majuscules produites par \MakeUppercase dans les titres (α-aminés -> Α)
\newunicodechar{Α}{\ensuremath{\alpha}}
\newunicodechar{Β}{\ensuremath{\beta}}
\newunicodechar{Γ}{\ensuremath{\Gamma}}
\newunicodechar{Δ}{\ensuremath{\Delta}}
\newunicodechar{Ε}{\ensuremath{\varepsilon}}
\newunicodechar{Θ}{\ensuremath{\Theta}}
\newunicodechar{Λ}{\ensuremath{\Lambda}}
\newunicodechar{Μ}{\ensuremath{\mu}}
\newunicodechar{Τ}{\ensuremath{\tau}}
\newunicodechar{Φ}{\ensuremath{\Phi}}
\newunicodechar{Ψ}{\ensuremath{\Psi}}
\newunicodechar{Ω}{\ensuremath{\Omega}}
\newunicodechar{↦}{\ensuremath{\mapsto}}
\newunicodechar{∈}{\ensuremath{\in}}
\newunicodechar{∉}{\ensuremath{\notin}}
\newunicodechar{∧}{\ensuremath{\wedge}}
\newunicodechar{⇔}{\ensuremath{\Leftrightarrow}}
\newunicodechar{⟺}{\ensuremath{\Longleftrightarrow}}
\newunicodechar{⇒}{\ensuremath{\Rightarrow}}
\newunicodechar{⟹}{\ensuremath{\Longrightarrow}}
\newunicodechar{∘}{\ensuremath{\circ}}
\newunicodechar{∪}{\ensuremath{\cup}}
\newunicodechar{∩}{\ensuremath{\cap}}
\newunicodechar{≡}{\ensuremath{\equiv}}
\newunicodechar{⊂}{\ensuremath{\subset}}
\newunicodechar{⊥}{\ensuremath{\perp}}
\newunicodechar{∃}{\ensuremath{\exists}}
\newunicodechar{∀}{\ensuremath{\forall}}
\newunicodechar{∛}{\ensuremath{\sqrt[3]{\,}}}
\newunicodechar{∜}{\ensuremath{\sqrt[4]{\,}}}
\newunicodechar{⃗}{\ensuremath{{}^{\to}}}
\newunicodechar{ⁿ}{\ifmmode^{n}\else\textsuperscript{\ensuremath{n}}\fi}
\newunicodechar{ˣ}{\ifmmode^{x}\else\textsuperscript{\ensuremath{x}}\fi}
\newunicodechar{ᵇ}{\ifmmode^{b}\else\textsuperscript{\ensuremath{b}}\fi}
\newunicodechar{ʳ}{\ifmmode^{r}\else\textsuperscript{\ensuremath{r}}\fi}
\newunicodechar{ᵘ}{\ifmmode^{u}\else\textsuperscript{\ensuremath{u}}\fi}
\newunicodechar{ʸ}{\ifmmode^{y}\else\textsuperscript{\ensuremath{y}}\fi}
\newunicodechar{ᵃ}{\ifmmode^{a}\else\textsuperscript{\ensuremath{a}}\fi}
\newunicodechar{ᵉ}{\ifmmode^{e}\else\textsuperscript{\ensuremath{e}}\fi}
\newunicodechar{ᵏ}{\ifmmode^{k}\else\textsuperscript{\ensuremath{k}}\fi}
\newunicodechar{ᵖ}{\ifmmode^{p}\else\textsuperscript{\ensuremath{p}}\fi}
\newunicodechar{ᴺ}{\ifmmode^{N}\else\textsuperscript{\ensuremath{N}}\fi}
\newunicodechar{ᶜ}{\ifmmode^{c}\else\textsuperscript{\ensuremath{c}}\fi}
\newunicodechar{ʰ}{\ifmmode^{h}\else\textsuperscript{\ensuremath{h}}\fi}
\newunicodechar{ᵠ}{\ifmmode^{q}\else\textsuperscript{\ensuremath{q}}\fi}
\newunicodechar{ₙ}{\ifmmode_{n}\else\textsubscript{\ensuremath{n}}\fi}
\newunicodechar{ₐ}{\ifmmode_{a}\else\textsubscript{\ensuremath{a}}\fi}
\newunicodechar{ᵢ}{\ifmmode_{i}\else\textsubscript{\ensuremath{i}}\fi}
\newunicodechar{ᵣ}{\ifmmode_{r}\else\textsubscript{\ensuremath{r}}\fi}
\newunicodechar{ₖ}{\ifmmode_{k}\else\textsubscript{\ensuremath{k}}\fi}
\newunicodechar{ᵦ}{\ifmmode_{\beta}\else\textsubscript{\ensuremath{\beta}}\fi}
\newunicodechar{ₓ}{\ifmmode_{x}\else\textsubscript{\ensuremath{x}}\fi}
\newfontfamily\popdisplay{ArchivoBlack-Regular.ttf}[Path=../../../fonts/]
\newfontfamily\poplogo{SpaceGrotesk-Bold.ttf}[Path=../../../fonts/]
\newfontfamily\popsemi{PlusJakartaSans-SemiBold.ttf}[Path=../../../fonts/]
\newfontfamily\popxbold{PlusJakartaSans-ExtraBold.ttf}[Path=../../../fonts/]
\newfontfamily\popxboldls{PlusJakartaSans-ExtraBold.ttf}[Path=../../../fonts/,LetterSpace=3.0]

\setlist[enumerate]{leftmargin=1.7em,itemsep=5pt,topsep=4pt}
\setlist[itemize]{leftmargin=1.7em,itemsep=4pt,topsep=4pt,label={\color{poporange}\textbullet}}
\setlength{\parindent}{0pt}
\setlength{\parskip}{5pt}
\renewcommand{\arraystretch}{1.25}

% pgfplots : style Pop par défaut
\pgfplotsset{
  every axis/.append style={axis line style={popink,line width=1pt},tick style={popink},
    grid style={popline,line width=0.5pt},label style={font=\small},tick label style={font=\footnotesize}},
  popcurve/.style={poporange,line width=1.8pt,no markers,smooth},
}
% Même style disponible dans un \draw TikZ simple (hors pgfplots)
\tikzset{popcurve/.style={poporange,line width=1.8pt}}
\tikzset{popgrid/.style={popline,very thin}}
\colorlet{popcurve}{poporange} % le nom du style est parfois utilisé comme couleur

% ============================================================
% Pill / squiggle
% ============================================================
\newcommand{\poppill}[3]{%
  \tikz[baseline=-0.55ex]\node[draw=popink,line width=1.2pt,rounded corners=2.6pt,%
    fill=#2,inner xsep=6.5pt,inner ysep=3pt]{%
    \popxboldls\fontsize{7}{7}\selectfont\color{#3}\MakeUppercase{#1}};%
}
\newcommand{\popsquiggle}[1]{%
  \tikz[baseline=-0.4ex]{
    \draw[#1,line width=2.6pt,line cap=round]
      plot [smooth] coordinates {(0,0.09) (0.34,0.01) (0.68,0.21) (1.02,0.01) (1.36,0.10)};
  }%
}
% Mot-clé surligné (marqueur orange sous le texte)
\newcommand{\cle}[1]{{\popxbold\color{popdark}#1}}

% ============================================================
% En-tête / pied
% ============================================================
\pagestyle{fancy}
\fancyhf{}
\renewcommand{\headrulewidth}{0pt}
\renewcommand{\footrulewidth}{0pt}
\lhead{\raisebox{-0.15ex}{\poplogo\fontsize{13}{13}\selectfont\color{popink}OnBuch\color{poporange}+}}
\rhead{\poppill{\DOCMATIERE\ \textbullet\ \DOCNIVEAU\ \textbullet\ Leçon \DOCLECON}{white}{popink}}
\lfoot{\popxbold\fontsize{7.6}{7.6}\selectfont\color{popmuted}\MakeUppercase{Cours signé OnBuch+}\ \textbullet\ {\popsemi\DOCTITRE}}
\rfoot{\tikz[baseline=-0.6ex]\node[rounded corners=8.5pt,fill=popink,minimum height=17pt,minimum width=17pt,inner xsep=5pt,inner sep=0]{\popxbold\fontsize{8}{8}\selectfont\color{popcream}\thepage\,/\,\pageref*{LastPage}};}

% ============================================================
% Titres
% ============================================================
\newcommand{\chaptitre}[2]{%
  \begin{tcolorbox}[enhanced,colback=poporange,colframe=popink,boxrule=2.6pt,arc=11pt,
      drop shadow=popink,left=15pt,right=15pt,top=12pt,bottom=11pt,before skip=3pt,after skip=18pt]
    \poppill{#2}{popink}{popcream}\par\vspace{9pt}
    {\raggedright\hyphenpenalty=10000\exhyphenpenalty=10000\popdisplay\fontsize{22}{24}\selectfont\color{popcream}\MakeUppercase{#1}\par}
  \end{tcolorbox}
}
% Section numérotée : pastille orange avec le numéro + titre Archivo
\newcounter{popsec}
\newcommand{\coursec}[1]{%
  \stepcounter{popsec}\setcounter{popsub}{0}%
  \par\vspace{16pt}\noindent
  \begin{minipage}{\linewidth}
  \tikz[baseline=-0.7ex]\node[draw=popink,line width=1.6pt,fill=poporange,rounded corners=4pt,minimum size=22pt,inner sep=0]{\popdisplay\fontsize{12}{12}\selectfont\color{popcream}\thepopsec};%
  \hspace{8pt}{\popdisplay\fontsize{15}{17}\selectfont\color{popink}\MakeUppercase{#1}}\par
  \vspace{4pt}\noindent{\color{poporange}\rule{36pt}{3.6pt}}
  \end{minipage}\par\nopagebreak\vspace{8pt}
}
\newcounter{popsub}
\newcommand{\courssub}[1]{%
  \stepcounter{popsub}%
  \par\vspace{9pt}\noindent{\popxbold\fontsize{11.5}{13}\selectfont\color{popink}{\color{poporange}\thepopsec.\thepopsub}\hspace{6pt}#1}\par\nopagebreak\vspace{4pt}
}

% ============================================================
% Boîtes pédagogiques — même recette : fond blanc, bordure épaisse colorée,
% ombre dure encre, badge plein. Titre optionnel [#1].
% ============================================================
\tcbset{popbase/.style={breakable,enhanced,colback=white,boxrule=2.3pt,arc=9pt,
  shadow={3pt}{-3pt}{0pt}{popink},left=14pt,right=14pt,top=11pt,bottom=11pt,before skip=11pt,after skip=15pt,
  fontupper=\popsemi\fontsize{10.6}{14.5}\selectfont\color{popink},
  fontlower=\popsemi\fontsize{10.6}{14.5}\selectfont\color{popink}}}
% Coche dessinée (le glyphe ✓ n'existe pas dans les polices du document)
\newcommand{\popcheck}[1]{\tikz[baseline=-0.5ex]\draw[#1,line width=1.6pt,line cap=round,line join=round](-0.13,0.0)--(-0.04,-0.09)--(0.13,0.1);}
\providecommand{\coche}{\popcheck{popgreen}}
\providecommand{\resultat}[1]{\par\smallskip\noindent\fcolorbox{popgreen}{popgreenL}{\parbox{\dimexpr\linewidth-2\fboxsep-2\fboxrule\relax}{\textbf{Résultat.} #1}}\par\smallskip}
\newcommand{\popboxhead}[3]{\poppill{#1}{#2}{white}\ifx\relax#3\relax\else\hspace{6pt}{\popxbold\fontsize{10.6}{12}\selectfont\color{popink}#3}\fi\par\vspace{7pt}}

\newtcolorbox{aretenir}[1][]{popbase,colframe=popgreen,before upper={\popboxhead{À retenir}{popgreen}{#1}}}
\newtcolorbox{exemplebox}[1][]{popbase,colframe=poporange,before upper={\popboxhead{Exemple}{poporange}{#1}}}
\newtcolorbox{definition}[1][]{popbase,colframe=popblue,before upper={\popboxhead{Définition}{popblue}{#1}}}
\newtcolorbox{propriete}[1][]{popbase,colframe=poppurple,before upper={\popboxhead{Propriété}{poppurple}{#1}}}
\newtcolorbox{methode}[1][]{popbase,colframe=popgold,before upper={\popboxhead{Méthode}{popgold}{#1}}}
\newtcolorbox{attention}[1][]{popbase,colframe=poppink,before upper={\popboxhead{Attention}{poppink}{#1}}}
\newtcolorbox{experience}[1][]{popbase,colframe=popblue,colback=popblueL,before upper={\popboxhead{Expérience}{popblue}{#1}}}
% « Le savais-tu ? » : carte violette pleine (contexte, culture, Cameroun)
\newtcolorbox{savaistu}[1][]{popbase,colback=poppurple,colframe=popink,
  fontupper=\popsemi\fontsize{10.4}{14.2}\selectfont\color{white},
  before upper={\poppill{Le savais-tu\,?}{white}{poppurple}\ifx\relax#1\relax\else\hspace{6pt}{\popxbold\color{white}#1}\fi\par\vspace{7pt}}}
% Exercice résolu : énoncé en haut, \tcblower puis la solution sur fond crème
\newcounter{popexo}\newcounter{popexr}
\newtcolorbox{exoresolu}[1][]{popbase,colframe=poporange,colbacklower=popcream,
  segmentation style={popink,line width=1.2pt,dashed},
  before upper={\stepcounter{popexr}\popboxhead{Exercice résolu \thepopexr}{poporange}{#1}},
  before lower={\poppill{Solution}{popink}{popcream}\par\vspace{6pt}}}
% Exercice (non résolu en place) — la correction est dans la partie Corrigés
\newtcolorbox{exercice}[1][]{popbase,colframe=popink,
  before upper={\stepcounter{popexo}\popboxhead{Exercice \thepopexo}{popink}{#1}}}
% Corrigé
\newtcolorbox{corrige}[1][]{popbase,colframe=popgreen,colback=popgreenL,
  before upper={\popboxhead{Corrigé}{popgreen}{#1}}}
% Objectifs / prérequis / bilan
\newtcolorbox{objectifs}{popbase,colframe=popink,colback=poporangeL,
  before upper={\popboxhead{Objectifs de la leçon}{popink}{}}}
\newtcolorbox{prerequis}{popbase,colframe=popink,colback=popblueL,
  before upper={\popboxhead{Prérequis}{popblue}{}}}
\newtcolorbox{bilan}{popbase,colframe=popink,colback=white,boxrule=2.8pt,
  before upper={\popboxhead{Fiche bilan}{poporange}{L'essentiel à connaître pour l'examen}}}
% Figure encadrée avec légende
\newtcolorbox{popfigure}[1][]{enhanced,breakable=false,colback=white,colframe=popline,boxrule=1.4pt,arc=8pt,
  left=10pt,right=10pt,top=10pt,bottom=8pt,before skip=10pt,after skip=12pt,halign=center,
  lower separated=false,halign lower=center,
  fontlower=\popsemi\fontsize{9}{11.5}\selectfont\color{popmuted},
  #1}
\newcommand{\legende}[1]{\par\vspace{6pt}{\popsemi\fontsize{9}{11.5}\selectfont\color{popmuted}#1}}

% ============================================================
% Signature OnBuch+
% ============================================================
\newcommand{\poplogoinline}[1]{{\poplogo\fontsize{#1}{#1}\selectfont\color{popink}OnBuch\color{poporange}+}}
\newcommand{\signatureonbuch}{%
  \begin{tcolorbox}[enhanced,colback=popink,colframe=popink,boxrule=2.2pt,arc=10pt,
      drop shadow=poporange,left=14pt,right=14pt,top=10pt,bottom=10pt,before skip=0pt,after skip=0pt]
    \tikz[baseline=-0.7ex]\node[circle,fill=popgreen,draw=popcream,line width=1.3pt,minimum size=22pt,inner sep=0]{\popcheck{white}};%
    \hspace{9pt}\begin{minipage}[c]{0.62\linewidth}
      {\popxbold\fontsize{12}{13}\selectfont\color{popcream}Signé\ \textbullet\ {\poplogo OnBuch\color{poporange}+}}\par\vspace{2pt}
      {\raggedright\popsemi\fontsize{8}{10.5}\selectfont\color{popline}\MakeUppercase{Équipe pédagogique OnBuch+}\\\MakeUppercase{Conforme au programme officiel MINESEC}\par}
    \end{minipage}\hfill
    {\poplogo\fontsize{22}{22}\selectfont\color{popcream}OnBuch\color{poporange}+}
  \end{tcolorbox}%
}

% ============================================================
% Page de garde
% #1 titre · #2 matière · #3 niveau · #4 module · #5 n° leçon · #6 durée ·
% #7 description / objectifs (texte ou itemize)
% ============================================================
\newcommand{\pagegarde}[7]{%
  \thispagestyle{empty}
  \newgeometry{a4paper,margin=1.7cm,top=1.6cm,bottom=1.4cm}%
  \begin{tikzpicture}[remember picture,overlay]
    \fill[poporange!18] ([xshift=-3.2cm,yshift=-3.6cm]current page.north east) circle (4.6cm);
    \fill[poppurple!14] ([xshift=2.4cm,yshift=7.6cm]current page.south west) circle (3.8cm);
  \end{tikzpicture}%
  {\poplogo\fontsize{30}{30}\selectfont\color{popink}OnBuch\color{poporange}+}\hfill
  \raisebox{0.6ex}{\poppill{Cours signé \textbullet\ Premium}{popink}{popcream}}\par
  \vspace{1pt}\popsquiggle{poppurple}\par
  \vspace{8pt}
  \begin{tcolorbox}[enhanced,colback=poporange,colframe=popink,boxrule=2.8pt,arc=13pt,
      drop shadow=popink,left=18pt,right=18pt,top=12pt,bottom=13pt,after skip=10pt]
    \poppill{#2 \textbullet\ #3}{popink}{popcream}\hspace{5pt}\poppill{#4}{white}{popink}\par\vspace{8pt}
    {\popdisplay\fontsize{14}{14}\selectfont\color{popink}LEÇON #5}\par\vspace{4pt}
    {\raggedright\hyphenpenalty=10000\exhyphenpenalty=10000\popdisplay\fontsize{25}{27}\selectfont\color{popcream}\MakeUppercase{#1}\par}
    \vspace{8pt}
    {\popxbold\fontsize{9.5}{9.5}\selectfont\color{popink}\MakeUppercase{Durée conseillée : #6}}
  \end{tcolorbox}
  \begin{tcolorbox}[enhanced,colback=white,colframe=popink,boxrule=2.2pt,arc=10pt,
      drop shadow=popink,left=16pt,right=16pt,top=10pt,bottom=10pt,after skip=8pt]
    \poppill{Dans cette leçon}{poppurple}{white}\par\vspace{5pt}
    \popsemi\fontsize{9.3}{12.6}\selectfont\color{popink}#7
  \end{tcolorbox}
  \noindent\begin{minipage}[t]{0.487\linewidth}
    \begin{tcolorbox}[enhanced,colback=popgreen,colframe=popink,boxrule=2.2pt,arc=10pt,
        drop shadow=popink,left=11pt,right=11pt,top=10pt,bottom=10pt,height=3.1cm,valign=center]
      \tcbox[colback=white,colframe=popink,boxrule=1.3pt,arc=5pt,boxsep=0pt,nobeforeafter,
        left=3pt,right=3pt,top=3pt,bottom=3pt,valign=center]{\qrcode[height=1.9cm]{https://play.google.com/store/apps/details?id=com.luvvix.onbuch&hl=fr}}%
      \hspace{8pt}\begin{minipage}[c]{0.5\linewidth}
        {\popdisplay\fontsize{11}{12.5}\selectfont\color{white}\MakeUppercase{Télécharge OnBuch}}\par\vspace{4pt}
        {\raggedright\popsemi\fontsize{8.4}{11}\selectfont\color{white}Gratuit \textbullet\ Google Play\\Tuteur IA, annales, concours, résultats\par}
      \end{minipage}
    \end{tcolorbox}
  \end{minipage}\hfill
  \begin{minipage}[t]{0.487\linewidth}
    \begin{tcolorbox}[enhanced,colback=poppurple,colframe=popink,boxrule=2.2pt,arc=10pt,
        drop shadow=popink,left=11pt,right=11pt,top=10pt,bottom=10pt,height=3.1cm,valign=center]
      \tikz[baseline=-0.7ex]\node[circle,fill=white,draw=popink,line width=1.3pt,minimum size=1.6cm,inner sep=0]{\popdisplay\fontsize{24}{24}\selectfont\color{poppurple}+};%
      \hspace{8pt}\begin{minipage}[c]{0.6\linewidth}
        {\popdisplay\fontsize{11}{12.5}\selectfont\color{white}\MakeUppercase{Passe à OnBuch+}}\par\vspace{4pt}
        {\raggedright\popsemi\fontsize{8.4}{11}\selectfont\color{white}Tous les cours signés, Tuteur IA illimité, fiches PDF — onglet OnBuch+ de l'app\par}
      \end{minipage}
    \end{tcolorbox}
  \end{minipage}
  \vspace{8pt}
  \popcontenu
  \vfill
  \signatureonbuch
  \restoregeometry
  \newpage
}

% Rangée « Ce cours contient » (page de garde)
\newcommand{\poptile}[3]{% #1 couleur #2 gros texte #3 légende
  \begin{tcolorbox}[enhanced,colback=#1,colframe=popink,boxrule=2pt,arc=9pt,shadow={2.5pt}{-2.5pt}{0pt}{popink},
      left=6pt,right=6pt,top=9pt,bottom=9pt,halign=center,valign=center,height=1.75cm,width=\linewidth,nobeforeafter]
    {\popdisplay\fontsize{12.5}{13}\selectfont\color{white}\MakeUppercase{#2}}\par\vspace{3pt}
    {\centering\popsemi\fontsize{7.8}{9.5}\selectfont\color{white}#3\par}
  \end{tcolorbox}}
\newcommand{\popcontenu}{%
  \par\noindent{\popxboldls\fontsize{7.5}{7.5}\selectfont\color{popmuted}CE COURS SIGNÉ CONTIENT}\par\vspace{7pt}
  \noindent\begin{minipage}[t]{0.235\linewidth}\poptile{popblue}{Cours}{complet, illustré, pas à pas}\end{minipage}\hfill
  \begin{minipage}[t]{0.235\linewidth}\poptile{poporange}{Méthodes}{et exercices résolus}\end{minipage}\hfill
  \begin{minipage}[t]{0.235\linewidth}\poptile{poppink}{Exercices}{type examen + corrigés}\end{minipage}\hfill
  \begin{minipage}[t]{0.235\linewidth}\poptile{popgreen}{Fiche bilan}{l'essentiel à retenir}\end{minipage}\par}

% Sommaire
\newtcolorbox{sommaire}{enhanced,colback=white,colframe=popink,boxrule=2.2pt,arc=10pt,
  drop shadow=popink,left=18pt,right=18pt,top=13pt,bottom=13pt,before skip=6pt,after skip=20pt,
  before upper={\poppill{Sommaire}{poppurple}{white}\par\vspace{9pt}\popsemi\fontsize{10.6}{14.5}\selectfont\color{popink}}}

% Signature de fin de cours — poussée tout en bas de la dernière page
\newcommand{\findecours}{%
  \par\vfill\nopagebreak
  \begin{center}\popsquiggle{poporange}\end{center}\vspace{4pt}
  \signatureonbuch
}
\AtBeginDocument{\def\llbracket{\mathopen{[\mkern-3.5mu[}}\def\rrbracket{\mathclose{]\mkern-3.5mu]}}} % intervalles d'entiers (glyphe ⟦ absent de la police)
% Glyphes absents de Fira Math : redéfinis à partir de glyphes disponibles
\AtBeginDocument{%
  \def\setminus{\mathbin{\backslash}}\let\smallsetminus\setminus
  \def\vdots{\mathord{\vbox{\baselineskip3.2pt\lineskiplimit0pt\kern4pt\hbox{.}\hbox{.}\hbox{.}}}}%
  \def\ddots{\mathinner{\mkern1mu\raise6.5pt\hbox{.}\mkern2mu\raise3.6pt\hbox{.}\mkern2mu\raise0.7pt\hbox{.}\mkern1mu}}%
  \def\triangle{\mathord{\scalebox{0.9}{$\symup{\Delta}$}}}%
  \def\nabla{\mathord{\raisebox{\depth}{\scalebox{1}[-1]{$\symup{\Delta}$}}}}%
  \def\checkmark{\popcheck{popdark}}%
  \def\star{\mathbin{\ast}}\def\bigstar{\mathord{\ast}}%
  \def\diamond{\mathbin{\scalebox{0.6}{$\Diamond$}}}%
  \def\bigcup{\mathop{\mathchoice{\raisebox{-0.3em}{\scalebox{1.7}{$\cup$}}}{\scalebox{1.25}{$\cup$}}{\cup}{\cup}}\displaylimits}%
  \def\bigcap{\mathop{\mathchoice{\raisebox{-0.3em}{\scalebox{1.7}{$\cap$}}}{\scalebox{1.25}{$\cap$}}{\cap}{\cap}}\displaylimits}%
  \def\lesseqqgtr{\mathrel{\lessgtr}}\def\bigoplus{\mathop{\oplus}\displaylimits}%
}
\providecommand{\ssi}{si et seulement si} % abréviation fréquente
\AtBeginDocument{\providecommand{\wideparen}{\overparen}} % arc de cercle

%% FIGFIX-BEGIN v4 — correctifs globaux des figures (docs/onbuch-generation/tools/figfix)
\usepackage{adjustbox}\usepackage{etoolbox}
% toute figure TikZ/pgfplots plus large que la ligne est réduite à la largeur disponible
\BeforeBeginEnvironment{tikzpicture}{\begin{adjustbox}{max width=\linewidth}}
\AfterEndEnvironment{tikzpicture}{\end{adjustbox}}
% légendes pgfplots lisibles par-dessus les courbes
\pgfplotsset{every axis/.append style={legend style={fill=white,fill opacity=0.92,text opacity=1,draw=black!15,inner sep=3pt}}}
% étiquettes d'arêtes (syntaxe edge["..."]) avec fond blanc
\tikzset{every edge quotes/.append style={fill=white,inner sep=1.2pt}}
%% FIGFIX-END
\begin{document}

\pagegarde{Le Cameroun dans la coopération bilatérale}{ECM}{Tle Tech}{ECM – classe de Terminale technique}{N04}{3 séances (fiche de progression harmonisée)}{Cette leçon étudie les relations bilatérales du Cameroun avec ses principaux partenaires : la France et la Grande-Bretagne (partenaires historiques), la Chine et le Japon (partenaires asiatiques montants), les USA et le Brésil (partenaires américains). Elle montre comment ces coopérations structurent le développement économique, culturel et diplomatique du Cameroun.
\vspace{4pt}
\begin{multicols}{2}\small
\begin{itemize}
\item À la fin de cette leçon, je sais définir la coopération bilatérale et distinguer ses formes (économique, culturelle, militaire, technique)
\item À la fin de cette leçon, je sais décrire les domaines et les réalisations concrètes de la coopération Cameroun-France
\item À la fin de cette leçon, je sais expliquer les spécificités de la coopération Cameroun-Grande-Bretagne liées à l'héritage anglophone
\item À la fin de cette leçon, je sais identifier les grands projets de la coopération sino-camerounaise et nippo-camerounaise
\item À la fin de cette leçon, je sais présenter les axes de la coopération du Cameroun avec les USA et le Brésil
\item À la fin de cette leçon, je sais lire et interpréter des documents (cartes, données chiffrées, photographies) sur les échanges du Cameroun
\end{itemize}\end{multicols}}

\begin{sommaire}
\begin{multicols}{2}
\begin{enumerate}
\item Généralités sur la coopération bilatérale du Cameroun
\item Le Cameroun et la France
\item Le Cameroun et la Grande-Bretagne
\item Le Cameroun et les pays asiatiques : la Chine et le Japon
\item Le Cameroun et les pays d'Amérique : USA et Brésil
\item Bilan : avantages, limites et enjeux de la coopération bilatérale
\item Activité d'intégration
\item Méthodes et exercices résolus
\item Exercices
\item Corrigés des exercices
\item Fiche bilan
\end{enumerate}
\end{multicols}
\end{sommaire}

\begin{objectifs}
\begin{itemize}
\item À la fin de cette leçon, je sais définir la coopération bilatérale et distinguer ses formes (économique, culturelle, militaire, technique)
\item À la fin de cette leçon, je sais décrire les domaines et les réalisations concrètes de la coopération Cameroun-France
\item À la fin de cette leçon, je sais expliquer les spécificités de la coopération Cameroun-Grande-Bretagne liées à l'héritage anglophone
\item À la fin de cette leçon, je sais identifier les grands projets de la coopération sino-camerounaise et nippo-camerounaise
\item À la fin de cette leçon, je sais présenter les axes de la coopération du Cameroun avec les USA et le Brésil
\item À la fin de cette leçon, je sais lire et interpréter des documents (cartes, données chiffrées, photographies) sur les échanges du Cameroun
\item À la fin de cette leçon, je sais construire un croquis ou un diagramme illustrant les partenaires bilatéraux du Cameroun
\item À la fin de cette leçon, je sais rédiger et justifier une argumentation sur les avantages et les limites de la coopération bilatérale
\item À la fin de cette leçon, je sais appliquer ces notions à des situations concrètes de la vie courante (emploi, infrastructures, bourses d'études)
\end{itemize}
\end{objectifs}

\begin{prerequis}
\begin{itemize}
\item Les grandes étapes de l'histoire du Cameroun : colonisation allemande, mandats français et britannique, indépendances (cours de 4e/3e)
\item La notion de coopération et les organisations internationales vues en début d'année de Terminale
\item La localisation des continents et des grandes puissances mondiales sur une carte du monde
\item Les notions d'aide au développement, d'investissement et d'échanges commerciaux
\item La lecture de documents géographiques simples (cartes, tableaux de données)
\end{itemize}
\end{prerequis}

\newpage

\coursec{Généralités sur la coopération bilatérale du Cameroun}

\textbf{Situation d'accroche.} À Douala, le port en eau profonde de Kribi a été construit avec des financements chinois ; à Yaoundé, le Palais des Sports porte la signature de la coopération chinoise, tandis que le lycée bilingue fonctionne grâce à l'héritage anglo-français. Un élève se demande : pourquoi le Cameroun entretient-il des relations privilégiées avec la France, la Grande-Bretagne, la Chine, le Japon, les USA ou le Brésil ? Que gagne-t-il concrètement de ces partenariats ? Et que donnent ces pays en échange ? Cette leçon répond à ces questions en analysant la \cle{coopération bilatérale} du Cameroun.

\textbf{Problématique.} Qu'est-ce que la coopération bilatérale ? Quelles formes prend-elle, avec quels instruments, quels partenaires et quels enjeux pour le Cameroun ?

\courssub{Définition de la coopération bilatérale}

\textbf{Intuition.} Quand deux voisins s'entraident — l'un prête sa brouette, l'autre partage sa récolte —, chacun y gagne. À l'échelle des États, c'est la même logique : deux pays décident librement d'échanger des biens, des savoirs, des financements ou une assistance, parce que chacun y trouve un intérêt.

\begin{definition}[Coopération bilatérale]
La \cle{coopération bilatérale} est l'ensemble des relations d'échanges et d'entraide établies entre deux \cle{États souverains} dans des domaines variés : économie, commerce, finances, culture, éducation, science, défense et technique. Elle repose sur des accords signés librement entre les deux partenaires, dans le respect de leur souveraineté et de leurs intérêts réciproques.
\end{definition}

Le mot « bilatéral » vient du latin \emph{bi} (deux) et \emph{latus} (côté) : il s'agit donc littéralement d'une relation « à deux côtés ». Le Cameroun, État indépendant depuis 1960, a ainsi noué des liens privilégiés avec de nombreux pays du monde.

\begin{attention}[Ne pas confondre]
Il ne faut pas confondre la \textbf{coopération bilatérale} (entre deux pays seulement, par exemple le Cameroun et la Chine) et la \textbf{coopération multilatérale} (au sein d'une organisation regroupant plusieurs pays, comme l'ONU, l'Union africaine — UA — ou la CEEAC). Au Bac, cette distinction est fréquemment demandée : cite toujours un exemple de chaque type pour montrer que tu maîtrises la différence.
\end{attention}

\courssub{Les formes de la coopération bilatérale}

La coopération entre le Cameroun et ses partenaires se déploie dans plusieurs domaines complémentaires.

\textbf{1. La coopération économique et commerciale.} Elle concerne les échanges de marchandises (exportations de pétrole, de cacao, de bois, de coton ; importations de machines, de véhicules, de produits manufacturés), les investissements des entreprises étrangères et la construction d'infrastructures (routes, ports, barrages).

\textbf{2. La coopération financière.} Elle prend la forme d'\cle{aides}, de \cle{dons} (argent ou matériel offerts sans contrepartie de remboursement) et de \cle{prêts} (argent prêté, souvent à taux réduit, que le Cameroun doit rembourser). Elle finance les grands projets de développement.

\textbf{3. La coopération culturelle et scientifique.} Elle comprend les bourses d'études, l'envoi d'enseignants et de chercheurs, l'appui aux écoles et universités, la promotion des langues (français, anglais, chinois...) et les échanges artistiques.

\textbf{4. La coopération militaire et technique.} Elle porte sur la formation des forces de défense et de sécurité, la fourniture d'équipements, l'envoi d'experts et de coopérants techniques dans l'administration, la santé ou l'agriculture.

\begin{popfigure}
\begin{center}
\begin{tikzpicture}[scale=1]
\node[circle, draw=popdark, fill=popblueL, minimum size=2.6cm, align=center, font=\bfseries] (cam) at (0,0) {CAMEROUN};
\node[draw=poporange, fill=poporangeL, rounded corners, align=center, font=\small] (eco) at (90:3.4) {Économique\\ et commerciale};
\node[draw=popgreen, fill=popgreenL, rounded corners, align=center, font=\small] (fin) at (0:4.1) {Financière\\ (aides, dons, prêts)};
\node[draw=poppurple, fill=poppurpleL, rounded corners, align=center, font=\small] (cul) at (270:3.4) {Culturelle\\ et scientifique};
\node[draw=popgold, fill=popgoldL, rounded corners, align=center, font=\small] (mil) at (180:4.1) {Militaire\\ et technique};
\draw[-{Stealth}, very thick, poporange] (cam) -- (eco);
\draw[-{Stealth}, very thick, popgreen] (cam) -- (fin);
\draw[-{Stealth}, very thick, poppurple] (cam) -- (cul);
\draw[-{Stealth}, very thick, popgold] (cam) -- (mil);
\end{tikzpicture}
\end{center}
\legende{Les quatre formes de la coopération bilatérale autour du Cameroun. Chaque forme répond à un besoin précis : échanger, financer, former, sécuriser.}
\end{popfigure}

\courssub{Les instruments de la coopération bilatérale}

Pour exister juridiquement et fonctionner au quotidien, la coopération bilatérale s'appuie sur des instruments précis.

\begin{itemize}
\item Les \cle{accords de coopération} : documents signés par les deux États qui définissent les domaines et les modalités de leur collaboration (par exemple un accord de coopération économique ou culturelle).
\item Les \cle{conventions} : textes plus précis qui organisent un secteur particulier (convention fiscale, convention de stage, convention universitaire).
\item Les \cle{traités} : engagements solennels et durables entre les deux pays (traité d'amitié, traité de défense).
\item Les \cle{ambassades} et les \cle{consulats} : représentations diplomatiques permanentes. L'ambassade, dirigée par un ambassadeur, entretient les relations politiques ; le consulat s'occupe des ressortissants (visas, état civil, assistance).
\end{itemize}

\begin{savaistu}[Le Cameroun, acteur diplomatique]
Le Cameroun entretient des relations diplomatiques avec plus de 100 pays et abrite à Yaoundé de nombreuses ambassades, ce qui en fait un acteur diplomatique notable en Afrique centrale. Sa position géographique, sa stabilité relative et son bilinguisme français-anglais en font un partenaire recherché.
\end{savaistu}

\courssub{Les principaux partenaires bilatéraux du Cameroun}

Le Cameroun diversifie ses partenaires pour ne dépendre d'aucun pays. On peut les présenter en trois grands ensembles.

\begin{itemize}
\item \textbf{Les partenaires européens historiques} : la \cle{France} (ancienne puissance mandataire, premier partenaire économique et culturel), la \cle{Grande-Bretagne} (liée à l'héritage anglophone) et l'\cle{Allemagne} (ancienne puissance coloniale, active dans les projets techniques).
\item \textbf{Les partenaires asiatiques} : la \cle{Chine} (grands travaux : routes, stades, port de Kribi, barrages) et le \cle{Japon} (aide technique, santé, équipements, dons).
\item \textbf{Les partenaires américains} : les \cle{USA} (formation, sécurité, santé, commerce via la loi AGOA) et le \cle{Brésil} (agriculture, élevage, formation professionnelle, affinités culturelles).
\end{itemize}

\begin{popfigure}
\begin{center}
\begin{tikzpicture}[scale=0.92]
% Fond océan
\fill[popblueL!50] (-7.5,-3.2) rectangle (7.5,3.2);
% Amérique du Nord (bloc simplifié)
\fill[popgreenL] (-6.8,0.6) -- (-4.6,1.0) -- (-4.2,2.6) -- (-6.9,2.8) -- cycle;
% Amérique du Sud
\fill[popgreenL] (-5.2,0.4) -- (-3.9,0.4) -- (-4.2,-2.6) -- (-5.0,-1.6) -- cycle;
% Europe
\fill[poporangeL] (-0.9,1.2) -- (1.4,1.2) -- (1.6,2.7) -- (-1.1,2.7) -- cycle;
% Afrique
\fill[popgoldL] (-0.7,1.0) -- (1.5,1.0) -- (1.2,-2.6) -- (0.1,-2.9) -- (-0.9,-1.2) -- cycle;
% Asie
\fill[poppurpleL] (1.9,0.5) -- (6.9,0.7) -- (7.1,2.8) -- (2.0,2.8) -- cycle;
% Japon (petit bloc)
\fill[poppurpleL] (6.6,0.0) rectangle (7.2,0.45);
% Points
\node[circle, fill=popdark, inner sep=1.6pt, label={[font=\scriptsize\bfseries]below:Cameroun}] (cam) at (0.3,-0.4) {};
\node[circle, fill=poporange, inner sep=1.4pt, label={[font=\scriptsize]above:France / G.-B. / All.}] (fr) at (0.2,2.0) {};
\node[circle, fill=poppurple, inner sep=1.4pt, label={[font=\scriptsize]above:Chine}] (cn) at (4.6,1.7) {};
\node[circle, fill=poppurple, inner sep=1.4pt, label={[font=\scriptsize]right:Japon}] (jp) at (6.9,0.2) {};
\node[circle, fill=popgreen, inner sep=1.4pt, label={[font=\scriptsize]above:USA}] (us) at (-5.6,1.8) {};
\node[circle, fill=popgreen, inner sep=1.4pt, label={[font=\scriptsize]below:Brésil}] (br) at (-4.5,-1.3) {};
% Flèches d'échanges
\draw[{Stealth}-{Stealth}, very thick, poporange] (cam) -- (fr);
\draw[{Stealth}-{Stealth}, very thick, poppurple] (cam) -- (cn);
\draw[{Stealth}-{Stealth}, very thick, poppurple] (cam) -- (jp);
\draw[{Stealth}-{Stealth}, very thick, popgreen] (cam) -- (us);
\draw[{Stealth}-{Stealth}, very thick, popgreen] (cam) -- (br);
% Nord
\draw[-{Stealth}, thick, popdark] (-7.1,-2.6) -- (-7.1,-1.6) node[above, font=\scriptsize\bfseries]{N};
\end{tikzpicture}
\end{center}
\legende{Planisphère très simplifié localisant les principaux partenaires bilatéraux du Cameroun (flèches = flux d'échanges : biens, financements, formations). Carte schématique, sans précision cartographique.}
\end{popfigure}

\courssub{Les enjeux de la coopération bilatérale pour le Cameroun}

Pourquoi le Cameroun recherche-t-il ces partenariats ? Quatre enjeux majeurs se dégagent.

\begin{enumerate}
\item \textbf{Le financement du développement} : les dons et les prêts permettent de construire les infrastructures (ports, routes, barrages, hôpitaux) que le budget national ne peut financer seul.
\item \textbf{Le transfert de technologies} : les entreprises et experts étrangers apportent des savoir-faire (génie civil, télécommunications, agriculture) que les techniciens camerounais apprennent et réutilisent.
\item \textbf{La formation} : bourses d'études, stages, envoi d'enseignants et coopérants renforcent les compétences des élèves, étudiants et fonctionnaires camerounais.
\item \textbf{L'ouverture diplomatique} : multiplier les partenaires donne au Cameroun une voix sur la scène internationale, attire les investisseurs et évite la dépendance envers un seul pays.
\end{enumerate}

\begin{methode}[Analyser une relation bilatérale]
Pour analyser une relation bilatérale (méthode exigible au Bac), procède en quatre étapes :
\begin{enumerate}
\item Identifier le \textbf{domaine} de coopération (économique, financier, culturel, militaire...).
\item Citer les \textbf{réalisations concrètes} (infrastructures, bourses, accords signés).
\item Dégager les \textbf{avantages pour chaque partenaire} (que gagne le Cameroun ? que gagne l'autre pays ?).
\item Relever les \textbf{limites éventuelles} (endettement, dépendance, contreparties exigées).
\end{enumerate}
\end{methode}

\begin{exoresolu}[Distinguer coopération bilatérale et multilatérale]
\textbf{Énoncé.} Voici quatre situations : (a) la Chine construit le Palais des Sports de Yaoundé ; (b) l'ONU envoie des casques bleus en Afrique centrale ; (c) la France accorde des bourses à des étudiants camerounais ; (d) la CEEAC organise un sommet des chefs d'État. Classe chaque situation selon qu'elle relève de la coopération bilatérale ou multilatérale, en justifiant.
\tcblower
\textbf{Solution détaillée.}
\begin{itemize}
\item (a) \textbf{Bilatérale} : elle engage deux États seulement, le Cameroun et la Chine (coopération financière et technique).
\item (b) \textbf{Multilatérale} : l'ONU est une organisation regroupant près de 193 États membres.
\item (c) \textbf{Bilatérale} : relation directe entre la France et le Cameroun (coopération culturelle et scientifique).
\item (d) \textbf{Multilatérale} : la CEEAC réunit onze pays d'Afrique centrale.
\end{itemize}
\textbf{Conclusion} : une coopération est bilatérale dès lors qu'elle lie deux États souverains ; elle est multilatérale dès lors qu'elle passe par une organisation internationale regroupant plusieurs pays.
\end{exoresolu}

\begin{aretenir}[L'essentiel de la section]
La coopération bilatérale est l'ensemble des relations d'échanges et d'entraide entre deux États souverains. Elle prend quatre formes (économique et commerciale, financière, culturelle et scientifique, militaire et technique), s'appuie sur des instruments (accords, conventions, traités, ambassades et consulats) et répond pour le Cameroun à quatre enjeux : financer le développement, transférer les technologies, former les hommes et s'ouvrir diplomatiquement. Ses principaux partenaires sont la France, la Grande-Bretagne, l'Allemagne, la Chine, le Japon, les USA et le Brésil.
\end{aretenir}

\begin{exercice}[À toi de jouer]
\begin{enumerate}
\item Définis la coopération bilatérale et donne deux exemples de ses instruments.
\item Cite les quatre formes de la coopération bilatérale en illustrant chacune par un exemple camerounais.
\item Pourquoi le Cameroun cherche-t-il à diversifier ses partenaires bilatéraux ? Rédige un paragraphe argumenté de dix lignes maximum.
\end{enumerate}
\end{exercice}

\begin{corrige}[Exercice : À toi de jouer]
\begin{enumerate}
\item La coopération bilatérale est l'ensemble des relations d'échanges et d'entraide entre deux États souverains dans des domaines variés. Instruments : les accords de coopération, les conventions, les traités, les ambassades et consulats (deux suffisent).
\item Économique et commerciale (exportation du cacao vers la France) ; financière (prêts chinois pour le port de Kribi) ; culturelle et scientifique (bourses d'études au Japon) ; militaire et technique (formation des forces de sécurité avec les USA).
\item Diversifier ses partenaires permet au Cameroun de ne pas dépendre d'un seul pays, de négocier de meilleures conditions, d'attirer davantage d'investissements et de technologies variées, et de renforcer son poids diplomatique sur la scène internationale.
\end{enumerate}
\end{corrige}

\coursec{Le Cameroun et la France}

\courssub{Transition : de la coopération générale au partenariat privilégié avec la France}

Dans la section précédente, nous avons défini la \cle{coopération bilatérale} comme une relation directe entre deux États souverains, fondée sur des accords juridiques et des intérêts partagés. Nous avons vu que le Cameroun, par sa position géostratégique en Afrique centrale et son histoire, entretient des relations bilatérales diversifiées. Parmi l'ensemble de ces partenaires, la \textbf{France} occupe une place singulière, à la fois historique, structurelle et contemporaine. Ce n'est pas seulement un partenaire parmi d'autres : c'est l'ancien pouvoir mandataire, le premier partenaire commercial historique, le principal pourvoyeur d'aide publique au développement et le vecteur principal de la \cle{francophonie} institutionnelle au Cameroun. Comprendre cette relation spécifique, c'est saisir une clé majeure de la politique étrangère, de l'économie et de l'identité culturelle du Cameroun contemporain.

\courssub{Fondements historiques : du mandat à l'indépendance}

\begin{definition}[Le mandat français sur le Cameroun]
Le \cle{mandat français} est l'administration confiée à la France par la \cle{Société des Nations (SDN)} en 1922 sur la partie orientale du Cameroun (environ 4/5 du territoire actuel), après la défaite de l'Allemagne lors de la Première Guerre mondiale (Traité de Versailles, 1919). Ce n'était pas une colonisation classique au sens juridique strict, mais une \emph{tutelle internationale} : la France avait l'obligation de préparer le territoire à l'autonomie, de respecter les droits des populations et de rendre des comptes à la SDN. En pratique, l'administration française a appliqué le droit colonial (code de l'indigénat, travail forcé, mise en valeur économique par les compagnies concessionnaires) tout en développant des infrastructures (chemin de fer Douala-Yaoundé, port de Douala, routes) et un système éducatif francophone.
\end{definition}

\begin{savaistu}[Héritage linguistique et institutionnel]
Le Cameroun est l'un des rares pays au monde à avoir \textbf{deux langues officielles} : le français et l'anglais. Cet héritage direct de la période des mandats (français à l'Est, britannique à l'Ouest) structure encore aujourd'hui l'administration, la justice (droit civil français à l'Est, common law à l'Ouest), l'éducation et la vie politique. L'adhésion à l'Organisation internationale de la Francophonie (OIF) en 1970 et au Commonwealth en 1995 consacre ce bilinguisme officiel unique en Afrique.
\end{savaistu}

L'accès à la \textbf{souveraineté internationale} intervient le \textbf{1er janvier 1960} pour le Cameroun oriental (République du Cameroun), sous la présidence d'\textbf{Ahmadou Ahidjo}. La réunification avec le Cameroun britannique (Southern Cameroons) n'interviendra qu'en octobre 1961 (plébiscite de février 1961), formant la République fédérale du Cameroun. Dès la veille de l'indépendance, le \textbf{13 novembre 1959}, puis le \textbf{27 février 1960}, Ahidjo signe avec le Général de Gaulle une série d'\textbf{accords de coopération} (culturels, techniques, militaires, économiques) qui instituent un cadre juridique privilégié, souvent qualifié de \emph{« coopération d'exception »}.

\begin{aretenir}[Les accords de 1959-1960 : le socle juridique]
Ces accords organisent :
\begin{itemize}
    \item La \textbf{coopération technique} : maintien de fonctionnaires français (coopérants) dans l'administration camerounaise (éducation, santé, finances, justice) pour assurer la continuité de l'État.
    \item La \textbf{coopération militaire} : accords de défense, formation des officiers, présence de forces françaises (éléments des Forces armées dans la zone de coopération - FAZC, base de Garoua).
    \item La \textbf{coopération monétaire} : maintien du \cle{Franc CFA} (Franc des Colonies Françaises d'Afrique, devenu Franc de la Coopération Financière en Afrique), arrimé au Franc français puis à l'Euro, avec garantie de convertibilité par le Trésor français.
    \item La \textbf{coopération culturelle} : développement de l'enseignement du français, création de lycées français, bourses d'études en France.
\end{itemize}
\end{aretenir}

\courssub{Coopération économique : un partenaire commercial et investisseur majeur}

La France reste, malgré la diversification des partenaires (Chine, USA, Inde), l'un des tout premiers partenaires commerciaux du Cameroun.

\begin{exemplebox}[Structure des échanges Cameroun-France (données simplifiées annuelles moyennes récentes)]
\begin{itemize}
    \item \textbf{Exportations camerounaises vers la France :} Produits pétroliers (brut), bois (grumes, sciages), cacao, café, coton, aluminium, bananes. La France est souvent le 1er ou 2e client pour le bois et le cacao.
    \item \textbf{Importations camerounaises depuis la France :} Matériel de transport (automobiles, pièces), machines et équipements industriels, produits pharmaceutiques, produits chimiques, céréales (blé), équipements électriques/électroniques.
    \item \textbf{Balance commerciale :} Structurellement excédentaire pour le Cameroun grâce au pétrole, mais volatile selon le cours du baril. Hors pétrole, le déficit est fréquent.
\end{itemize}
\end{exemplebox}

\begin{popfigure}
\begin{tikzpicture}
\begin{axis}[
    ybar,
    bar width=18pt,
    width=\linewidth,
    height=10cm,
    enlarge x limits=0.25,
    ylabel={Valeur (milliards FCFA)},
    symbolic x coords={2018,2019,2020,2021,2022,2023},
    xtick=data,
    nodes near coords,
    nodes near coords align={vertical},
    every node near coord/.append style={font=\tiny, rotate=90, anchor=west, yshift=2pt},
    legend style={at={(0.5,-0.15)},anchor=north,legend columns=-1},
    ymin=0,
    ymax=1800,
    tick label style={font=\small},
    label style={font=\small},
    grid=major,
    grid style={popline},
]
\addplot[popblue!80, fill=popblue!30] coordinates {(2018,950) (2019,1020) (2020,780) (2021,1100) (2022,1350) (2023,1250)};
\addplot[poporange!80, fill=poporange!30] coordinates {(2018,420) (2019,450) (2020,380) (2021,480) (2022,550) (2023,580)};
\legend{Exportations Cameroun $\rightarrow$ France, Importations Cameroun $\leftarrow$ France}
\end{axis}
\end{tikzpicture}
\legende{Figure 1 : Évolution simplifiée des échanges commerciaux Cameroun-France (source : Douanes Cameroun / DG Trésor France, ordres de grandeur arrondis). La chute de 2020 reflète la crise Covid-19 (pétrole). La reprise 2021-2022 suit la hausse des cours des matières premières.}
\end{popfigure}

\begin{exemplebox}[Les grands groupes français implantés au Cameroun]
La présence économique française ne se limite pas au commerce ; elle est structurelle via des filiales de grands groupes (souvent leaders sur leur segment) :
\begin{itemize}
    \item \textbf{Bolloré Africa Logistics :} Gestion du terminal à conteneurs du port de Douala (concession), logistique ferroviaire (Camrail), hydrocarbures.
    \item \textbf{TotalEnergies Marketing Cameroun :} Leader de la distribution de produits pétroliers (stations-service, aviation, lubrifiants), exploration/production (amont) via TotalEnergies EP Cameroun.
    \item \textbf{Orange Cameroun :} Opérateur historique de téléphonie mobile et internet (partenaire de l'État), services financiers (Orange Money).
    \item \textbf{Société Générale Cameroun :} Banque de référence, financement de l'économie, trésorerie de l'État.
    \item \textbf{Autres :} CFAO (automobile, équipements, consommateurs), Vinci (construction, concessions), Veolia / Suez (eau, déchets - historique), Sanofi (pharmacie), Michelin (pneumatiques, hévéa via Socapalm/HEVECAM), Castel (brasseries, boissons).
\end{itemize}
Ces entreprises sont de gros employeurs, contributeurs majeurs au budget de l'État (impôts, douanes) et vecteurs de transfert de technologie et de formation.
\end{exemplebox}

\courssub{Coopération financière : l'AFD, bras armé du développement}

L'aide publique au développement (APD) française est canalisée principalement par l'\textbf{Agence Française de Développement (AFD)}, établissement public financier. Elle se décline en \textbf{dons} (subventions, expertises) et \textbf{prêts} (à taux concessionnels ou de marché), souvent délégués à l'État camerounais ou à des structures publiques.

\begin{propriete}[Les secteurs d'intervention prioritaires de l'AFD au Cameroun (2015-2025)]
L'alignement se fait sur la \cle{Stratégie Nationale de Développement 2020-2030 (SND30)} :
\begin{enumerate}
    \item \textbf{Eau et Assainissement :} Projet d'alimentation en eau de Yaoundé (AEP Yaoundé - rivière Sanaga), projet d'assainissement de Douala (PADE), projets ruraux (PNDP).
    \item \textbf{Energie :} Barrage de Lom Pangar (financement majeur), transport d'énergie (lignes haute tension), accès à l'électricité rurale (PERG), appui à la transition énergétique (solaire).
    \item \textbf{Transport / Infrastructures :} Réhabilitation route Nationale 1 (Douala-Yaoundé), route Douala-Bafoussam, appui au port de Kribi (accès), transport urbain (bus Yaoundé/Douala).
    \item \textbf{Agriculture / Développement rural :} Programme d'appui aux chaînes de valeur (cacao, coton, maïs), résilience climatique (Sahel).
    \item \textbf{Gouvernance / Secteur financier :} Appui budgétaire (réformes finances publiques), inclusion financière, appui à la BEAC/COBAC.
    \item \textbf{Education / Formation professionnelle :} Construction/réhabilitation d'écoles techniques, appui à l'insertion professionnelle (PAJER, PROJET).
\end{enumerate}
\end{propriete}

\begin{popfigure}
\begin{tikzpicture}[scale=0.95, every node/.style={font=\small}]
% Contour très simplifié du Cameroun (polygone approximatif)
\draw[popdark, thick] (0,0) -- (5,0.5) -- (6,3) -- (5.5,6) -- (3,7) -- (0.5,6.5) -- (-0.5,4) -- (-0.2,1.5) -- cycle;
% Villes repères
\node[circle, fill=popblue, inner sep=1.5pt, label=above right:\tiny Yaoundé] (Y) at (2.5, 3.5) {};
\node[circle, fill=popblue, inner sep=1.5pt, label=below left:\tiny Douala] (D) at (1.5, 1.8) {};
\node[circle, fill=popblue, inner sep=1.5pt, label=above:\tiny Garoua] (G) at (3.5, 5.5) {};
\node[circle, fill=popblue, inner sep=1.5pt, label=above right:\tiny Bertoua] (B) at (4.5, 4) {};
\node[circle, fill=popblue, inner sep=1.5pt, label=below right:\tiny Kribi] (K) at (0.5, 1) {};
\node[circle, fill=popblue, inner sep=1.5pt, label=left:\tiny Bafoussam] (Bf) at (1.2, 3.2) {};
\node[circle, fill=popblue, inner sep=1.5pt, label=right:\tiny Maroua] (M) at (4.2, 5.8) {};

% Projets AFD (carrés orange)
\node[rectangle, fill=poporange, inner sep=2pt, label=above:\tiny Lom Pangar (Énergie)] at (3.8, 4.5) {};
\node[rectangle, fill=poporange, inner sep=2pt, label=right:\tiny AEP Yaoundé (Eau)] at (2.8, 3.8) {};
\node[rectangle, fill=poporange, inner sep=2pt, label=below:\tiny Assainissement Douala] at (1.5, 1.3) {};
\node[rectangle, fill=poporange, inner sep=2pt, label=left:\tiny Route N1 / N6] at (1.8, 2.8) {};
\node[rectangle, fill=poporange, inner sep=2pt, label=above:\tiny Énergie Rural (Nord)] at (3.5, 5.8) {};
\node[rectangle, fill=poporange, inner sep=2pt, label=below:\tiny Port Kribi (Accès)] at (0.2, 0.8) {};

% Entreprises françaises (losanges verts)
\node[diamond, fill=popgreen, inner sep=2pt, label=right:\tiny Bolloré (Port Douala)] at (1.3, 1.6) {};
\node[diamond, fill=popgreen, inner sep=2pt, label=above:\tiny TotalEnergies (Dépôts)] at (2.0, 2.0) {};
\node[diamond, fill=popgreen, inner sep=2pt, label=above:\tiny Orange (Siège Yaoundé)] at (2.7, 3.3) {};
\node[diamond, fill=popgreen, inner sep=2pt, label=left:\tiny SOCAPALM/HEVECAM (Michelin)] at (0.8, 2.5) {};
\node[diamond, fill=popgreen, inner sep=2pt, label=right:\tiny CFAO/Vinci (Chantiers)] at (2.2, 3.0) {};

% Légende manuelle
\begin{scope}[shift={(6.5, 5.5)}]
\node[rectangle, fill=poporange, inner sep=2pt] at (0,0.3) {}; \node[anchor=west] at (0.3,0.3) {\tiny Projets AFD (Eau, Énergie, Routes)};
\node[diamond, fill=popgreen, inner sep=2pt] at (0,0) {}; \node[anchor=west] at (0.3,0) {\tiny Implantations Entreprises FR};
\node[circle, fill=popblue, inner sep=1.5pt] at (0,-0.3) {}; \node[anchor=west] at (0.3,-0.3) {\tiny Villes principales};
\end{scope}

\node[anchor=north west, font=\small\bfseries, text=popdark] at (-0.5, 7.2) {Carte schématique : Projets AFD \& Entreprises françaises au Cameroun};
\draw[->, thick] (6.5, 7) -- (6.5, 6.5) node[midway, right] {Nord};
\end{tikzpicture}
\legende{Figure 2 : Localisation schématique des principaux projets financés par l'AFD (carrés orange) et des implantations majeures d'entreprises françaises (losanges verts). Le Nord (Garoua, Maroua) bénéficie d'appuis énergie/eau/rural ; le triangle Douala-Yaoundé-Bafoussam concentre les infrastructures de transport, l'industrie et les sièges sociaux.}
\end{popfigure}

\courssub{Coopération culturelle, éducative et linguistique : le ciment de la Francophonie}

C'est le volet le plus visible par les populations et le plus structurant pour l'élite camerounaise.

\begin{itemize}
    \item \textbf{La langue française :} Langue de travail de l'administration, de la justice (zone francophone), de l'enseignement secondaire et supérieur (majoritaire), des médias d'État. Outil d'intégration nationale face à la diversité linguistique (250+ langues locales).
    \item \textbf{Réseau scolaire français :} \emph{Lycée Leclerc} (Yaoundé), \emph{Lycée Français de Douala}, \emph{Lycée Français de Garoua}, \emph{École Française Internationale} (Bertoua, etc.). Homologués par l'Éducation nationale française, ils préparent au Baccalauréat français et accueillent élèves camerounais et expatriés. Vecteurs d'excellence et de mobilité.
    \item \textbf{Instituts Français (IF) :} Yaoundé, Douala, Garoua, Bafoussam, Ngaoundéré. Centres culturels, médiathèques, cours de langue (FLE), programmation artistique, appui aux industries culturelles locales (musique, cinéma, littérature).
    \item \textbf{Bourses et mobilité étudiante :} Programme \emph{Eiffel}, bourses du gouvernement français, cotutelles de thèse, partenariats universitaires (Université de Yaoundé I/II, Ngaoundéré, Dschang, Douala avec Montpellier, Bordeaux, Paris-Sorbonne, etc.). La France reste la 1ère destination des étudiants camerounais à l'étranger (hors Afrique).
\end{itemize}

\begin{attention}[Piège à éviter : confusion « Aide » vs « Coopération culturelle »]
Ne pas réduire la coopération culturelle à de l'\emph{aide humanitaire}. C'est un \textbf{investissement stratégique} (soft power) : former les futures élites, diffuser la langue et les normes juridiques/techniques françaises, créer des réseaux d'influence durables (alumni). Le Cameroun y trouve son compte : formation de haut niveau, diplômes reconnus internationalement, ouverture sur la recherche mondiale. C'est un \emph{échange de valeurs}, pas un don unilatéral.
\end{attention}

\courssub{Coopération militaire et technique : défense et formation}

\begin{definition}[Accords de défense (1960, rénovés 1974, 2019)]
Ces accords juridiques lient les deux États en matière de sécurité. Ils prévoient :
\begin{itemize}
    \item L'\textbf{assistance militaire} : formation des officiers (École de guerre, Saint-Cyr, écoles d'application), sous-officiers et soldats (École militaire interarmes - EMIA Yaoundé, centres d'instruction).
    \item La \textbf{mise à disposition de coopérants militaires} : officiers français en poste de conseil ou d'instruction dans les écoles et états-majors camerounais.
    \item La \textbf{présence prépositionnée} : détachements (ex: base aérienne de Garoua, éléments à Douala/Yaoundé) pour appui logistique, renseignement, évacuation de ressortissants, formation au profit des forces de la sous-région (force multinationale mixte - FMM contre Boko Haram).
    \item L'\textbf{intervention en cas d'agression extérieure} (clause de défense), bien que l'application politique reste souveraine.
\end{itemize}
\end{definition}

\begin{exemplebox}[Opérationnalité : Lutte contre Boko Haram et sécurité maritime]
La coopération militaire s'est intensifiée depuis 2014 face à l'insurrection de Boko Haram dans l'Extrême-Nord :
\begin{itemize}
    \item Appui renseignement (drones, imagerie satellite partagée).
    \item Formation des unités spéciales (BIR - Bataillon d'Intervention Rapide) et appui logistique (véhicules, transmissions, médical).
    \item Sécurité maritime : Golfe de Guinée (piraterie), appui à la marine camerounaise (patrouilleurs, formation, exercices conjoints type \emph{NEMO} / \emph{Grand African NEMO}).
\end{itemize}
\end{exemplebox}

\courssub{Les limites, débats et enjeux contemporains : au-delà du « couple privilégié »}

La relation n'est pas exempte de tensions structurelles et de critiques vives, tant au Cameroun qu'en France.

\begin{aretenir}[Trois débats majeurs structurants]
\begin{enumerate}
    \item \textbf{La dépendance économique et le « Franc CFA » :}
    \begin{itemize}
        \item Le Franc CFA (XAF) garantit la \emph{stabilité monétaire} (inflation maîtrisée, change fixe vs Euro), atout pour les investisseurs et l'importation.
        \item \textbf{Critique :} Perte de souveraineté monétaire (taux de change non ajustable aux chocs asymétriques, réserves déposées au Trésor français à 50\% historiquement, réforme 2019 : fin de la centralisation des réserves à 50\% et retrait des administrateurs français au sein de la BEAC, mais parité fixe maintenue). Débat sur la compétitivité des exportations (taux de change surévalué ?) et la capacité de financement endogène.
    \end{itemize}
    \item \textbf{La « Françafrique » : réseaux d'influence et opacité :}
    \begin{itemize}
        \item Terme désignant les \emph{réseaux informels} (politiques, affairistes, militaires, services secrets) qui auraient perpétué une domination néocoloniale au-delà des accords officiels : soutien à des régimes autoritaires, contrats publics opaques (marchés de gré à gré), évasion fiscale, ingérence électorale.
        \item \textbf{Réponse officielle :} Discours de rupture (Sarkozy 2007, Hollande 2012, Macron 2017 : « La Françafrique est terminée »), transparence accrue (loi Sapin II, registres des bénéficiaires effectifs), mais persistance de soupçons sur certains dossiers (port de Kribi, barrages, marchés de défense).
    \end{itemize}
    \item \textbf{Déséquilibre commercial et transfert de technologie :}
    \begin{itemize}
        \item Le Cameroun exporte mostly des matières premières brutes (pétrole, bois, cacao) et importe des biens manufacturés à forte valeur ajoutée.
        \item Faible part de la transformation locale dans les filières tenues par des groupes français (ex: grumes vs meubles, cacao brut vs chocolat, pétrole brut vs raffiné). Enjeu : clauses de transformation locale dans les conventions d'établissement, formation d'ingénieurs camerounais aux postes de direction.
    \end{itemize}
\end{enumerate}
\end{aretenir}

\begin{exoresolu}[Analyse critique d'un article de presse : « Le Cameroun diversifie ses partenariats »]
\textbf{Énoncé :} Un éditorial affirme : « Le Cameroun tourne le dos à la France pour se tourner vers la Chine et la Russie. La coopération française est morte. » En vous appuyant sur les données du cours, nuancez cette affirmation.

\textbf{Solution détaillée :}
\begin{enumerate}
    \item \textbf{Faux : « Tourner le dos ».} Les chiffres du commerce (Fig. 1) montrent une interdépendance massive. La France reste un client majeur pour le pétrole/bois/cacao et un fournisseur incontournable de machines/médicaments. Les stocks d'IDE (Investissements Directs Étrangers) français sont colossaux (Bolloré, Total, Orange, SG). On ne « tourne le dos » à son premier créancier bilatéral (AFD) et à son premier partenaire éducatif du jour au lendemain.
    \item \textbf{Vrai : « Diversification ».} La part de la France dans le commerce total baisse (autour de 10-12\% aujourd'hui vs 40\% dans les années 80). La Chine est devenue le 1er fournisseur (machines, textiles, BTP), la Russie un partenaire militaire/énergétique (Rosatom, instructeurs). C'est une \textbf{stratégie rationnelle de diversification} (ne pas mettre tous les œufs dans le même panier), pas une rupture.
    \item \textbf{Nuance : « Coopération morte ».} Les signatures récentes (Cadre de partenariat 2023-2027 AFD, accords défense 2019, sommet Afrique-France 2021 Montpellier, visite Macron 2023) prouvent une \textbf{reconfiguration}, pas une fin. Le passage est d'une \emph{coopération d'exception} (quasi exclusive, paternaliste) à un \emph{partenariat d'égaux} (contractuel, concurrentiel, aligné sur SND30).
    \item \textbf{Conclusion :} Le Cameroun exerce sa \cle{souveraineté} en élargissant son portefeuille de partenaires. La France doit désormais \emph{concourir} (qualité/prix/délais) là où elle était \emph{privilégiée}. C'est une maturation de la relation, pas sa mort.
\end{enumerate}
\end{exoresolu}

\begin{methode}[Comment analyser un accord de coopération bilatérale (grille de lecture pour le Bac)]
Pour tout texte d'accord ou article de presse sur la coopération Cameroun-France :
\begin{enumerate}
    \item \textbf{Identifier le cadre juridique :} Traité ? Accord-cadre ? Protocole d'accord ? Échange de lettres ? (Force juridique différente).
    \item \textbf{Repérer les secteurs couverts :} Économie/Finance ? Défense/Sécurité ? Culture/Éducation ? Recherche/Innovation ? Gouvernance ?
    \item \textbf{Analyser la réciprocité :} Qu'est-ce que le Cameroun donne ? (Marché, ressources, base militaire, voix à l'ONU, alignement diplomatique). Qu'est-ce que la France donne ? (Financement, technologie, formation, garantie monétaire, accès au marché UE, appui diplomatique).
    \item \textbf{Évaluer la souveraineté :} Clauses de résiliation ? Juridiction de règlement des différends ? Conditionnalités politiques (droits de l'homme, gouvernance) ? Alignement sur SND30 ?
    \item \textbf{Situer dans le temps :} Rupture ou continuité par rapport aux accords 1960 ? Adaptation aux nouveaux enjeux (climat, numérique, terrorisme, jeunesse) ?
\end{enumerate}
\end{methode}

\begin{exercice}[Entraînement : La réforme du Franc CFA (Eco)]
\textbf{Consigne :} En 2019, la France et l'UEMOA (zone CFA Ouest) annoncent la réforme du Franc CFA vers l'\textbf{Eco} (changement de nom, fin de la centralisation des réserves à 50\% au Trésor français, retrait des administrateurs français). La zone CEMAC (Cameroun inclus) n'a pas suivi immédiatement.
\begin{enumerate}
    \item Rappeler les 3 fonctions de la monnaie et dire lesquelles le Franc CFA remplit bien / mal pour le Cameroun.
    \item Expliquer pourquoi le Cameroun (zone CEMAC) a été plus prudent que la zone UEMOA sur cette réforme.
    \item Donner deux arguments \textbf{pour} et deux arguments \textbf{contre} le maintien de la parité fixe Euro/FCFA pour l'économie camerounaise.
\end{enumerate}
\end{exercice}

\begin{corrige}[Exercice n° : La réforme du Franc CFA (Eco)]
\begin{enumerate}
    \item \textbf{Fonctions :} Unité de compte (oui, stable), Moyen d'échange (oui, convertible garanti), Réserve de valeur (oui, inflation faible). \textbf{Mal remplie :} Instrument d'ajustement externe (taux de change fixe empêche la dévaluation pour regagner compétitivité après choc pétrolier/cacao) ; Prêteur en dernier ressort national (pas de banque centrale nationale, BEAC régionale).
    \item \textbf{Prudence CEMAC :} Hétérogénéité économique plus forte (pétroliers vs non-pétroliers) ; Gouvernance de la BEAC perfectible (scandales passés) ; Risque de fragmentation monétaire en Afrique centrale ; Besoin de convergence économique préalable (critères de convergence non tous respectés).
    \item \textbf{Pour :} Stabilité des prix (ancrage anti-inflation), Confiance investisseurs (pas de risque de change), Facilite importations/énergie/médicaments, Intégration régionale (monnaie unique CEMAC).
    \item \textbf{Contre :} Perte compétitivité exportations (taux surévalué si inflation partenaires > inflation locale), Politique monétaire inadaptée (taux directeurs calqués sur BCE, pas sur cycle camerounais), Contrainte budgétaire (rigueur pour défendre parité), Frein à l'industrialisation par substitution aux importations.
\end{enumerate}
\end{corrige}

\coursec{Le Cameroun et la Grande-Bretagne}

Après avoir étudié la relation privilégiée entre le Cameroun et la France, héritière du mandat français sur la plus grande partie du territoire, nous nous tournons maintenant vers l'autre grande puissance mandataire : la \cle{Grande-Bretagne}. Si la présence britannique au Cameroun fut géographiquement plus limitée que la présence française, elle n'en a pas moins laissé des traces profondes et durables : la langue anglaise, le système éducatif anglo-saxon, la Common Law et l'adhésion du Cameroun au Commonwealth. Comprendre la coopération camerouno-britannique, c'est d'abord comprendre pourquoi le Cameroun est aujourd'hui un pays officiellement bilingue et biculturel.

\courssub{Les fondements historiques : le mandat britannique (1922-1961)}

\textbf{De la colonie allemande au partage franco-britannique.}

De 1884 à 1916, le Cameroun est une colonie allemande : le \emph{Kamerun}. Pendant la Première Guerre mondiale, les troupes françaises et britanniques chassent les Allemands du territoire (campagne du Cameroun, 1914-1916). À la fin de la guerre, la Société des Nations (SDN) confie l'administration du territoire aux deux vainqueurs sous forme de \cle{mandats} (1922), transformés en \cle{tutelles} de l'ONU après 1945.

\begin{definition}[Mandat et tutelle]
Un \cle{mandat} est une mission d'administration d'un territoire confiée par la Société des Nations à une puissance, qui doit préparer ce territoire à l'autonomie. Après 1945, l'ONU reprend le même principe sous le nom de \cle{tutelle}. Le Cameroun fut ainsi partagé : environ les quatre cinquièmes du territoire revinrent à la France, et deux bandes occidentales, le long de la frontière nigériane, à la Grande-Bretagne.
\end{definition}

La Grande-Bretagne administre donc deux zones distinctes :
\begin{itemize}
\item le \cle{Southern Cameroons} (Cameroun méridional britannique), au sud, avec pour chef-lieu Buéa ;
\item le \cle{Northern Cameroons} (Cameroun septentrional britannique), au nord.
\end{itemize}
Ces deux zones ne sont pas administrées directement comme des colonies : elles sont rattachées administrativement au \cle{Nigeria} britannique voisin, ce qui aura des conséquences décisives en 1961.

\begin{attention}[Piège fréquent]
Le Cameroun n'a \textbf{jamais} été entièrement colonisé par la Grande-Bretagne. Seules deux bandes occidentales (Southern et Northern Cameroons) étaient sous administration britannique, entre 1922 et 1961. Dire que << le Cameroun était une colonie anglaise >> est une erreur grave à l'examen : la majeure partie du territoire était sous mandat français.
\end{attention}

\begin{popfigure}
\begin{center}
\begin{tikzpicture}[scale=0.9]
% Contour simplifié du Cameroun
\draw[thick,popdark,fill=popcream] (2.2,6.2) -- (3.4,6.6) -- (4.6,6.1) -- (5.2,5.2) -- (6.4,4.6) -- (7.2,3.6) -- (6.6,2.4) -- (5.6,1.4) -- (4.4,0.6) -- (3.0,0.2) -- (1.8,0.8) -- (1.2,2.0) -- (0.8,3.2) -- (1.0,4.4) -- (1.6,5.4) -- cycle;
% Zone britannique (bande occidentale)
\draw[thick,popblue,fill=popblue!35] (1.2,2.0) -- (0.8,3.2) -- (1.0,4.4) -- (1.6,5.4) -- (2.2,6.2) -- (2.6,5.2) -- (2.4,4.0) -- (2.2,2.8) -- (1.8,1.6) -- cycle;
% Frontière Nigeria
\draw[dashed,popmuted] (0.4,6.4) -- (0.2,0.4);
\node[popmuted,rotate=85,font=\small] at (0.05,3.4) {Nigeria (britannique)};
% Labels
\node[popblue,font=\small\bfseries,align=center] at (1.55,4.6) {Northern\\Cameroons};
\node[popblue,font=\small\bfseries,align=center] at (1.5,2.4) {Southern\\Cameroons};
\node[popdark,font=\small\bfseries,align=center] at (4.6,3.4) {Mandat\\français\\(environ 4/5\\du territoire)};
\node[popblue,font=\footnotesize] at (3.1,5.6) {Mandat britannique};
\draw[->,popblue] (2.7,5.5) -- (2.0,5.3);
% Nord
\draw[->,thick,popdark] (7.6,5.6) -- (7.6,6.4);
\node[popdark,font=\small] at (7.6,6.7) {N};
% Villes
\fill[popink] (1.35,1.5) circle (1.5pt); \node[right,font=\footnotesize] at (1.45,1.5) {Buéa};
\fill[popink] (4.0,2.2) circle (1.5pt); \node[right,font=\footnotesize] at (4.1,2.2) {Yaoundé};
\fill[popink] (2.6,1.1) circle (1.5pt); \node[right,font=\footnotesize] at (2.7,1.1) {Douala};
\end{tikzpicture}
\end{center}
\legende{Carte schématique du partage du Cameroun allemand (1916-1922) : le mandat français (en clair) couvre l'essentiel du territoire ; le mandat britannique (en bleu) se limite à deux bandes occidentales le long du Nigeria. Croquis simplifié, sans précision cartographique.}
\end{popfigure}

\courssub{Le plébiscite du 11 février 1961 : deux destinées différentes}

\textbf{Le problème posé.} Le 1\textsuperscript{er} janvier 1960, le Cameroun français accède à l'indépendance et devient la République du Cameroun. Le Nigeria devient indépendant le 1\textsuperscript{er} octobre 1960. Que faire alors des deux zones britanniques, administrées depuis le Nigeria ? L'ONU décide de consulter directement les populations par un \cle{plébiscite}.

\begin{definition}[Plébiscite]
Un \cle{plébiscite} est une consultation populaire par laquelle les habitants d'un territoire décident eux-mêmes, par vote, de leur destinée politique. Ici, la question posée était : souhaitez-vous rejoindre la République du Cameroun ou la République fédérale du Nigeria ?
\end{definition}

\textbf{Les résultats du 11 février 1961.} Les deux zones font des choix opposés :
\begin{itemize}
\item le \cle{Southern Cameroons} vote majoritairement pour la réunification avec la République du Cameroun (environ 233 571 voix contre 97 741 pour le Nigeria) ; la réunification est effective le 1\textsuperscript{er} octobre 1961, date de naissance de la République fédérale du Cameroun ;
\item le \cle{Northern Cameroons} vote majoritairement pour le rattachement au Nigeria (environ 146 296 voix contre 97 659), rattachement effectif le 1\textsuperscript{er} juin 1961.
\end{itemize}

\begin{methode}[Expliquer le bilinguisme camerounais]
Pour expliquer correctement le bilinguisme officiel du Cameroun, procède toujours en trois étapes :
\begin{enumerate}
\item remonter au \textbf{double héritage colonial} : mandat français sur l'essentiel du territoire, mandat britannique sur deux bandes occidentales ;
\item citer le \textbf{plébiscite du 11 février 1961} : le Southern Cameroons choisit de rejoindre la République du Cameroun, apportant avec lui sa langue et ses institutions anglaises ;
\item conclure par la \textbf{Constitution} : le Cameroun est officiellement bilingue (français et anglais, les deux langues ayant le même statut), héritage direct de cette histoire.
\end{enumerate}
Cette démarche chronologique et causale est celle attendue au Probatoire et au Baccalauréat technique.
\end{methode}

\begin{popfigure}
\begin{center}
\begin{tikzpicture}[scale=0.85]
% Cameroun simplifié
\draw[thick,popdark,fill=popcream] (2.6,5.6) -- (4.2,6.0) -- (5.4,5.0) -- (6.6,4.2) -- (7.0,3.2) -- (6.2,2.0) -- (5.0,1.0) -- (3.6,0.4) -- (2.4,0.8) -- (1.8,2.0) -- (1.6,3.4) -- (2.0,4.6) -- cycle;
% Southern Cameroons
\draw[thick,popgreen,fill=popgreen!35] (1.8,2.0) -- (1.6,3.4) -- (2.0,4.6) -- (2.6,5.6) -- (2.9,4.4) -- (2.7,3.0) -- (2.3,1.8) -- cycle;
\node[popgreen!60!black,font=\small\bfseries,align=center] at (2.15,3.6) {Southern\\Cameroons};
% Flèche vers République du Cameroun
\draw[->,very thick,popgreen!60!black] (2.6,2.6) -- (4.2,2.6);
\node[popgreen!60!black,font=\footnotesize,align=center] at (3.4,2.15) {1\textsuperscript{er} oct. 1961 :\\République du Cameroun};
% Northern Cameroons (au nord, séparé)
\draw[thick,poporange,fill=poporange!35] (2.6,6.6) -- (4.0,7.0) -- (5.0,6.6) -- (4.6,6.1) -- (3.2,6.15) -- cycle;
\node[poporange!70!black,font=\small\bfseries] at (3.8,6.55) {Northern Cameroons};
% Nigeria
\draw[thick,popmuted,fill=popmuted!20] (0.2,5.8) -- (1.4,6.2) -- (2.4,6.4) -- (2.6,7.2) -- (0.2,7.4) -- cycle;
\node[popmuted,font=\small\bfseries] at (1.2,6.8) {Nigeria};
% Flèche Northern -> Nigeria
\draw[->,very thick,poporange!70!black] (2.7,6.75) -- (1.7,6.9);
\node[poporange!70!black,font=\footnotesize,align=center] at (2.0,7.5) {1\textsuperscript{er} juin 1961 :\\rattachement au Nigeria};
% Nord
\draw[->,thick,popdark] (7.4,6.6) -- (7.4,7.4);
\node[popdark,font=\small] at (7.4,7.7) {N};
% Titre du vote
\node[popdark,font=\small\bfseries,align=center] at (4.6,0.0) {Plébiscite du 11 février 1961 (ONU)};
\end{tikzpicture}
\end{center}
\legende{Croquis du plébiscite de 1961 : le Southern Cameroons (vert) choisit la République du Cameroun ; le Northern Cameroons (orange) choisit le Nigeria. Croquis simplifié.}
\end{popfigure}

\courssub{L'héritage anglophone : langue, éducation et droit}

Le retour du Southern Cameroons dans la nation camerounaise a apporté un héritage institutionnel spécifique, encore visible aujourd'hui dans les régions du \cle{Nord-Ouest} et du \cle{Sud-Ouest}.

\textbf{1. La langue anglaise.} L'anglais est, avec le français, langue officielle du Cameroun. Il est langue d'enseignement, d'administration et de justice dans les régions anglophones, et langue de communication dans les affaires internationales.

\textbf{2. Le système éducatif anglo-saxon.} Les régions anglophones conservent un sous-système éducatif distinct : le \cle{GCE} (\emph{General Certificate of Education}), avec l'\emph{Ordinary Level} et l'\emph{Advanced Level}, sanctionne les études secondaires, sous le contrôle du \emph{GCE Board} créé en 1993. Ce système coexiste avec le sous-système francophone (BEPC, Probatoire, Baccalauréat), dont relèvent les élèves de l'enseignement technique francophone.

\textbf{3. Le droit : Common Law et droit coutumier.} Dans les régions du Nord-Ouest et du Sud-Ouest s'applique la \cle{Common Law}, droit d'origine britannique fondé sur la jurisprudence (les décisions des juges font règle), par opposition au droit civil d'inspiration française appliqué dans les régions francophones. S'y ajoutent les tribunaux coutumiers qui jugent selon les coutumes locales (mariages, successions, terres). Le Cameroun est donc un État \cle{bijuridique} : deux traditions juridiques coexistent sur son territoire.

\begin{aretenir}[L'essentiel sur l'héritage anglophone]
Le Cameroun tire de son histoire trois spécificités rares en Afrique : le \textbf{bilinguisme} officiel (français-anglais), le \textbf{biculturalisme} éducatif (deux sous-systèmes scolaires) et le \textbf{bijuridisme} (droit civil et Common Law). Toutes trois découlent du mandat britannique et du plébiscite de 1961.
\end{aretenir}

\courssub{La coopération économique : échanges, entreprises et Commonwealth}

\textbf{Les échanges commerciaux.} La Grande-Bretagne est un partenaire commercial historique du Cameroun : elle importe des produits camerounais (pétrole brut, cacao, café, bananes, bois) et exporte vers le Cameroun des machines, des véhicules, des produits pharmaceutiques et manufacturés. Des entreprises britanniques ou à capitaux britanniques ont longtemps été présentes dans l'agro-industrie (plantations du Sud-Ouest), la banque et le commerce.

\textbf{L'adhésion au Commonwealth (1995).} En 1995, le Cameroun adhère au \cle{Commonwealth}, devenant l'un des rares membres dont la majeure partie du territoire n'a jamais été colonie britannique.

\begin{definition}[Le Commonwealth]
Le \cle{Commonwealth} est une association volontaire de 56 pays, pour la plupart d'anciennes colonies britanniques, qui coopèrent dans les domaines de la démocratie, du développement, de l'éducation et de la culture. Le Cameroun en est membre depuis \textbf{1995}. L'adhésion s'explique par l'héritage du Southern Cameroons et par la volonté du Cameroun de valoriser sa composante anglophone sur la scène internationale.
\end{definition}

L'adhésion au Commonwealth offre au Cameroun des avantages concrets : accès à des programmes de formation et de bourses, participation aux Jeux du Commonwealth, coopération technique (renforcement des institutions, observation électorale), et ouverture sur les marchés anglophones.

\begin{savaistu}[Un pays << charnière >>]
Le Cameroun est l'un des très rares pays au monde membre à la fois de la \textbf{Francophonie} et du \textbf{Commonwealth} (avec le Rwanda, le Mozambique et le Gabon notamment). Ce double ancrage fait du Cameroun une véritable << charnière >> entre l'Afrique francophone et l'Afrique anglophone : un atout diplomatique et économique unique, souvent résumé par la formule << l'Afrique en miniature >>.
\end{savaistu}

\courssub{La coopération éducative et culturelle}

\textbf{Le British Council.} Présent à Yaoundé, le \cle{British Council} est l'agence britannique de coopération culturelle et éducative. Il organise les examens de certification en anglais (comme l'\emph{IELTS}, souvent exigé pour étudier ou travailler dans les pays anglophones), forme les enseignants d'anglais, met à disposition des ressources pédagogiques et soutient les échanges culturels (arts, débats, clubs de lecture).

\textbf{Les bourses Chevening.} Le programme \cle{Chevening} offre chaque année des bourses à de jeunes Camerounais méritants pour poursuivre des études de master dans les universités britanniques. Les boursiers, appelés << Chevening Scholars >>, reviennent ensuite occuper des postes de responsabilité dans l'administration, les entreprises ou la société civile camerounaise.

\textbf{L'enseignement de l'anglais.} La Grande-Bretagne appuie la promotion de l'anglais au Cameroun : formation des enseignants, fourniture de manuels, programmes d'échanges. Pour un élève de Terminale technique, la maîtrise de l'anglais est un atout professionnel majeur : nombre d'entreprises et de formations techniques supérieures exigent la pratique des deux langues officielles.

\courssub{La coopération dans d'autres domaines}

La coopération camerouno-britannique dépasse l'économie et l'éducation :
\begin{itemize}
\item \textbf{Gouvernance et appui institutionnel} : le Royaume-Uni soutient la décentralisation, la modernisation de l'administration, la lutte contre la corruption et le renforcement des processus électoraux, notamment à travers les programmes du Commonwealth et de la coopération britannique au développement ;
\item \textbf{Coopération militaire ponctuelle} : des formations et des exercices conjoints ont lieu périodiquement, notamment dans le cadre de la lutte contre le terrorisme dans la région du Lac Tchad et de la sécurisation du golfe de Guinée (piraterie maritime) ;
\item \textbf{Santé et action humanitaire} : appuis ponctuels lors de crises sanitaires ou humanitaires, en lien avec les organisations internationales.
\end{itemize}

\begin{exoresolu}[Expliquer l'adhésion du Cameroun au Commonwealth]
\textbf{Énoncé.} Un élève affirme : << Le Cameroun a adhéré au Commonwealth en 1995 parce qu'il était une colonie anglaise. >> Corrige cette affirmation et explique les vraies raisons de l'adhésion du Cameroun au Commonwealth.
\tcblower
\textbf{Solution détaillée.} L'affirmation contient une erreur historique : le Cameroun n'a jamais été entièrement une colonie anglaise. Après la Première Guerre mondiale, seules deux bandes occidentales, le Southern Cameroons et le Northern Cameroons, furent placées sous mandat puis tutelle britannique (1922-1961), l'essentiel du territoire étant sous administration française. Lors du plébiscite du 11 février 1961, le Southern Cameroons a choisi de rejoindre la République du Cameroun, apportant à la nation un héritage anglophone : langue anglaise, système éducatif anglo-saxon (GCE) et Common Law dans les régions du Nord-Ouest et du Sud-Ouest. C'est cet héritage qui fonde la légitimité de l'adhésion au Commonwealth en 1995. À cela s'ajoutent des raisons stratégiques : valoriser la composante anglophone du pays, accéder aux programmes de formation et de bourses du Commonwealth, diversifier les partenaires internationaux du Cameroun et renforcer son statut de charnière entre Afrique francophone et anglophone. \textbf{Conclusion} : l'adhésion de 1995 s'explique par l'héritage britannique d'une partie du territoire et par un choix diplomatique d'ouverture, non par une colonisation anglaise de l'ensemble du pays.
\end{exoresolu}

\begin{exercice}[À faire]
\begin{enumerate}
\item Définis : mandat, plébiscite, Commonwealth, Common Law.
\item Reconstitue la chronologie : 1916, 1922, 1960, 11 février 1961, 1\textsuperscript{er} octobre 1961, 1995.
\item Rédige un paragraphe de dix lignes expliquant pourquoi le Cameroun est un pays bilingue et bijuridique.
\end{enumerate}
\end{exercice}

\begin{corrige}[Exercice]
\textbf{1.} Mandat : mission d'administration d'un territoire confiée par la SDN à une puissance pour préparer ce territoire à l'autonomie. Plébiscite : consultation populaire par laquelle les habitants d'un territoire décident par vote de leur destinée politique. Commonwealth : association volontaire de 56 pays, pour la plupart anciennes colonies britanniques, dont le Cameroun est membre depuis 1995. Common Law : droit d'origine britannique fondé sur la jurisprudence, appliqué dans les régions du Nord-Ouest et du Sud-Ouest.

\textbf{2.} 1916 : fin de la colonisation allemande, occupation franco-britannique. 1922 : mandats de la SDN (France et Grande-Bretagne). 1960 : indépendance du Cameroun français (1\textsuperscript{er} janvier) et du Nigeria (1\textsuperscript{er} octobre). 11 février 1961 : plébiscite dans les Cameroons britanniques. 1\textsuperscript{er} octobre 1961 : réunification du Southern Cameroons avec la République du Cameroun (République fédérale). 1995 : adhésion du Cameroun au Commonwealth.

\textbf{3.} Éléments attendus : double héritage colonial franco-britannique ; mandat britannique sur le Southern Cameroons ; plébiscite du 11 février 1961 et réunification du 1\textsuperscript{er} octobre 1961 ; conséquences : bilinguisme officiel (français et anglais), deux sous-systèmes éducatifs (GCE et Baccalauréat), coexistence du droit civil français et de la Common Law britannique. Le Cameroun est donc à la fois bilingue, biculturel et bijuridique.
\end{corrige}

\coursec{Le Cameroun et les pays asiatiques : la Chine et le Japon}

Après avoir étudié les relations du Cameroun avec la Grande-Bretagne, héritées d'une histoire coloniale partagée, nous changeons de continent pour nous tourner vers l'Asie. Depuis quelques décennies, deux grandes puissances asiatiques sont devenues des partenaires incontournables du Cameroun : la \cle{Chine}, visible partout à travers ses grands chantiers, et le \cle{Japon}, plus discret mais très actif dans le développement humain. Qui n'a jamais vu, à Yaoundé ou à Douala, un engin de chantier portant un nom chinois, ou acheté un produit estampillé \emph{Made in China} ? Cette présence quotidienne est le signe d'une coopération bilatérale intense qu'il faut comprendre, analyser et évaluer avec un regard critique.

\courssub{L'essor de la coopération sino-camerounaise}

\textbf{Intuition et point de départ.} Imagine un jeune technicien de Kribi qui travaille sur le chantier du grand port : son employeur est une entreprise chinoise, le financement vient d'une banque chinoise, et les grues portent des inscriptions en mandarin. Comment un pays situé à plus de \SI{10000}{km} est-il devenu l'un des premiers bâtisseurs du Cameroun ? C'est toute l'histoire de la coopération sino-camerounaise.

\begin{definition}[Coopération sino-camerounaise]
La coopération sino-camerounaise désigne l'ensemble des relations diplomatiques, économiques, commerciales et culturelles établies entre la République du Cameroun et la République populaire de Chine depuis l'établissement des \cle{relations diplomatiques} en \textbf{1971}, quelques années après l'indépendance du Cameroun.
\end{definition}

\textbf{Les étapes de l'essor.} La coopération entre les deux pays s'est construite progressivement :

\begin{enumerate}
\item \textbf{1971} : établissement des relations diplomatiques officielles ; le Cameroun reconnaît la République populaire de Chine.
\item \textbf{Années 1970-1980} : premières réalisations emblématiques offertes ou financées par la Chine, notamment le \cle{Palais des Congrès de Yaoundé} et le \cle{Palais des Sports} de Yaoundé, ainsi que des hôpitaux (hôpital de Mbalmayo, hôpital gynéco-obstétrique et pédiatrique de Yaoundé).
\item \textbf{Années 1990-2000} : montée en puissance des échanges commerciaux ; les produits manufacturés chinois inondent les marchés camerounais.
\item \textbf{Depuis les années 2000} : la Chine devient le partenaire des \cle{grandes infrastructures} : stades, routes, barrages hydroélectriques, port en eau profonde de Kribi. Des visites officielles de haut niveau (comme celle du président chinois Hu Jintao en 2007) scellent cette relation privilégiée.
\end{enumerate}

\begin{savaistu}[La Chine, un partenaire désormais incontournable]
La Chine est devenue l'un des tout premiers partenaires commerciaux du Cameroun, concurrençant désormais la \textbf{France}, partenaire historique depuis l'indépendance. En quelques décennies, les produits chinois (appareils électroménagers, téléphones, motos, tissus, matériaux de construction) sont devenus omniprésents dans les marchés et les foyers camerounais, au point que l'expression \emph{made in China} fait partie du langage courant.
\end{savaistu}

\courssub{Les grandes réalisations chinoises au Cameroun}

La coopération chinoise se reconnaît d'abord à ses ouvrages, visibles dans tout le pays. On peut les classer en plusieurs catégories :

\begin{itemize}
\item \textbf{Équipements sportifs et culturels} : le Palais des Congrès et le Palais des Sports de Yaoundé ; les stades modernes de \cle{Japoma} (Douala) et de \cle{Bafoussam}, construits par des entreprises chinoises notamment pour la Coupe d'Afrique des Nations, auxquels s'ajoutent d'autres équipements sportifs réalisés dans le cadre de cette compétition (comme le stade d'Olembe à Yaoundé).
\item \textbf{Infrastructures de transport} : de nombreux tronçons routiers (routes nationales, échangeurs urbains) et surtout le \cle{port en eau profonde de Kribi}, capable d'accueillir de très grands navires.
\item \textbf{Énergie} : les barrages hydroélectriques de \cle{Mekin} (sur le Dja, région du Sud) et de \cle{Memve'ele} (sur le Ntem, région du Sud), qui augmentent la production d'électricité du pays.
\item \textbf{Santé} : construction et équipement d'hôpitaux, envoi de missions médicales chinoises qui soignent des populations.
\end{itemize}

\begin{exemplebox}[Le port en eau profonde de Kribi]
Le port en eau profonde de Kribi, inauguré en 2018, a été financé majoritairement par un \textbf{prêt chinois} de la banque \emph{Eximbank of China} et construit par l'entreprise d'État chinoise \textbf{CHEC} (China Harbour Engineering Company). Il est devenu l'un des principaux \cle{hubs portuaires} d'Afrique centrale : il désengorge le port de Douala, accueille les plus grands porte-conteneurs et constitue le cœur d'une zone industrielle en développement. C'est l'exemple type de la coopération chinoise : un grand ouvrage, un financement par prêt, une entreprise chinoise comme constructeur.
\end{exemplebox}

\begin{popfigure}
\begin{center}
\begin{tikzpicture}[scale=1.0, every node/.style={font=\small}]
% Contour très simplifié du Cameroun
\draw[thick, popdark, fill=popcream] (1.2,0) -- (3.4,0) -- (4.6,0.8) -- (4.4,2.2) -- (5.2,3.2) -- (4.2,4.6) -- (3.4,5.6) -- (2.6,6.6) -- (1.8,6.2) -- (1.2,5.0) -- (0.4,4.2) -- (0.8,3.0) -- (0.2,1.8) -- cycle;
% Flèche du nord
\draw[-{Stealth}, thick, popdark] (5.6,5.6) -- (5.6,6.4) node[above]{N};
% Villes et projets
\fill[popblue] (1.0,1.2) circle (2.2pt) node[below left, popdark]{Douala};
\fill[popblue] (2.2,1.6) circle (2.2pt) node[right, popdark]{Yaoundé};
\fill[popink] (0.6,0.5) circle (2.6pt) node[below, popdark]{Kribi : port en eau profonde};
\fill[poporange] (2.9,1.0) circle (2.6pt) node[right, popdark]{Mekin / Memve'ele (barrages)};
\fill[poppurple] (1.3,2.6) circle (2.6pt) node[left, popdark]{Bafoussam : stade};
\fill[poppurple] (1.35,1.55) circle (2.6pt);
\node[popdark] at (0.35,2.0) {Japoma};
\fill[popgreen] (2.2,2.6) circle (2.6pt) node[right, popdark]{Projets JICA (riz, santé, formation)};
\end{tikzpicture}
\end{center}
\legende{Carte schématique du Cameroun localisant les grands projets de la coopération asiatique : en violet, les stades construits avec la Chine ; en orange, les barrages hydroélectriques ; en rouge, le port de Kribi ; en vert, les zones d'appui japonais (JICA). Croquis approximatif, sans prétention de précision cartographique.}
\end{popfigure}

\courssub{Les échanges commerciaux et les instruments chinois}

\textbf{Un commerce déséquilibré.} Les échanges commerciaux sino-camerounais reposent sur un schéma classique entre un pays industrialisé et un pays en développement :

\begin{itemize}
\item le Cameroun \textbf{importe} de Chine des \cle{produits manufacturés} : appareils électroniques, motos, textiles, chaussures, matériaux de construction, machines ;
\item le Cameroun \textbf{exporte} vers la Chine des \cle{matières premières} : bois, pétrole brut, coton, minerais.
\end{itemize}

Ce déséquilibre (le Cameroun vend des produits bruts peu transformés et achète des produits finis à forte valeur ajoutée) est une limite importante de cette coopération.

\begin{definition}[Prêt concessionnel]
Un \cle{prêt concessionnel} est un prêt accordé à des conditions avantageuses (faible taux d'intérêt, longue durée de remboursement, parfois un différé de plusieurs années avant le premier remboursement) par un État ou une institution à un pays en développement. Il reste néanmoins un prêt : il doit être \textbf{remboursé}.
\end{definition}

\textbf{Les instruments de la coopération chinoise.} La Chine utilise plusieurs outils :

\begin{enumerate}
\item les \textbf{prêts concessionnels} accordés par des banques d'État comme \emph{Eximbank of China} pour financer les grands ouvrages ;
\item les \textbf{dons} : certains bâtiments (Palais des Congrès, équipements hospitaliers) ont été offerts ;
\item les \textbf{entreprises d'État chinoises} comme \textbf{CHEC} (génie portuaire) ou \textbf{Sinohydro} (barrages), qui réalisent les chantiers avec une main-d'œuvre en partie chinoise ;
\item les \textbf{annulations de dettes} : la Chine a plusieurs fois effacé une partie des dettes du Cameroun, geste très médiatisé qui renforce l'amitié entre les deux pays.
\end{enumerate}

\begin{attention}[La coopération chinoise n'est pas une aide gratuite]
Piège fréquent dans les copies : présenter la coopération chinoise comme une aide gratuite ou une générosité désintéressée. En réalité, la plupart des grands projets sont financés par des \textbf{prêts à rembourser}, souvent garantis par les ressources naturelles du pays. Cela alimente le débat sur l'\cle{endettement} du Cameroun : les infrastructures sont utiles, mais leur remboursement pèse sur le budget de l'État pendant des décennies. Au Baccalauréat technique, il faut savoir présenter les deux faces : les avantages (infrastructures modernes, prix compétitifs, rapidité d'exécution) et les limites (dette, main-d'œuvre étrangère, concurrence aux entreprises locales).
\end{attention}

\courssub{La coopération nippo-camerounaise}

\textbf{Intuition.} Contrairement à la Chine, le Japon ne construit pas de stades géants au Cameroun. Pourtant, dans une rizière de Yagoua, dans un centre de santé de village ou dans une salle de formation professionnelle, on trouve sa trace. Le Japon a choisi un autre visage de la coopération : celui du \cle{développement humain} et technique.

\begin{definition}[La JICA]
La \cle{JICA} (\emph{Japan International Cooperation Agency}) est l'agence japonaise de coopération internationale. Elle met en œuvre l'aide officielle du Japon aux pays en développement : projets techniques, envoi d'experts et de volontaires, dons d'équipements, bourses de formation.
\end{definition}

\textbf{Les domaines d'intervention du Japon au Cameroun :}

\begin{itemize}
\item \textbf{Agriculture} : appui à la \cle{riziculture} (amélioration des rendements, formation des producteurs, périmètres irrigués dans le Nord et l'Extrême-Nord), ce qui aide le Cameroun à réduire ses importations de riz ;
\item \textbf{Santé} : dons d'équipements médicaux, appui à la vaccination et à la lutte contre les épidémies ;
\item \textbf{Éducation et formation} : construction de salles de classe, appui à la \cle{formation professionnelle} (domaine qui concerne directement les élèves de l'enseignement technique), envoi de volontaires japonais ;
\item \textbf{Bourses d'études} : chaque année, des étudiants et des fonctionnaires camerounais partent se former au Japon grâce à des bourses du gouvernement japonais ;
\item \textbf{Développement rural} : projets d'adduction d'eau, petits équipements communautaires, appui aux groupements de producteurs.
\end{itemize}

\begin{exemplebox}[Deux styles de coopération]
Pour réhabiliter la riziculture dans la plaine des Mbo, la JICA envoie des experts japonais qui forment les agriculteurs camerounais aux techniques de repiquage et de gestion de l'eau : l'objectif est de transmettre un \textbf{savoir-faire durable}. À la même époque, à quelques centaines de kilomètres, une entreprise chinoise construit un barrage financé par un prêt : l'objectif est de livrer un \textbf{ouvrage}. Les deux approches sont complémentaires mais ne répondent pas aux mêmes besoins.
\end{exemplebox}

\courssub{Bilan comparé et illustration statistique}

\begin{popfigure}
\begin{center}
\begin{tikzpicture}[scale=1.0]
% Diagramme circulaire simplifié : 30 % = 108 degres, 20 % = 72 degres
\draw[fill=popblue!70] (0,0) -- (0:2.2) arc (0:108:2.2) -- cycle;
\draw[fill=poporange!80] (0,0) -- (108:2.2) arc (108:180:2.2) -- cycle;
\draw[fill=popgreen!70] (0,0) -- (180:2.2) arc (180:252:2.2) -- cycle;
\draw[fill=poppurple!70] (0,0) -- (252:2.2) arc (252:360:2.2) -- cycle;
% Légende
\node[right, font=\small] at (3.0,1.5) {\textcolor{popblue}{\rule{10pt}{10pt}} Chine : environ 30 \%};
\node[right, font=\small] at (3.0,0.9) {\textcolor{poporange}{\rule{10pt}{10pt}} France : environ 20 \%};
\node[right, font=\small] at (3.0,0.3) {\textcolor{popgreen}{\rule{10pt}{10pt}} Autres partenaires UE : environ 20 \%};
\node[right, font=\small] at (3.0,-0.3) {\textcolor{poppurple}{\rule{10pt}{10pt}} Reste du monde : environ 30 \%};
\end{tikzpicture}
\end{center}
\legende{Répartition indicative des principaux partenaires dans les importations du Cameroun (ordres de grandeur, sources : données commerciales récentes). La Chine occupe désormais la première place, devant la France, partenaire historique : c'est le signe de la montée en puissance asiatique dans le commerce extérieur camerounais.}
\end{popfigure}

\begin{exoresolu}[Analyser et comparer deux partenariats]
\textbf{Énoncé.} À partir de tes connaissances, présente en quelques lignes la différence entre la coopération du Cameroun avec la Chine et celle avec le Japon, puis dégage un avantage et une limite de la coopération sino-camerounaise.
\tcblower
\textbf{Solution détaillée.}
\begin{enumerate}
\item \emph{Différence} : la Chine est le partenaire des \textbf{grandes infrastructures} (port de Kribi, barrages de Mekin et Memve'ele, stades de Japoma et de Bafoussam, routes), financées surtout par des prêts et réalisées par des entreprises d'État chinoises. Le Japon, à travers la JICA, est le partenaire du \textbf{développement humain et technique} : riziculture, santé, éducation, formation professionnelle, bourses d'études.
\item \emph{Avantage de la coopération chinoise} : elle fournit rapidement au Cameroun des infrastructures modernes dont le pays avait besoin (énergie, transport, équipements sportifs).
\item \emph{Limite} : ces projets sont financés par des prêts qui alourdissent la dette publique, et les chantiers emploient souvent une main-d'œuvre et des entreprises étrangères, ce qui profite moins à l'économie locale.
\end{enumerate}
\emph{Conclusion} : les deux coopérations sont complémentaires ; le citoyen éclairé doit savoir en apprécier les apports tout en restant critique sur leurs conditions.
\end{exoresolu}

\begin{aretenir}[L'essentiel de la section]
\begin{itemize}
\item Les relations diplomatiques entre le Cameroun et la Chine datent de \textbf{1971} ; la Chine est devenue l'un des premiers partenaires commerciaux du Cameroun.
\item Réalisations chinoises majeures : Palais des Congrès, Palais des Sports, stades (Japoma, Bafoussam), port en eau profonde de \textbf{Kribi}, barrages de \textbf{Mekin} et \textbf{Memve'ele}.
\item Instruments chinois : prêts concessionnels (Eximbank), dons, entreprises d'État (CHEC, Sinohydro), annulations de dettes.
\item Le Japon agit surtout à travers la \textbf{JICA} : riziculture, santé, éducation, formation professionnelle, bourses.
\item Formule de synthèse : la Chine, partenaire des \textbf{grandes infrastructures} ; le Japon, partenaire du \textbf{développement humain et technique}.
\end{itemize}
\end{aretenir}

\begin{exercice}[Pour vérifier tes acquis]
\begin{enumerate}
\item Définis l'expression \emph{prêt concessionnel} et explique pourquoi il ne s'agit pas d'un don.
\item Cite trois grandes réalisations chinoises au Cameroun et précise leur utilité pour le pays.
\item Présente deux domaines d'intervention de la JICA au Cameroun.
\item Rédige un paragraphe argumenté répondant à la question : la coopération sino-camerounaise est-elle une chance ou un risque pour le Cameroun ?
\end{enumerate}
\end{exercice}

\begin{corrige}[Exercice : Pour vérifier tes acquis]
\begin{enumerate}
\item Un prêt concessionnel est un prêt accordé à des conditions avantageuses (faible taux d'intérêt, longue durée de remboursement) par un État à un pays en développement. Contrairement à un don, il doit être remboursé : il crée donc une dette.
\item Exemples attendus : le port en eau profonde de Kribi (désengorge Douala, hub régional), les barrages de Mekin et Memve'ele (production d'électricité), les stades de Japoma et de Bafoussam (équipements sportifs), le Palais des Congrès (vie politique et culturelle).
\item La JICA intervient notamment dans la riziculture (amélioration des rendements), la santé (dons d'équipements), l'éducation et la formation professionnelle, et l'octroi de bourses d'études au Japon.
\item Paragraphe attendu : thèse nuancée. Chance : infrastructures modernes, financements accessibles, rapidité d'exécution, partenaire commercial majeur. Risque : endettement croissant, dépendance économique, concurrence faite aux entreprises et à la main-d'œuvre locales, exportation de matières premières brutes sans transformation. Conclusion : une coopération utile si l'État camerounais négocie de bonnes conditions et veille au remboursement maîtrisé de la dette.
\end{enumerate}
\end{corrige}

\coursec{Le Cameroun et les pays d'Amérique : USA et Brésil}

Dans la section précédente, nous avons vu comment le Cameroun s'est rapproché des puissances asiatiques : la Chine, devenue le grand bailleur d'infrastructures (routes, barrages, stades), et le Japon, partenaire technique discret mais efficace. Nous traversons maintenant l'océan Atlantique pour étudier deux partenaires américains du Cameroun, très différents l'un de l'autre : les \cle{États-Unis d'Amérique} (USA), superpuissance du Nord, et le \cle{Brésil}, grande puissance émergente du Sud. Comprendre ces deux coopérations, c'est comprendre deux logiques distinctes : la \emph{coopération Nord-Sud}, marquée par l'aide et la sécurité, et la \emph{coopération Sud-Sud}, fondée sur l'échange entre pays en développement.

\courssub{La coopération américano-camerounaise}

\textbf{Des relations diplomatiques anciennes}\par
Les États-Unis ont reconnu le Cameroun dès son indépendance en 1960 et ont ouvert une ambassade à Yaoundé. Les deux pays entretiennent depuis des relations diplomatiques continues, matérialisées par des visites officielles, des accords de coopération et la présence d'une ambassade des États-Unis à Yaoundé et d'une ambassade du Cameroun à Washington. Les États-Unis sont pour le Cameroun un partenaire \cle{stratégique} : leurs interventions touchent la santé, l'éducation, la sécurité et la promotion de la démocratie et des droits humains.

\textbf{Le Peace Corps : des volontaires au service des populations}\par
Le \cle{Peace Corps} est un programme américain qui envoie des volontaires dans les pays en développement. Présent au Cameroun depuis 1962, il déploie de jeunes Américains dans les villages et les lycées, où ils enseignent l'anglais, appuient les activités de santé communautaire et participent à des projets de développement local. Pour un élève de Terminale technique, c'est une réalité concrète : beaucoup ont eu un professeur d'anglais volontaire du Peace Corps dans leur établissement.

\textbf{L'USAID : l'agence américaine d'aide au développement}\par
L'\cle{USAID} (United States Agency for International Development) est l'agence du gouvernement américain chargée de l'aide au développement. Au Cameroun, elle finance de grands programmes, notamment :
\begin{itemize}
\item la \textbf{santé publique} : lutte contre le paludisme (distribution de moustiquaires imprégnées), lutte contre le \cle{VIH/sida} (dépistage, traitements antirétroviraux, notamment via le plan présidentiel américain PEPFAR), vaccination et santé de la mère et de l'enfant ;
\item l'\textbf{éducation} : appui à la lecture et à l'alphabétisation dans les écoles primaires, notamment dans les régions affectées par l'insécurité ;
\item l'\textbf{aide humanitaire} : assistance aux réfugiés et aux déplacés internes de l'Extrême-Nord et des régions anglophones.
\end{itemize}

\begin{exemplebox}[Un exemple chiffré]
L'USAID finance au Cameroun des programmes de santé publique (paludisme, VIH/sida, vaccination) à hauteur de \textbf{dizaines de millions de dollars par an}. À titre indicatif, l'ensemble de l'aide américaine au Cameroun (santé, éducation, aide humanitaire, sécurité) dépasse couramment 100 millions de dollars annuels, ce qui place les États-Unis parmi les premiers bailleurs bilatéraux du pays dans le domaine social.
\end{exemplebox}

\courssub{La coopération militaire et sécuritaire avec les USA}

C'est le volet le plus visible de la coopération américano-camerounaise depuis les années 2010.

\textbf{La lutte contre le terrorisme de Boko Haram}\par
Depuis 2014, la secte terroriste \cle{Boko Haram} multiplie les attentats dans la région de l'Extrême-Nord du Cameroun. Les États-Unis appuient le Cameroun par :
\begin{itemize}
\item la \textbf{formation} des forces de défense camerounaises, notamment le Bataillon d'intervention rapide (BIR) ;
\item la fourniture d'\textbf{équipements} : véhicules, matériel de communication, moyens de surveillance et de renseignement ;
\item l'envoi de \textbf{conseillers militaires} et le partage de renseignements dans le cadre de la lutte commune contre le terrorisme dans le bassin du lac Tchad.
\end{itemize}

\textbf{La lutte contre la piraterie dans le golfe de Guinée}\par
Le \cle{golfe de Guinée} est l'une des zones maritimes les plus dangereuses du monde (piraterie, vols à main armée en mer, trafics). Les États-Unis coopèrent avec la marine camerounaise à travers des exercices militaires conjoints (comme l'exercice \emph{Obangame Express}), des formations et des dons d'équipements pour sécuriser les côtes et le port de Douala, poumon économique du pays.

\begin{attention}[Ne pas réduire la coopération avec les USA au commerce]
La coopération américano-camerounaise est fortement marquée par le \cle{volet sécuritaire} (lutte contre Boko Haram dans l'Extrême-Nord, sécurisation du golfe de Guinée) et par le \cle{volet humanitaire et sanitaire} (USAID, PEPFAR). Au Bac, une erreur fréquente consiste à ne citer que les échanges commerciaux, qui sont en réalité modestes. Pensez toujours à la triade : \textbf{sécurité, santé, éducation}.
\end{attention}

\courssub{La coopération culturelle, éducative et commerciale avec les USA}

\textbf{Les bourses et les programmes d'échange}\par
Les États-Unis développent une coopération culturelle active :
\begin{itemize}
\item les bourses \cle{Fulbright}, qui permettent à des enseignants, chercheurs et étudiants camerounais de poursuivre des études ou des recherches aux États-Unis ;
\item les programmes d'échange comme le \emph{Young African Leaders Initiative} (YALI), qui forme de jeunes leaders africains ;
\item les \cle{English Language Programs} : formation des enseignants d'anglais, American Spaces (bibliothèques et centres culturels américains) à Yaoundé et dans les régions.
\end{itemize}

\textbf{Le programme AGOA}\par
L'\cle{AGOA} (African Growth and Opportunity Act), adopté en 2000, est une loi américaine qui autorise les pays africains éligibles, dont le Cameroun, à exporter certains produits vers le marché américain \textbf{sans droits de douane}. C'est une ouverture commerciale importante : le Cameroun peut exporter notamment des produits agricoles et pétroliers vers les USA. Toutefois, ce dispositif est conditionné au respect de critères (bonne gouvernance, droits humains) ; le Cameroun a perdu son éligibilité en 2020 (décision annoncée fin 2019) avant une procédure de réexamen — preuve que la coopération commerciale avec les USA n'est jamais totalement désintéressée.

\begin{popfigure}
\begin{center}
\begin{tikzpicture}[scale=1,font=\small]
% Océan Atlantique
\fill[popblueL] (0,0) rectangle (10,5.2);
% Amérique du Nord (USA) - bloc simplifié
\fill[popgreen!45] (0,2.6) -- (0,5.2) -- (2.6,5.2) -- (2.4,4.2) -- (1.6,3.4) -- (1.2,2.6) -- cycle;
\node[popdark] at (1.2,4.4) {\textbf{USA}};
\node[popdark,scale=0.8] at (1.2,4.0) {Washington};
% Amérique du Sud (Brésil) - bloc simplifié
\fill[popgold!50] (0,0) -- (0,2.4) -- (1.4,2.2) -- (2.4,1.4) -- (2.2,0.4) -- (1.4,0) -- cycle;
\node[popdark] at (1.0,1.1) {\textbf{Brésil}};
\node[popdark,scale=0.8] at (1.0,0.7) {Brasilia};
% Afrique (Cameroun) - bloc simplifié
\fill[poporange!45] (7.6,0.6) -- (7.2,2.2) -- (7.4,4.4) -- (8.2,5.2) -- (10,5.2) -- (10,0.6) -- cycle;
\node[popdark] at (9.0,4.4) {\textbf{Afrique}};
% Point Cameroun
\fill[popink] (7.55,2.55) circle (0.09);
\node[popdark,anchor=west,scale=0.9] at (7.7,2.55) {\textbf{Cameroun}};
% Flèches d'échanges
\draw[->,line width=2.2pt,popblue] (7.4,2.9) .. controls (5.2,4.6) and (3.4,4.6) .. (2.1,4.1);
\draw[->,line width=2.2pt,popblue,dashed] (2.1,3.9) .. controls (3.6,3.6) and (5.4,3.4) .. (7.3,2.75);
\draw[->,line width=2.2pt,poporange] (7.3,2.3) .. controls (5.4,1.0) and (3.8,0.9) .. (2.5,1.3);
\draw[->,line width=2.2pt,poporange,dashed] (2.6,1.55) .. controls (4.0,1.9) and (5.6,2.0) .. (7.35,2.15);
% Légende des flux
\node[popblue,scale=0.8,align=center] at (4.7,4.9) {aide, sécurité, santé, AGOA};
\node[poporange,scale=0.8,align=center] at (4.7,0.55) {agriculture, élevage, formation technique};
% Nord
\draw[->,thick,popdark] (9.6,0.5) -- (9.6,1.3) node[above,scale=0.8]{N};
% Titre océan
\node[popblue!60!popdark,scale=0.85] at (5.0,2.6) {\emph{Océan Atlantique}};
\end{tikzpicture}
\end{center}
\legende{Carte schématique de l'Atlantique : axes d'échanges Cameroun--USA (flèches bleues, au Nord) et Cameroun--Brésil (flèches oranges, au Sud). Croquis simplifié, sans précision cartographique.}
\end{popfigure}

\courssub{La coopération brésilo-camerounaise : une coopération Sud-Sud}

\textbf{Des liens historiques profonds}\par
Le Brésil et le Cameroun sont séparés par l'Atlantique, mais unis par l'histoire. Pendant la \cle{traite atlantique} (du XVI\textsuperscript{e} au XIX\textsuperscript{e} siècle), des millions d'Africains ont été déportés vers le Brésil, dont beaucoup originaires du golfe de Guinée et des côtes camerounaises. Leurs descendants, les \cle{Afro-Brésiliens}, constituent aujourd'hui une part majeure de la population brésilienne et entretiennent des cultures (musique, religions, cuisine) marquées par cet héritage africain. Cette mémoire commune donne à la coopération brésilo-camerounaise une dimension symbolique forte.

\begin{savaistu}[Le Brésil, « fils » de l'Afrique]
Le Brésil est le pays qui compte la plus grande population d'origine africaine hors d'Afrique. De nombreux Afro-Brésiliens ont des ancêtres originaires du golfe de Guinée. Certaines villes brésiliennes, comme Salvador de Bahia, conservent des traditions (candomblé, capoeira) directement héritées de l'Afrique de l'Ouest et centrale. La coopération actuelle entre Brasilia et Yaoundé renoue ainsi, deux siècles plus tard, des liens que la traite avait brisés.
\end{savaistu}

\begin{definition}[La coopération Sud-Sud]
La \cle{coopération Sud-Sud} désigne les échanges et l'entraide entre \textbf{pays en développement} (pays du « Sud »), par opposition à la coopération Nord-Sud entre pays riches et pays pauvres. Elle repose sur le partage d'expériences et de techniques adaptées à des réalités proches. La coopération entre le Cameroun et le Brésil — deux pays en développement à agriculture tropicale — en est un exemple type.
\end{definition}

\textbf{Les domaines de la coopération avec le Brésil}\par
\begin{itemize}
\item \textbf{L'agriculture et l'agronomie} : le Brésil est devenu une puissance agricole mondiale grâce à sa recherche, incarnée par l'\cle{EMBRAPA} (Entreprise brésilienne de recherche agricole). Le Cameroun bénéficie de l'expertise brésilienne pour améliorer ses cultures (maïs, manioc, canne à sucre) et ses techniques agricoles tropicales.
\item \textbf{L'élevage} : le Brésil est l'un des premiers exportateurs mondiaux de viande bovine ; il appuie le Cameroun dans l'amélioration des races, de l'alimentation animale et de la santé vétérinaire.
\item \textbf{La formation technique et professionnelle} : accueil de cadres, d'enseignants et de techniciens camerounais dans les instituts brésiliens, appui à l'enseignement agricole.
\item \textbf{Les échanges commerciaux} : le Cameroun importe du Brésil du \textbf{sucre}, du \textbf{maïs} et de la \textbf{viande} (notamment de la volaille), tandis que ses exportations vers le Brésil restent modestes.
\end{itemize}

\courssub{Comparaison : deux partenaires, deux logiques}

\begin{center}
\begin{tabularx}{\linewidth}{|l|X|X|}
\hline
\rowcolor{popblueL} \textbf{Critère} & \textbf{États-Unis (USA)} & \textbf{Brésil} \\
\hline
Type de coopération & Nord-Sud (aide au développement) & Sud-Sud (entraide entre pays en développement) \\
\hline
Domaines privilégiés & Sécurité (Boko Haram, piraterie), santé (VIH/sida, paludisme), éducation & Agriculture, élevage, formation technique, commerce \\
\hline
Outils principaux & USAID, Peace Corps, PEPFAR, AGOA, bourses Fulbright & EMBRAPA, formation de techniciens, échanges commerciaux \\
\hline
Réalisations marquantes & Formation du BIR, moustiquaires, antirétroviraux, American Spaces & Appui à l'agronomie tropicale, importations de sucre, maïs, viande \\
\hline
Avantages pour le Cameroun & Sécurisation du territoire, amélioration de la santé, accès au marché américain & Techniques agricoles adaptées au climat tropical, partenariat sans conditionnalité lourde \\
\hline
Statut du partenaire & Partenaire \textbf{stratégique et sécuritaire} & Partenaire \textbf{agricole et technique du « Sud »} \\
\hline
\end{tabularx}
\end{center}

\begin{aretenir}[L'essentiel de la section]
Les États-Unis sont le partenaire \cle{stratégique et sécuritaire} du Cameroun : aide sanitaire massive (USAID, lutte contre le VIH/sida et le paludisme), coopération militaire contre Boko Haram et la piraterie dans le golfe de Guinée, bourses et programmes éducatifs, ouverture commerciale via l'AGOA. Le Brésil est le partenaire \cle{agricole et technique} du Cameroun, dans le cadre d'une coopération \cle{Sud-Sud} nourrie par des liens historiques afro-brésiliens : agronomie (EMBRAPA), élevage, formation technique et commerce (sucre, maïs, viande).
\end{aretenir}

\begin{exoresolu}[Exercice résolu : distinguer les deux coopérations]
\textbf{Énoncé.} Un élève affirme : « La coopération entre le Cameroun et le Brésil est une coopération Nord-Sud, comme celle avec les États-Unis. » Corrige cette affirmation en deux arguments, puis cite deux réalisations concrètes de chaque coopération.
\tcblower
\textbf{Solution détaillée.} L'affirmation est fausse. \textbf{Argument 1} : le Brésil est, comme le Cameroun, un pays en développement ; leurs échanges relèvent donc de la \cle{coopération Sud-Sud} (entraide entre pays du Sud). \textbf{Argument 2} : les États-Unis sont une superpuissance développée ; leur coopération avec le Cameroun est une coopération \cle{Nord-Sud}, fondée sur l'aide au développement. \textbf{Réalisations avec les USA} : financement de la lutte contre le VIH/sida et le paludisme par l'USAID ; formation du BIR contre Boko Haram ; bourses Fulbright. \textbf{Réalisations avec le Brésil} : appui de l'EMBRAPA à l'agronomie tropicale ; importations de sucre, de maïs et de viande ; formation de techniciens agricoles camerounais.
\end{exoresolu}

\begin{exercice}[À faire]
\begin{enumerate}
\item Définis la coopération Sud-Sud et explique pourquoi la coopération Cameroun--Brésil en est un exemple.
\item Cite trois domaines d'intervention de l'USAID au Cameroun.
\item Pourquoi dit-on que la coopération avec les États-Unis ne se limite pas au commerce ? Donne deux arguments.
\end{enumerate}
\end{exercice}

\begin{corrige}[Exercice]
\begin{enumerate}
\item La coopération Sud-Sud est l'entraide entre pays en développement. Le Brésil et le Cameroun sont tous deux des pays du Sud à agriculture tropicale : ils échangent des techniques adaptées (agronomie, élevage) sans rapport de domination Nord-Sud.
\item Santé publique (paludisme, VIH/sida, vaccination), éducation (lecture, alphabétisation), aide humanitaire aux réfugiés et déplacés.
\item Parce qu'elle comprend un volet sécuritaire majeur (lutte contre Boko Haram, sécurisation du golfe de Guinée) et un volet sanitaire et éducatif (USAID, Peace Corps, Fulbright), alors que les échanges commerciaux restent modestes.
\end{enumerate}
\end{corrige}

\coursec{Bilan : avantages, limites et enjeux de la coopération bilatérale}

Après avoir étudié les relations du Cameroun avec la France et la Grande-Bretagne, puis avec la Chine et le Japon, enfin avec les États-Unis et le Brésil, nous pouvons prendre du recul. Toutes ces coopérations bilatérales forment un ensemble cohérent : un \cle{réseau de partenariats} que le Cameroun entretient avec de nombreux pays. Il est temps de dresser le \cle{bilan} : que gagne le Cameroun ? Que gagnent ses partenaires ? Quels sont les dangers ? Et comment le Cameroun peut-il tirer le meilleur parti de ces relations ? C'est à ces questions que répond cette dernière partie, qui te prépare aussi à un type de sujet fréquent au Baccalauréat technique : la dissertation argumentée sur la coopération.

\courssub{Les avantages de la coopération bilatérale pour le Cameroun}

Imaginons d'abord une situation simple. Un jeune technicien de Douala travaille sur le chantier d'un barrage financé par un partenaire étranger ; sa sœur a obtenu une bourse pour étudier à l'étranger ; leur mère se soigne dans un hôpital équipé grâce à un don international. Sans le savoir, cette famille bénéficie chaque jour de la coopération bilatérale. Cette observation concrète nous permet de dégager les principaux avantages pour le Cameroun.

\begin{definition}[Avantages de la coopération bilatérale pour le Cameroun]
La \cle{coopération bilatérale} apporte au Cameroun des ressources financières, techniques et humaines qu'il ne pourrait pas mobiliser seul à court terme : financement des infrastructures, formation des cadres, transferts de technologie, appuis sanitaires et éducatifs, ouverture sur le monde.
\end{definition}

Détaillons ces avantages :

\begin{itemize}
\item \textbf{Le financement des infrastructures.} Les routes, ponts, ports, barrages et stades coûtent très cher. Grâce à des prêts et des dons de partenaires comme la Chine ou la France, le Cameroun a pu réaliser de grands projets : le port en eau profonde de Kribi, le barrage hydroélectrique de Memve'ele sur la Ntem, l'autopont de la Wouri à Douala, les stades construits ou rénovés pour la CAN 2021 (Olembe, Japoma, Bafoussam, Limbé).
\item \textbf{La formation des cadres.} Les partenaires offrent des \cle{bourses d'études} et des stages. Des milliers de Camerounais ont été formés en France, en Chine, au Japon, aux États-Unis ou en Grande-Bretagne, dans des domaines comme la médecine, l'ingénierie ou l'agriculture.
\item \textbf{Les transferts de technologie.} Les entreprises étrangères installées au Cameroun apportent des machines, des méthodes de travail et des savoir-faire que les techniciens camerounais apprennent sur le tas. C'est le \cle{transfert de technologie}.
\item \textbf{Les appuis sanitaires et éducatifs.} Construction d'hôpitaux et d'écoles, campagnes de vaccination, dons de matériel médical et scolaire : les partenaires soutiennent les secteurs sociaux.
\item \textbf{L'ouverture sur le monde.} La coopération donne au Cameroun une visibilité internationale, facilite les échanges culturels et permet à ses produits (cacao, café, bois, pétrole) d'être connus et vendus à l'étranger.
\end{itemize}

\begin{exemplebox}[Un avantage visible au quotidien]
Le port en eau profonde de Kribi, construit avec la coopération chinoise et inauguré en 2018, permet aujourd'hui d'accueillir de grands navires que le port de Douala ne pouvait pas recevoir. Il crée des emplois (dockers, techniciens, douaniers), fait baisser le coût des importations et dynamise toute la région du Sud. C'est un exemple concret d'infrastructure issue de la coopération bilatérale utilisée quotidiennement.
\end{exemplebox}

\courssub{Les avantages pour les pays partenaires}

Une erreur de raisonnement serait de croire que seul le Cameroun profite de la coopération. Observons les faits : pourquoi des pays riches consacrent-ils des milliards à la coopération avec le Cameroun ? Parce qu'ils y trouvent aussi leur intérêt. La coopération est un \emph{échange}, pas un cadeau à sens unique.

\begin{itemize}
\item \textbf{Des débouchés commerciaux.} Le Cameroun est un marché de près de 30 millions de consommateurs. Les entreprises françaises, chinoises ou américaines y vendent leurs produits : voitures, machines, médicaments, téléphones, matériaux de construction.
\item \textbf{L'accès aux matières premières.} Le Cameroun exporte vers ses partenaires le \cle{pétrole}, le \cle{bois}, le \cle{cacao}, le café, le coton et l'aluminium. Ces ressources sont indispensables aux industries des pays partenaires.
\item \textbf{L'influence diplomatique et culturelle.} En aidant le Cameroun, un partenaire gagne un ami qui peut soutenir ses positions dans les organisations internationales (ONU, Francophonie, Commonwealth). La diffusion de la langue et de la culture (Instituts français, British Council, Instituts Confucius) renforce cette influence, appelée \emph{soft power}.
\end{itemize}

\begin{attention}[Erreur fréquente au Bac]
Ne jamais affirmer que la coopération bilatérale ne profite qu'au Cameroun. Chaque partenaire y trouve aussi des intérêts : des marchés pour ses produits, l'accès aux matières premières (pétrole, bois, cacao) et une influence diplomatique et culturelle. Une bonne copie présente toujours les deux côtés de l'échange.
\end{attention}

\courssub{Les limites de la coopération bilatérale}

Pour être honnête, le bilan comporte aussi des zones d'ombre. Reprenons notre démarche d'observation : si la coopération était parfaite, le Cameroun serait déjà un pays émergent. Or des difficultés persistent. Analysons-les.

\begin{itemize}
\item \textbf{L'endettement.} Beaucoup de projets sont financés par des \cle{prêts} qu'il faut rembourser avec des intérêts. La dette extérieure du Cameroun a fortement augmenté, ce qui pèse sur le budget de l'État : une partie des recettes publiques sert à rembourser au lieu de financer les écoles ou les hôpitaux.
\item \textbf{La dépendance économique.} Quand un pays dépend trop d'un partenaire pour ses financements, ses équipements ou ses débouchés, il perd une partie de sa liberté de décision. C'est la \cle{dépendance économique}.
\item \textbf{La concurrence des produits importés.} Les produits étrangers bon marché (riz, huile, tissus, poulets, appareils) concurrencent la production locale et peuvent décourager les producteurs et artisans camerounais.
\item \textbf{Des conditions parfois désavantageuses.} Certains contrats prévoient que les travaux sont réalisés par les entreprises et la main-d'œuvre du partenaire, ou que les matières premières sont exportées brutes, sans transformation locale. Le Cameroun gagne alors moins que ce qu'il pourrait gagner.
\end{itemize}

\begin{propriete}[Condition d'une coopération réussie]
Une coopération bilatérale équilibrée doit être \cle{mutuellement avantageuse} (gagnant-gagnant). Pour le Cameroun, l'enjeu est de transformer l'aide et les prêts en \cle{investissements productifs} : usines, infrastructures rentables, transformation locale des matières premières, afin de créer des richesses durables et des emplois plutôt que de la dette.
\end{propriete}

\courssub{La diversification des partenaires : une stratégie d'ouverture}

Face au risque de dépendance, le Cameroun a adopté une stratégie claire : ne pas mettre tous ses œufs dans le même panier. C'est la \cle{diversification des partenaires}.

\begin{definition}[Diversification des partenaires]
La \cle{diversification des partenaires} est la stratégie qui consiste, pour un pays, à multiplier ses partenaires de coopération afin de ne dépendre d'aucun partenaire unique et de pouvoir négocier de meilleures conditions.
\end{definition}

Longtemps tourné vers la France et les pays occidentaux, le Cameroun s'est ouvert depuis les années 2000 à de nouveaux partenaires : la \textbf{Chine} (infrastructures), le \textbf{Japon} (appui technique), les \textbf{États-Unis} et le \textbf{Brésil} (agriculture, formation), mais aussi la \textbf{Turquie} (construction, éducation) et l'\textbf{Inde} (technologies, santé). Cette ouverture permet au Cameroun de comparer les offres, de choisir les financements les plus avantageux et de renforcer son indépendance.

\begin{savaistu}[Le Cameroun émergent à l'horizon 2035]
La stratégie de diversification des partenaires s'inscrit dans la vision « Cameroun émergent à l'horizon 2035 », déclinée dans la Stratégie Nationale de Développement 2020-2030 (SND30), qui a succédé au Document de Stratégie pour la Croissance et l'Emploi (DSCE). L'objectif : faire du Cameroun un pays émergent, démocratique et uni dans sa diversité, en s'appuyant sur des partenariats internationaux équilibrés et sur la transformation locale des richesses nationales.
\end{savaistu}

\begin{popfigure}
\begin{center}
\begin{tikzpicture}[scale=1]
% Cameroun au centre
\node[circle,draw=popdark,fill=popgoldL,thick,minimum size=2.2cm,align=center,font=\bfseries] (cam) at (0,0) {CAMEROUN};
% Partenaires
\node[rectangle,rounded corners,draw=popblue,fill=popblueL,thick,minimum width=2.1cm,align=center] (fr) at (-4,2.2) {France};
\node[rectangle,rounded corners,draw=popblue,fill=popblueL,thick,minimum width=2.1cm,align=center] (gb) at (-4,-2.2) {Grande-\\Bretagne};
\node[rectangle,rounded corners,draw=poppink,fill=poppinkL,thick,minimum width=2.1cm,align=center] (cn) at (4,2.6) {Chine};
\node[rectangle,rounded corners,draw=poppink,fill=poppinkL,thick,minimum width=2.1cm,align=center] (jp) at (4,0.6) {Japon};
\node[rectangle,rounded corners,draw=popgreen,fill=popgreenL,thick,minimum width=2.1cm,align=center] (us) at (4,-1.4) {USA};
\node[rectangle,rounded corners,draw=popgreen,fill=popgreenL,thick,minimum width=2.1cm,align=center] (br) at (4,-3.2) {Brésil};
\node[rectangle,rounded corners,draw=poporange,fill=poporangeL,thick,minimum width=2.6cm,align=center] (div) at (-4,0) {Turquie, Inde,\\Corée...};
% Flèches aller-retour
\draw[<->,>=stealth,very thick,popblue] (cam) -- (fr);
\draw[<->,>=stealth,very thick,popblue] (cam) -- (gb);
\draw[<->,>=stealth,very thick,poppink] (cam) -- (cn);
\draw[<->,>=stealth,very thick,poppink] (cam) -- (jp);
\draw[<->,>=stealth,very thick,popgreen] (cam) -- (us);
\draw[<->,>=stealth,very thick,popgreen] (cam) -- (br);
\draw[<->,>=stealth,very thick,poporange] (cam) -- (div);
% Légende des flux
\node[align=left,font=\small,fill=popcream,draw=popline,rounded corners,inner sep=4pt] at (0,-4.6) {Vers l'extérieur : pétrole, bois, cacao, café, coton \quad / \quad Vers le Cameroun : financements, technologies, bourses, appuis};
\end{tikzpicture}
\end{center}
\legende{Croquis de synthèse : le Cameroun au centre d'un réseau de partenaires. Les doubles flèches symbolisent les échanges aller-retour : le Cameroun exporte ses matières premières et reçoit en retour financements, technologies, formations et appuis. La diversification (Turquie, Inde, Corée...) complète les partenaires traditionnels.}
\end{popfigure}

\courssub{Tableau-bilan : forces et faiblesses de la coopération bilatérale}

\begin{popfigure}
\begin{center}
\begin{tabularx}{\linewidth}{|X|X|}
\hline
\rowcolor{popgreenL} \textbf{Avantages (forces)} & \textbf{Limites (faiblesses)} \\
\hline
Financement des infrastructures (port de Kribi, barrages, routes, stades) & Endettement croissant de l'État \\
\hline
Formation des cadres (bourses, stages à l'étranger) & Dépendance économique vis-à-vis des partenaires \\
\hline
Transferts de technologie et emplois créés par les entreprises étrangères & Concurrence des produits importés à la production locale \\
\hline
Appuis sanitaires et éducatifs (hôpitaux, écoles, vaccinations) & Conditions parfois désavantageuses (matières premières exportées brutes) \\
\hline
Ouverture sur le monde et influence internationale du Cameroun & Risque de perte de souveraineté dans les décisions \\
\hline
\rowcolor{popblueL} \multicolumn{2}{|l|}{\textbf{Intérêts des partenaires :} débouchés commerciaux, matières premières (pétrole, bois, cacao), influence diplomatique et culturelle.} \\
\hline
\end{tabularx}
\end{center}
\legende{Tableau-bilan des forces et faiblesses de la coopération bilatérale du Cameroun, avec le rappel des intérêts des partenaires.}
\end{popfigure}

\courssub{La coopération bilatérale dans ta vie quotidienne}

La coopération n'est pas une affaire lointaine de diplomates : elle touche directement la vie des jeunes techniciens camerounais.

\begin{itemize}
\item \textbf{Les emplois.} Les entreprises étrangères (BTP, télécommunications, agro-industrie) emploient des milliers de techniciens camerounais : conducteurs d'engins, électriciens, informaticiens, comptables.
\item \textbf{Les bourses d'études.} Chaque année, des concours de bourses permettent à des élèves de poursuivre leurs études en Chine, en Russie, en Turquie, au Maroc ou en France.
\item \textbf{Les infrastructures utilisées chaque jour.} Le pont que tu traverses, la route qui mène à ton lycée, l'hôpital où l'on te soigne, le réseau électrique : beaucoup de ces équipements existent grâce à la coopération bilatérale.
\end{itemize}

\begin{exoresolu}[Exercice résolu : analyser une situation concrète]
\textbf{Énoncé.} Une entreprise chinoise construit une route entre deux villes camerounaises, avec un prêt accordé par la Chine au Cameroun. Elle emploie 300 ouvriers camerounais. Une fois achevée, la route permet aux planteurs de la zone d'écouler leur cacao vers le port de Douala. Identifie deux avantages et une limite de cette coopération pour le Cameroun.
\tcblower
\textbf{Solution détaillée.}
\begin{enumerate}
\item \emph{Avantage 1 — infrastructure et développement :} la route désenclave la zone et permet aux planteurs d'exporter leur cacao, ce qui augmente leurs revenus et les recettes d'exportation du pays.
\item \emph{Avantage 2 — emplois et formation :} 300 ouvriers camerounais sont employés et acquièrent des savoir-faire techniques (transfert de technologie).
\item \emph{Limite — endettement :} le projet est financé par un prêt que l'État camerounais devra rembourser avec intérêts ; si la route ne génère pas assez de richesses, la dette pèsera sur le budget national.
\end{enumerate}
\textbf{Conclusion :} cette coopération est utile à condition que l'infrastructure devienne un investissement productif rentable.
\end{exoresolu}

\courssub{Synthèse guidée : rédiger une conclusion argumentée}

Au Baccalauréat technique, on te demandera souvent de conclure un devoir par une prise de position argumentée. Voici la méthode, appliquée à la question : \emph{« La coopération bilatérale est-elle une chance pour le Cameroun ? »}

\begin{methode}[Rédiger une conclusion argumentée au Bac]
\begin{enumerate}
\item \textbf{Rappeler la thèse} : reformuler la question et annoncer ta position (ex. : la coopération bilatérale est globalement une chance pour le Cameroun, à certaines conditions).
\item \textbf{Développer deux arguments étayés d'exemples camerounais} : par exemple le financement des infrastructures (port de Kribi) et la formation des cadres (bourses).
\item \textbf{Nuance (limites)} : reconnaître les dangers (endettement, dépendance, concurrence des importations) pour montrer un esprit critique.
\item \textbf{Ouverture} : élargir la réflexion (ex. : l'avenir dépend de la capacité du Cameroun à transformer l'aide en investissement productif, dans la perspective de l'émergence en 2035).
\end{enumerate}
\end{methode}

\begin{exemplebox}[Modèle de conclusion argumentée]
« En définitive, la coopération bilatérale apparaît comme une chance pour le Cameroun. Elle a d'abord permis le financement d'infrastructures majeures, comme le port en eau profonde de Kribi réalisé avec la Chine, qui dynamise le commerce national. Elle assure ensuite la formation de milliers de cadres grâce aux bourses offertes par la France, le Japon ou les États-Unis. Toutefois, cette coopération comporte des limites : l'endettement croissant, la dépendance économique et la concurrence des produits importés fragilisent l'économie locale. La coopération n'est donc une chance que si elle est équilibrée et mutuellement avantageuse. C'est pourquoi le Cameroun diversifie ses partenaires et cherche à transformer l'aide en investissement productif, condition essentielle pour atteindre l'émergence à l'horizon 2035. »
\end{exemplebox}

\begin{aretenir}[L'essentiel du bilan]
\begin{itemize}
\item La coopération bilatérale apporte au Cameroun financements, infrastructures, formation, technologies et appuis sociaux.
\item Les partenaires y gagnent des marchés, des matières premières (pétrole, bois, cacao) et une influence diplomatique et culturelle.
\item Ses limites : endettement, dépendance, concurrence des importations, conditions parfois désavantageuses.
\item La diversification des partenaires (Chine, Turquie, Inde...) évite la dépendance à un partenaire unique.
\item Une coopération réussie est mutuellement avantageuse et transforme l'aide en investissement productif, dans la vision « Cameroun émergent à l'horizon 2035 » (DSCE/SND30).
\end{itemize}
\end{aretenir}

\begin{exercice}[Pour t'entraîner]
\textbf{Exercice 1.} Définis les expressions suivantes : coopération bilatérale ; diversification des partenaires ; transfert de technologie.

\textbf{Exercice 2.} Cite trois avantages de la coopération bilatérale pour le Cameroun et deux avantages pour les pays partenaires.

\textbf{Exercice 3.} Rédige une conclusion argumentée de dix lignes répondant à la question : « La coopération bilatérale est-elle une chance pour le Cameroun ? » en suivant le plan : thèse, deux arguments avec exemples, nuance, ouverture.
\end{exercice}

\begin{corrige}[Exercice 1]
La \cle{coopération bilatérale} est l'ensemble des relations d'aide et d'échanges (financiers, techniques, culturels, militaires) entre deux États. La \cle{diversification des partenaires} est la stratégie qui consiste à multiplier les partenaires de coopération pour ne dépendre d'aucun partenaire unique. Le \cle{transfert de technologie} est l'acquisition par les nationaux des savoir-faire, machines et méthodes apportés par les partenaires étrangers.
\end{corrige}

\begin{corrige}[Exercice 2]
Pour le Cameroun (trois au choix) : financement des infrastructures (port de Kribi, barrages), formation des cadres (bourses d'études), transferts de technologie, appuis sanitaires et éducatifs, ouverture sur le monde. Pour les partenaires (deux au choix) : débouchés commerciaux pour leurs entreprises, accès aux matières premières (pétrole, bois, cacao), influence diplomatique et culturelle.
\end{corrige}

\begin{corrige}[Exercice 3]
Barème indicatif : \textbf{Thèse} (position claire sur la question) ; \textbf{argument 1} avec exemple camerounais précis (ex. : infrastructures financées par la Chine) ; \textbf{argument 2} avec exemple (ex. : bourses de formation, appuis sanitaires) ; \textbf{nuance} (endettement, dépendance, concurrence des importations) ; \textbf{ouverture} (diversification des partenaires, transformation de l'aide en investissement productif, émergence 2035). La copie doit être rédigée, cohérente et justifiée.
\end{corrige}

\newpage

\coursec{Activité d'intégration}

\coursec{Activité d'intégration : Le Cameroun dans la coopération bilatérale}

\courssub{Objectifs de l'activité}
\noindent Cette activité d'intégration vise à :
\begin{itemize}
    \item Mobiliser les connaissances sur les six partenaires bilatéraux majeurs du Cameroun (France, Grande-Bretagne, Chine, Japon, USA, Brésil) pour produire une synthèse graphique et tabulaire.
    \item Développer la compétence « \cle{Analyser une situation géographique et politique} » en construisant un croquis commenté respectant les règles de la sémiologie graphique (flèches proportionnelles, légende organisée).
    \item Exercer la compétence « \cle{Argumenter et prendre position} » en rédigeant une note de cadrage argumentée à destination d'un décideur public (le MINREX), en s'appuyant sur des données chiffrées et des réalisations concrètes.
    \item Travailler en équipe (collaboration, répartition des tâches, restitution orale) puis confronter les productions lors d'une correction collective guidée.
\end{itemize}

\courssub{Situation}
\noindent \textbf{Contexte authentique :} Le Ministère des Relations Extérieures (MINREX) prépare la visite d'État d'un partenaire stratégique. L'objectif affiché est la signature d'un accord-cadre pour le financement et l'équipement d'un \textbf{nouveau lycée technique industriel de référence} (filières : maintenance industrielle, génie civil, énergies renouvelables) à \textbf{Maroua} (région de l'Extrême-Nord), zone prioritaire pour la décentralisation et la lutte contre la déscolarisation.

\noindent En tant qu'analystes juniors au Département de la Coopération Bilatérale, vous recevez un \textbf{dossier documentaire} composé de :
\begin{enumerate}
    \item Une \textbf{carte du monde} centrée sur le Cameroun.
    \item Un \textbf{tableau de données chiffrées} (ordres de grandeur indicatifs, sources : INS, Douanes camerounaises, ambassades, rapports de coopération, 2022-2023) :
    \begin{center}
    \begin{tabularx}{\linewidth}{|l|X|X|X|X|}\hline
    \rowcolor{popblueL}
    \textbf{Partenaire} & \textbf{Échanges commerciaux (Mrd FCFA)} & \textbf{Investissements directs (stock, Mrd FCFA)} & \textbf{Coopération technique / éducation (projets récents)} & \textbf{Domaines prioritaires déclarés} \\ \hline
    France & 1 250 (1\textsuperscript{er} fournisseur historique, client majeur) & 1 800 & Lycées bilingues (Yaoundé, Douala, Garoua), AFD : formation professionnelle, Campus France & Éducation, santé, gouvernance, défense, francophonie \\ \hline
    Grande-Bretagne & 320 & 450 & British Council, bourses Chevening, appui au secteur privé agro-industriel & Anglais des affaires, justice, climat, sécurité maritime \\ \hline
    Chine & 2 100 (1\textsuperscript{er} partenaire commercial) & 3 500 & Port de Kribi (phases 1 et 2), barrage de Memve'ele, Palais des Congrès de Yaoundé, stade d'Olembé, lycées techniques (Bertoua, Ebolowa) & Infrastructures, mines, énergie, agriculture, formation professionnelle \\ \hline
    Japon & 180 & 120 & Projet rizicole (SEMRY, Extrême-Nord), réhabilitation du pont sur le Wouri, bourses MEXT, équipements médicaux & Agriculture, infrastructures, santé, ressources humaines, TIC \\ \hline
    USA & 450 & 600 & PEPFAR (santé), AGOA (exportations textile/agriculture), YALI, Peace Corps, appui sécuritaire (BIR) & Santé, sécurité, gouvernance, entrepreneuriat, énergie \\ \hline
    Brésil & 95 & 80 & Coopération Sud-Sud (Embrapa : coton, soja), bourses PEC-G, bioéthanol, coopération militaire & Agriculture tropicale, bioénergie, défense, éducation technique \\ \hline
    \end{tabularx}
    \end{center}
    \item Trois \textbf{photographies légendées} :
    \begin{itemize}
        \item \textit{Photo A :} Vue aérienne du \textbf{Port en eau profonde de Kribi} (phase 1 opérationnelle, phase 2 en cours), réalisé par \textit{China Harbour Engineering Company (CHEC)}, financé par Exim Bank of China. Légende : « Porte d'entrée du corridor Centre-Afrique ».
        \item \textit{Photo B :} \textbf{Palais des Congrès de Yaoundé}, construit puis rénové avec l'appui chinois. Légende : « Diplomatie d'infrastructures, vitrine de la coopération sino-camerounaise ».
        \item \textit{Photo C :} \textbf{Lycée bilingue de Nkolbisson (Yaoundé)}, réhabilité via le \textit{Fonds de solidarité prioritaire (FSP) français}. Légende : « Excellence éducative, modèle du bilinguisme officiel, appui à la réforme du système éducatif ».
    \end{itemize}
\end{enumerate}

\noindent \textbf{Situation-problème :} \emph{« Le MINREX doit prioriser un partenaire pour financer le nouveau lycée technique de Maroua. Sur la base du dossier, quel partenaire choisir et pourquoi ? »}

\courssub{Consignes}
\noindent Travail en groupes de 4 à 5 élèves (durée : 2 séances de 50 min + 1 séance de restitution). Chaque groupe rend un dossier unique contenant :
\begin{enumerate}
    \item \textbf{Croquis géographique (page A3 ou A4 paysage) :} « \textit{Le Cameroun et ses partenaires bilatéraux majeurs} ».
        \begin{itemize}
            \item Représenter le Cameroun (position centrale) et les 6 pays partenaires (localisation approximative).
            \item Tracer des \textbf{flèches proportionnelles} (largeur proportionnelle au volume d'échanges commerciaux) orientées Cameroun $\leftrightarrow$ Partenaire (double sens).
            \item Utiliser un \textbf{code couleur} par domaine dominant de coopération (ex. : \textcolor{popblue}{Bleu} = Infrastructures/Énergie ; \textcolor{popgreen}{Vert} = Agriculture/Environnement ; \textcolor{poporange}{Orange} = Éducation/Formation ; \textcolor{poppurple}{Violet} = Santé/Gouvernance ; \textcolor{poppink}{Rose} = Sécurité/Défense).
            \item Légende complète (titre, auteur, date, échelle, orientation, symboles, flèches, couleurs).
        \end{itemize}
    \item \textbf{Tableau comparatif synthétique (format A4) :} Colonnes : Partenaire | Domaines majeurs | 3 Réalisations emblématiques | Avantages pour le Cameroun | Limites / Points de vigilance.
    \item \textbf{Note d'argumentation (15 lignes $\pm$ 2) :} Réponse formelle à la situation-problème. Structure imposée : \textbf{Choix} $\rightarrow$ \textbf{3 arguments chiffrés/factuels} $\rightarrow$ \textbf{Adéquation avec le projet (lycée technique, Maroua, Extrême-Nord)} $\rightarrow$ \textbf{Mesure d'accompagnement proposée} $\rightarrow$ \textbf{Conclusion engageante}.
\end{enumerate}

\courssub{Ressources à mobiliser}
\noindent Rappel des notions-clés du cours (Le Cameroun dans la coopération bilatérale) :
\begin{itemize}
    \item \textbf{Définition :} La \cle{coopération bilatérale} est une relation d'entraide entre deux États souverains, formalisée par des accords (cadres ou spécifiques), reposant sur la réciprocité et l'intérêt mutuel.
    \item \textbf{Typologie des partenaires :}
    \begin{itemize}
        \item \textit{Partenaires historiques / traditionnels :} \textbf{France} (héritage colonial, langue, franc CFA, accords de coopération et de défense, premier investisseur en stock), \textbf{Grande-Bretagne} (zone anglophone, Commonwealth, droit, éducation).
        \item \textit{Partenaires émergents / stratégiques :} \textbf{Chine} (premier partenaire commercial, infrastructures lourdes, financements, Forum FOCAC), \textbf{Japon} (qualité technique, TICAD, appui agricole et sanitaire, dons et prêts concessionnels), \textbf{USA} (puissance économique et militaire, AGOA, santé/sécurité, soft power), \textbf{Brésil} (puissance agricole tropicale, coopération Sud-Sud, Embrapa, similitudes agro-climatiques malgré l'éloignement géographique).
    \end{itemize}
    \item \textbf{Critères d'évaluation d'un partenariat :} volume des échanges, stock d'IDE, alignement sur la \cle{SND30} (Stratégie Nationale de Développement 2020-2030), conditionnalité (gouvernance, droits humains), transfert de technologie, formation des ressources humaines, impact territorial (décentralisation).
    \item \textbf{Outils méthodologiques :} construction d'un croquis (sémiologie graphique), lecture critique de tableaux statistiques, rédaction administrative (note de cadrage, synthèse).
\end{itemize}

\courssub{Démarche et corrigé commenté}
\begin{exoresolu}{Réalisation complète de l'activité (Guide enseignant / Corrigé type)}
\noindent \textbf{Étape 1 : Construction du croquis (sémiologie graphique)}
\medskip
\noindent \textit{Choix de l'échelle des flèches :} volume maximal = Chine (2 100 Mrd FCFA). Échelle proposée : \textbf{1 mm de largeur = 100 Mrd FCFA}.
\begin{itemize}
    \item Chine : 21 mm (trait très épais, double flèche \textcolor{popblue}{Bleu} — Infrastructures/Énergie dominant).
    \item France : 12,5 mm (trait épais, double flèche \textcolor{poporange}{Orange} — Éducation/Formation dominant).
    \item USA : 4,5 mm (trait moyen, double flèche \textcolor{poppurple}{Violet} — Santé/Gouvernance dominant).
    \item Grande-Bretagne : 3,2 mm (trait fin, double flèche \textcolor{poporange}{Orange} — Éducation/Justice).
    \item Japon : 1,8 mm (trait fin, double flèche \textcolor{popgreen}{Vert} — Agriculture dominant).
    \item Brésil : 1 mm (trait très fin, double flèche \textcolor{popgreen}{Vert} — Agriculture/Éducation technique).
\end{itemize}
\noindent \textit{Implantation :} Cameroun au centre (polygone simplifié). Flèches courbes vers : France (nord-ouest), Grande-Bretagne (nord-ouest, au nord de la France), Chine (nord-est, lointaine), Japon (est, encore plus lointain), USA (ouest, traversée de l'Atlantique Nord), Brésil (sud-ouest, traversée de l'Atlantique Sud).
\medskip
\noindent \textit{Vérification de la proportionnalité :} $\dfrac{2100}{1250} \approx 1{,}7$ : la flèche Chine doit être environ 1,7 fois plus large que la flèche France ($21 \approx 1{,}7 \times 12{,}5$) ; $\dfrac{1250}{95} \approx 13$ : la flèche France est environ 13 fois plus large que la flèche Brésil ($12{,}5 \approx 13 \times 1$). La hiérarchie visuelle est donc conforme aux données.
\tcblower
\noindent \textbf{Modèle de croquis attendu (reproduit à main levée par les élèves) :}
\begin{popfigure}
\begin{tikzpicture}[scale=0.62, every node/.style={transform shape}]
% Continents simplifiés (contours très approximatifs)
\fill[popmuted!20] (-4.5,-3.5) -- (-3,2) -- (-3,4.5) -- (-7,4.5) -- (-7,-4.5) -- cycle;
\fill[popmuted!20] (-13,-2) rectangle (-9,2.5);
\fill[popmuted!20] (-12,-6.5) rectangle (-8,-2.5);
\fill[popmuted!20] (-2,5) rectangle (2,8);
\fill[popmuted!20] (6,0.5) rectangle (13,6);
% Cameroun (centré)
\node[circle,fill=popdark,minimum size=7pt,label=below:{\footnotesize\textbf{Cameroun}}] (CM) at (0,0) {};
% Partenaires
\node[circle,fill=popblue,minimum size=5pt,label=above:{\footnotesize\textbf{France}}] (FR) at (-0.5,6) {};
\node[circle,fill=popblue,minimum size=5pt,label=above:{\footnotesize\textbf{R.-U.}}] (UK) at (-1.5,7) {};
\node[circle,fill=poporange,minimum size=5pt,label=above:{\footnotesize\textbf{Chine}}] (CN) at (8,3) {};
\node[circle,fill=poporange,minimum size=5pt,label=right:{\footnotesize\textbf{Japon}}] (JP) at (12,4) {};
\node[circle,fill=popgreen,minimum size=5pt,label=below:{\footnotesize\textbf{USA}}] (US) at (-11,0.5) {};
\node[circle,fill=popgreen,minimum size=5pt,label=below:{\footnotesize\textbf{Brésil}}] (BR) at (-10,-4.5) {};
% Flèches proportionnelles (épaisseur proportionnelle aux échanges)
\draw[popblue, line width=21pt, -{Stealth[length=10pt]}, opacity=0.55] (CM) .. controls (3,1) and (6,2) .. (CN);
\draw[poporange, line width=12.5pt, -{Stealth[length=8pt]}, opacity=0.55] (CM) .. controls (-1,2.5) and (-1,4.5) .. (FR);
\draw[poppurple, line width=4.5pt, -{Stealth[length=6pt]}, opacity=0.6] (CM) .. controls (-4,-0.5) and (-8,0) .. (US);
\draw[poporange, line width=3.2pt, -{Stealth[length=5pt]}, opacity=0.6] (CM) .. controls (-2,3) and (-2,5.5) .. (UK);
\draw[popgreen, line width=1.8pt, -{Stealth[length=4pt]}, opacity=0.7] (CM) .. controls (4,2) and (9,3.5) .. (JP);
\draw[popgreen, line width=1pt, -{Stealth[length=3pt]}, opacity=0.7] (CM) .. controls (-4,-2) and (-7,-3.5) .. (BR);
% Orientation
\draw[-{Stealth}, popdark, thick] (13.5,-6) -- (13.5,-4.5) node[above]{\footnotesize\textbf{N}};
\end{tikzpicture}
\legende{Croquis schématique « Le Cameroun et ses partenaires bilatéraux majeurs ». Les continents sont figurés par des blocs grisés très simplifiés (Afrique au centre, Amérique du Nord et du Sud à l'ouest, Europe au nord, Asie à l'est). L'épaisseur de chaque flèche est proportionnelle au volume des échanges commerciaux (échelle : 1 mm = 100 Mrd FCFA ; Chine 21 mm, France 12,5 mm, USA 4,5 mm, Royaume-Uni 3,2 mm, Japon 1,8 mm, Brésil 1 mm). La couleur indique le domaine dominant de coopération : bleu = infrastructures/énergie, orange = éducation/formation, violet = santé/gouvernance, vert = agriculture. Échelle et distances purement indicatives.}
\end{popfigure}

\medskip
\noindent \textbf{Étape 2 : Tableau comparatif synthétique (corrigé type)}
\begin{center}
\begin{tabularx}{\linewidth}{|l|X|X|X|X|}\hline
\rowcolor{popblueL}
\textbf{Partenaire} & \textbf{Domaines majeurs} & \textbf{3 Réalisations emblématiques} & \textbf{Avantages} & \textbf{Limites / Vigilance} \\ \hline
France & Éducation, santé, gouvernance, défense, francophonie & 1. Réhabilitation de lycées bilingues (FSP) 2. Appui au secteur de la santé (C2D) 3. Formation des officiers (EMIA) & Proximité culturelle et juridique, langue, monnaie commune, réseau d'anciens élèves dense, expertise pédagogique & Conditionnalité de gouvernance, perception « néocoloniale », désengagement budgétaire relatif \\ \hline
Grande-Bretagne & Justice, anglais des affaires, sécurité maritime, climat & 1. Bourses Chevening (cadres) 2. Appui à l'agro-industrie 3. Programme de sécurité dans le golfe de Guinée & Expertise juridique (Common Law), langue anglaise, soft power éducatif (British Council) & Volume financier faible, focus sur la zone anglophone, visibilité réduite depuis le Brexit \\ \hline
Chine & Infrastructures, mines, énergie, agriculture, formation professionnelle & 1. Port de Kribi (phases 1 et 2) 2. Barrages de Memve'ele et Lom Pangar 3. Lycées techniques (Bertoua, Ebolowa) & Financements massifs et rapides, non-conditionnalité politique, transfert technique (BTP), alignement SND30 & Risque d'endettement, main-d'œuvre chinoise importée, maintenance variable, asymétrie commerciale \\ \hline
Japon & Agriculture (riz), infrastructures, santé, TIC, ressources humaines & 1. Riziculture SEMRY (Extrême-Nord) 2. Réhabilitation du pont sur le Wouri 3. Bourses MEXT / volontaires JICA & Excellence technique, fiabilité, dons et prêts très concessionnels, formation pointue & Volume limité, procédures lentes, barrière linguistique, projets ciblés (pas d'appui budgétaire général) \\ \hline
USA & Santé (PEPFAR), sécurité (BIR), entrepreneuriat (YALI/AGOA), énergie & 1. PEPFAR (lutte contre le VIH/SIDA) 2. AGOA (exportations textile/agriculture) 3. Appui au BIR (Extrême-Nord) & Puissance financière et technologique, innovation, appui sécuritaire en zone fragile, diaspora dynamique & Conditionnalité forte (droits humains, gouvernance), engagement variable selon les administrations, distance culturelle \\ \hline
Brésil & Agriculture tropicale (Embrapa), bioénergie, défense, éducation technique & 1. Appui coton/soja (Nord, Adamaoua) 2. Bourses PEC-G (étudiants) 3. Coopération militaire (avions Super Tucano) & Similitudes agro-climatiques (savane), expertise tropicale unique, coopération Sud-Sud égalitaire & Distance géographique, faible volume d'échanges, barrière linguistique (portugais), instabilité politique interne \\ \hline
\end{tabularx}
\end{center}

\medskip
\noindent \textbf{Étape 3 : Note d'argumentation (corrigé type, environ 15 lignes)}
\begin{quote}
\textbf{Objet : choix du partenaire prioritaire pour le financement du Lycée Technique Industriel de Maroua.}

\noindent \textbf{Choix :} la \textbf{République populaire de Chine}.

\noindent \textbf{Arguments :}
\begin{enumerate}
    \item \textbf{Capacité financière et rapidité d'exécution inégalées :} avec un stock d'IDE d'environ 3 500 Mrd FCFA et des échanges de 2 100 Mrd FCFA (premier partenaire commercial), la Chine dispose d'instruments (Exim Bank, FOCAC) capables de mobiliser rapidement les 15 à 20 Mrd FCFA nécessaires (bâtiments, ateliers équipés, centrale solaire), comme l'ont montré le port de Kribi et les lycées techniques de Bertoua et Ebolowa livrés clés en main.
    \item \textbf{Adéquation technique au projet :} la coopération sino-camerounaise inclut explicitement la \textbf{formation professionnelle} (construction d'écoles de métiers, dotation en machines-outils, ateliers de formation). L'expertise chinoise en génie civil, maintenance industrielle et énergies renouvelables (hydroélectricité, solaire) correspond exactement aux filières ciblées.
    \item \textbf{Impact territorial et stratégique (Extrême-Nord) :} la Chine finance déjà des routes et des barrages dans le Grand Nord. Implanter ce lycée à Maroua ancre la coopération dans une zone fragile (insécurité, pauvreté), répond à l'axe d'inclusion sociale de la SND30 et forme la main-d'œuvre locale pour les chantiers environnants.
\end{enumerate}
\noindent \textbf{Mesure d'accompagnement :} négocier une \textbf{clause de transfert de compétences obligatoire} (formation de formateurs camerounais, part significative de main-d'œuvre locale sur le chantier) et un \textbf{contrat de maintenance pluriannuel} financé sur le prêt, afin d'éviter l'écueil de l'équipement non entretenu observé sur certains ouvrages.

\noindent \textbf{Conclusion :} la Chine offre la meilleure combinaison \textbf{volume / rapidité / adéquation technique / ancrage territorial} pour ce projet structurant. Il est recommandé d'inscrire ce dossier à l'ordre du jour de la prochaine session de la Commission mixte de coopération Cameroun-Chine.
\end{quote}

\medskip
\noindent \textbf{Commentaire pédagogique sur le corrigé :}
\begin{itemize}
    \item Le croquis respecte la proportionnalité (Chine $\approx 1{,}7 \times$ France) et la double lecture (géographique par la position, sectorielle par la couleur).
    \item Le tableau synthétise sans recopier le dossier ; il hiérarchise l'information (3 réalisations maximum).
    \item L'argumentation suit la structure imposée, utilise des \textbf{chiffres du dossier} (IDE, échanges, réalisations citées), fait le lien explicite avec le \textbf{contexte local (Maroua, Extrême-Nord, SND30)} et anticipe une \textbf{limite} (maintenance, transfert) en proposant une \textbf{mesure corrective concrète}. C'est une rédaction de niveau Probatoire / Baccalauréat technique : argumentée, chiffrée, contextualisée.
\end{itemize}
\end{exoresolu}

\begin{attention}[Pièges fréquents à signaler aux élèves]
\begin{itemize}
    \item Confondre \cle{premier partenaire commercial} (la Chine, pour le volume des échanges) et \cle{premier investisseur en stock} (la France, pour les IDE accumulés) : les deux classements coexistent et doivent être distingués.
    \item Tracer des flèches simples alors que la coopération est réciproque : la consigne exige des \textbf{flèches doubles} (Cameroun $\leftrightarrow$ partenaire).
    \item Utiliser la couleur pour le quantitatif : règle de sémiologie — « \textit{l'épaisseur pour le quantitatif, la couleur pour le qualitatif} ».
    \item Rédiger une note qui liste le cours sans répondre à la situation : le MINREX attend un choix \textbf{justifié et situé} (Maroua, Extrême-Nord, SND30), pas un catalogue.
    \item Présenter les chiffres du dossier comme des données officielles exactes : ce sont des \textbf{ordres de grandeur indicatifs} à manier avec prudence.
\end{itemize}
\end{attention}

\courssub{Bilan de l'activité}
\noindent Cette activité a permis de mettre en évidence :
\begin{itemize}
    \item \textbf{La maîtrise inégale de la sémiologie graphique :} la construction de flèches proportionnelles (échelle, double sens) et le choix d'un code couleur pertinent (non redondant avec l'épaisseur) sont des compétences exigeantes. La correction collective a permis de réviser la règle : « \textit{l'épaisseur pour le quantitatif, la couleur pour le qualitatif} ».
    \item \textbf{La nécessaire contextualisation de l'argumentation :} les premières productions citaient des chiffres mais oubliaient le \textbf{territoire} (Maroua, Extrême-Nord, insécurité, SND30). Le corrigé type montre qu'un décideur attend une réponse \textbf{située}, pas une liste de cours.
    \item \textbf{La complexité du choix « rationnel » :} aucun partenaire n'est parfait. L'exercice force à \textbf{hiérarchiser les critères} (ici : vitesse, volume, adéquation technique) et à assumer un \textbf{arbitrage} en proposant des garde-fous (clause de transfert). C'est l'essence de la compétence « décider en citoyen responsable ».
    \item \textbf{L'interdisciplinarité effective :} géographie (croquis, flux), économie (IDE, échanges commerciaux), droit (accords-cadres), histoire (héritage colonial), langues (francophonie, anglais, portugais, chinois), enseignement technique (filières industrielles).
\end{itemize}
\noindent \textbf{Prolongement possible :} simulation d'une audition parlementaire (Commission des Relations Extérieures) où chaque groupe défend son choix face aux questions des « députés » (autres groupes), évalué sur la grille : \textit{argumentation factuelle (5 pts) + réactivité et clarté orale (3 pts) + respect du temps (2 pts)}.

\newpage

\coursec{Méthodes et exercices résolus}

\coursec{Le Cameroun dans la coopération bilatérale}

\courssub{Cadre conceptuel : définition et cadre juridique}

\begin{definition}[Coopération bilatérale]
La \cle{coopération bilatérale} désigne l'ensemble des relations d'échanges, d'assistance et de partenariat établies \textbf{directement entre deux États souverains}, sur la base d'accords juridiques (conventions, traités, protocoles d'accord). Elle se distingue de la coopération multilatérale qui implique plusieurs États ou des organisations internationales (ONU, UA, BAD, UE).
\end{definition}

\begin{aretenir}[Principes directeurs (Charte de l'ONU, Constitution du Cameroun)]
\begin{itemize}
\item \textbf{Égalité souveraine} : les deux partenaires négocient d'égal à égal.
\item \textbf{Non-ingérence} : respect des affaires intérieures de l'autre partie.
\item \textbf{Intérêt mutuel} : recherche d'avantages partagés (développement, sécurité, influence).
\item \textbf{Respect des engagements} : \emph{pacta sunt servanda} (les accords doivent être honorés).
\end{itemize}
\end{aretenir}

\begin{savaistu}[Le positionnement diplomatique du Cameroun]
Depuis l'indépendance (1960), le Cameroun pratique une politique de \textbf{non-alignement} et de \textbf{diversification des partenaires}. Cette stratégie vise à éviter la dépendance exclusive vis-à-vis de l'ancienne puissance administrante (France) et à capter des ressources (financements, technologies, marchés) auprès de puissances émergentes (Chine, Brésil) et traditionnelles (Japon, USA, Grande-Bretagne). La \cle{vision 2035} (Cameroun émergent) sert de boussole pour négocier les accords.
\end{savaistu}

\courssub{Les six partenaires bilatéraux prioritaires : fiche de synthèse}

\begin{exemplebox}[Récapitulatif par pays partenaire — à maîtriser pour le Bac]
\begin{center}
\begin{tabularx}{\linewidth}{|l|X|X|}
\hline
\rowcolor{popblueL} \textbf{Pays} & \textbf{Domaines privilégiés \& Acteurs clés} & \textbf{Réalisations / Exemples emblématiques} \\
\hline
\textbf{France} & \cle{Coopération historique} (mandat SDN/ONU 1919--1960). 
Domaines : Éducation/Formation (lycées techniques, ENS), Défense (accords 1960, forces françaises), Monnaie (zone FCFA, garantie Trésor français), Culture (Institut français), Entreprises (Bolloré Logistics, TotalEnergies, Orange Cameroun, CFAO). & Annulation de dette (C2D), Construction/réhabilitation lycées techniques (Douala, Garoua, Bertoua), Appui budgétaire, Formation militaire (EMIA), Francophonie. \\
\hline
\textbf{Grande-Bretagne} & \cle{Histoire partagée} (Southern/Northern Cameroons, mandat/tutelle 1919--1961). 
Domaines : Commonwealth (adhésion 1995), Langue anglaise, Éducation bilingue (British Council), Gouvernance/Droits de l'homme, Commerce. & Bourses Chevening, Appui à la décentralisation, Programme Justice/Sécurité, Investissements (CDC, Pamol - histoire), Commonwealth Games bid. \\
\hline
\textbf{Chine} & \cle{Partenaire stratégique} (relations 1971). 
Domaines : Infrastructures lourdes (BTP), Télécoms/Numérique (Huawei, ZTE), Santé (équipes médicales), Agriculture (riziculture), Formation (bourses universitaires). Financement : Prêts concessionnels (Exim Bank), Dons. & Port en eau profonde de \cle{Kribi} (China Harbour), Barrage de Memve'ele, Routes (Yaoundé-Nsimalen, Kribi-Lolodorf), Palais des Congrès (Yaoundé), Palais des Sports (Yaoundé), Stade d'Olembé, Siège Assemblée Nationale, Réseau 4G/Backbone (Huawei). \\
\hline
\textbf{Japon} & \cle{Coopération "qualitative"} (JICA). 
Domaines : Dons (équipements), Infrastructures ciblées (ponts, routes rurales), Agriculture (riz, machinisme), Santé (hôpitaux, lutte paludisme), Ressources humaines (bourses JICA, jeunes volontaires JOCV). & \textbf{Second pont sur le Wouri} (Douala, prêt JICA), Pont sur la Logone (Yagoua), Centre de formation agricole (Nkolbisson), Équipements hôpitaux (Laquintinie, Gynéco-Obstétrique Yaoundé), Projet rizicole Ndop/Santchou. \\
\hline
\textbf{USA} & \cle{Partenaire stratégique} (AGOA, Sécurité). 
Domaines : Commerce (\cle{AGOA} -- African Growth and Opportunity Act), Santé mondiale (\cle{PEPFAR} VIH/sida, Paludisme PMI), Sécurité/Défense (formation BIR, exercices), Éducation (Peace Corps, Fulbright), Gouvernance. & Accès duty-free au marché US (bois, textile, cacao), Programmes PEPFAR (ARV gratuits), Formation officiers (USAFCENT), Volontaires Peace Corps (éducation rurale), Appui processus électoral. \\
\hline
\textbf{Brésil} & \cle{Coopération Sud-Sud} (affinités tropicales). 
Domaines : Agriculture tropicale (Embrapa -- mecanisation, variétés), Élevage, Formation technique/professionnelle (SENAI), Sports (football), Défense (avions Embraer, blindés). & Projet \cle{Agropoles} (transfert tech. brésilien), Centre de formation professionnelle (Douala), Fourniture avions Super Tucano / blindés Guarani, Coopération universitaire (CAPES/CNPq). \\
\hline
\end{tabularx}
\end{center}
\legende{Tableau de synthèse des six partenaires bilatéraux prioritaires du programme Terminale Technique.}
\end{exemplebox}

% --- MÉTHODES ET EXERCICES RÉSOLUS (Revus et corrigés) ---

\courssub{Méthode 1 : Appliquer les notions de coopération bilatérale à une situation concrète}

\begin{methode}[Analyser une situation de coopération bilatérale]
Pour traiter une situation concrète (chantier, don, accord signé, formation en entreprise) :
\begin{enumerate}
\item \textbf{Identifier les deux partenaires} : le Cameroun et l'État étranger concerné (France, Grande-Bretagne, Chine, Japon, USA, Brésil).
\item \textbf{Qualifier la relation} : s'agit-il d'une coopération \emph{économique}, \emph{financière}, \emph{technique}, \emph{culturelle}, \emph{militaire} ou \emph{sanitaire} ?
\item \textbf{Repérer les acteurs} : États, ambassades, agences d'exécution (JICA, AFD, USAID, CICA), entreprises (ex. Bolloré, China Harbour, Huawei), ONG, populations bénéficiaires.
\item \textbf{Dégager les avantages pour le Cameroun} : infrastructures, emplois, transfert de technologies, formation, devises, désenclavement.
\item \textbf{Repérer les limites éventuelles} : endettement, dépendance technologique, concurrence des produits locaux, conditions politiques des accords, main-d'œuvre importée.
\item \textbf{Conclure par un jugement équilibré} : la coopération bilatérale est profitable si elle respecte la \cle{souveraineté} et sert le \cle{développement national} (Vision 2035).
\end{enumerate}
\textbf{Réflexe} : toujours nommer précisément le pays partenaire et le domaine de coopération ; une réponse vague (« les Blancs nous aident », « les Chinois construisent ») ne vaut aucun point.
\end{methode}

\begin{exoresolu}[Le port en eau profonde de Kribi]
Énoncé : En 2018, le port en eau profonde de Kribi (phase 1), construit avec le concours de l'entreprise \textbf{China Harbour Engineering Company (CHEC)} et un financement de la \textbf{Exim Bank of China}, est entré en service. Un élève déclare : « Ce projet n'a rien à voir avec la coopération bilatérale, c'est juste une affaire d'entreprises. »
\begin{enumerate}
\item Montre que ce projet relève bien de la coopération bilatérale Cameroun--Chine.
\item Dégage deux avantages de ce projet pour le Cameroun et une limite possible.
\end{enumerate}
\tcblower
\textbf{Solution détaillée.}
\begin{enumerate}
\item \textbf{Identification des partenaires} : les deux États sont le Cameroun et la République populaire de Chine. Même si une entreprise publique chinoise (CHEC) réalise les travaux, le projet repose sur un \textbf{accord intergouvernemental} et sur un \textbf{prêt concessionnel de l'État chinois} (Exim Bank). Or la coopération bilatérale désigne précisément l'ensemble des relations d'aide et d'échanges entre deux États. L'affirmation de l'élève est donc fausse : la présence d'entreprises privées ou publiques n'efface pas le caractère étatique de l'accord-cadre.
\item \textbf{Avantages} : (a) Le port désengorge le port de Douala et permet d'accueillir des navires de grande capacité (jusqu'à 100 000 DWT), ce qui dynamise le commerce extérieur et la compétitivité logistique ; (b) Le chantier a créé des emplois directs/indirects et favorise l'industrialisation de la région du Sud (zone industrielle de Kribi, usine d'aluminium, gazoduc GNL). \textbf{Limite} : Le financement par prêt (même concessionnel) accroît la \cle{dette extérieure} du Cameroun envers la Chine, ce qui peut créer une \cle{dépendance financière} si les recettes portuaires tardent à couvrir le service de la dette.
\end{enumerate}
\textbf{Conclusion} : Le port de Kribi illustre une coopération bilatérale sino-camerounaise de type \emph{économique et financier}, source d'infrastructures structurantes mais qui impose une gestion rigoureuse de l'endettement et de la maîtrise d'ouvrage.
\end{exoresolu}

\courssub{Méthode 2 : Rédiger et justifier une réponse argumentée (Dissertation courte)}

\begin{methode}[Construire une argumentation en ECM -- Standard Bac]
Pour rédiger une réponse justifiée (question de synthèse, dissertation 15--20 lignes) :
\begin{enumerate}
\item \textbf{Poser une thèse claire} dès l'introduction (une phrase qui répond directement à la question).
\item \textbf{Définir les mots-clés} du sujet (coopération bilatérale, partenariat, aide au développement, souveraineté).
\item \textbf{Développer 2 ou 3 arguments}, chacun illustré par un \textbf{exemple précis, daté et localisé} tiré du cours (ex. France : C2D + Lycées techniques ; Chine : Kribi + Backbone numérique ; Japon : Pont Wouri + Riziculture ; USA : AGOA + PEPFAR).
\item \textbf{Envisager la thèse adverse (nuance)} : limites (dette, conditionnalité, dépendance), cela montre l'esprit critique attendu.
\item \textbf{Conclure} en répondant à nouveau à la question posée, sans introduire d'idée nouvelle, en ouvrant sur la Vision 2035.
\item \textbf{Soigner la langue} : phrases courtes, vocabulaire du cours (\cle{transfert de technologie}, \cle{conditionnalité}, \cle{désenclavement}), aucune abréviation de type SMS.
\end{enumerate}
\textbf{Formule à retenir} : \cle{Thèse} + \cle{Définitions} + \cle{Arguments exemplifiés} + \cle{Nuance} + \cle{Conclusion}.
\end{methode}

\begin{exoresolu}[Sujet de synthèse type Probatoire/Bac]
Énoncé : « La coopération bilatérale est-elle toujours profitable au Cameroun ? » Rédige une réponse argumentée d'une quinzaine de lignes.
\tcblower
\textbf{Solution détaillée (corrigé rédigé).}

La coopération bilatérale désigne les relations d'échanges et d'assistance que le Cameroun entretient avec un autre État souverain (France, Chine, Japon, USA, etc.). On peut affirmer qu'elle est \textbf{globalement profitable, mais sous conditions de négociation équitable et de bonne gouvernance}.

D'abord, elle apporte des \textbf{financements et des infrastructures structurantes} : la coopération avec la \cle{Chine} a permis la réalisation du port en eau profonde de Kribi, du barrage de Memve'ele, du backbone numérique national et de nombreux édifices publics (Palais des Congrès, Assemblée Nationale). Ensuite, elle favorise les \textbf{transferts de compétences et la formation} : la coopération \cle{franco-camerounaise} soutient massivement l'enseignement technique (réhabilitation lycées, formation enseignants) et la défense ; le \cle{Japon} (via JICA) offre des bourses, des équipements médicaux/agricoles et a financé le second pont sur le Wouri. Enfin, elle ouvre des \textbf{marchés extérieurs} : les \cle{USA}, à travers l'\cle{AGOA}, permettent au Cameroun d'exporter bois, cacao, coton, textile sans droits de douane ; le \cle{Brésil} partage son expertise agricole tropicale (Embrapa, Agropoles).

Cependant, cette coopération comporte des \textbf{limites structurelles} : l'endettement croissant envers la Chine (risque de "piège à dette"), la concurrence des produits importés qui étouffe parfois la production locale (riz, poisson, textile), la \cle{conditionnalité} politique ou économique de certains partenaires occidentaux (gouvernance, droits de l'homme, ouverture marchés), et la persistance de schémas néocoloniaux avec l'ancienne puissance administrante (monnaie, défense, entreprises).

En conclusion, la coopération bilatérale est profitable au Cameroun lorsqu'elle est \textbf{négociée dans le respect de la souveraineté nationale}, orientée vers les priorités du \cle{PND 2030} et de la \cle{Vision 2035} (industrialisation, capital humain), et diversifiée pour éviter toute dépendance unique.
\end{exoresolu}

\courssub{Méthode 3 : Identifier le domaine et les réalisations d'un partenaire bilatéral (QCM / Questions courtes)}

\begin{methode}[Associer un pays à ses domaines de coopération -- Réflexes de rapidité]
Pour répondre vite et juste aux questions objectives (QCM, Vrai/Faux, Compléter) :
\begin{enumerate}
\item \textbf{France} : Ancienne puissance administrante (4/5 du territoire) ; Monnaie (FCFA), Défense, Éducation technique, Francophonie, Grandes entreprises (Bolloré, TotalEnergies, Orange, CFAO). \textit{Mot-clé : "Partenaire historique et multiforme"}.
\item \textbf{Grande-Bretagne} : Histoire (Southern/Northern Cameroons) ; Commonwealth, Langue anglaise, British Council, Gouvernance/Droits humains, Éducation bilingue. \textit{Mot-clé : "Lien linguistique et Commonwealth"}.
\item \textbf{Chine} : Infrastructures lourdes (BTP), Télécoms (Huawei/ZTE), Prêts (Exim Bank), Santé (équipes médicales), Agriculture. \textit{Mot-clé : "Infrastructures contre ressources / Prêts"}.
\item \textbf{Japon} : \cle{Dons} (non remboursables), \cle{Bourses} (JICA), Équipements agricoles/sanitaires, Infrastructures ciblées (ponts, routes rurales), Volontaires (JOCV). \textit{Mot-clé : "Coopération qualitative, don, capital humain"}.
\item \textbf{USA} : \cle{AGOA} (accès marché), \cle{PEPFAR/PMI} (Santé), \cle{Peace Corps} (Éducation), Sécurité (BIR), Gouvernance. \textit{Mot-clé : "Commerce préférentiel + Sécurité/Santé"}.
\item \textbf{Brésil} : \cle{Coopération Sud-Sud}, Agriculture tropicale (Embrapa, mécanisation), Formation pro (SENAI), Sports, Défense (Embraer). \textit{Mot-clé : "Affinités tropicales, agriculture, Sud-Sud"}.
\end{enumerate}
\textbf{Réflexe} : Apprendre \textbf{un exemple concret, daté et localisé par pays} ; c'est l'exemple précis qui fait la différence entre la moyenne et l'excellence.
\end{methode}

\begin{exoresolu}[QCM justifié -- Entraînement]
Énoncé : Pour chaque affirmation, dis si elle est vraie ou fausse et justifie en une phrase précise.
\begin{enumerate}
\item La Grande-Bretagne a administré la totalité du territoire camerounais avant l'indépendance.
\item Le Japon intervient au Cameroun notamment par des dons d'équipements et des bourses de formation via la JICA.
\item L'AGOA est un accord commercial bilatéral signé entre le Cameroun et le Brésil.
\item La zone FCFA garantit la convertibilité du franc CFA par le Trésor français, dans le cadre de la coopération monétaire avec la France.
\item Le port de Kribi a été financé intégralement par des dons de la Chine.
\end{enumerate}
\tcblower
\textbf{Solution détaillée.}
\begin{enumerate}
\item \textbf{Faux.} Après la Première Guerre mondiale, le Cameroun allemand est partagé : la \textbf{France} administre la plus grande partie (environ 4/5, Cameroun oriental) sous mandat SDN puis tutelle ONU ; la \textbf{Grande-Bretagne} administre seulement deux bandes occidentales (Southern Cameroons et Northern Cameroons), soit environ 1/5 du territoire.
\item \textbf{Vrai.} La coopération nippo-camerounaise est pilotée par la \cle{JICA} (Japan International Cooperation Agency) : dons d'équipements (agricoles, médicaux, routiers), bourses de master/doctorat pour cadres camerounais, appui à des infrastructures ciblées (ex. second pont sur le Wouri, pont de Yagoua, centres de formation agricole).
\item \textbf{Faux.} L'\cle{AGOA} (\emph{African Growth and Opportunity Act}) est une \textbf{loi unilatérale des États-Unis} (votée en 2000, renouvelée) qui offre un accès préférentiel (droits de douane nuls) au marché américain pour des produits éligibles de pays africains partenaires, dont le Cameroun. Le Brésil n'est pas partie à l'AGOA ; sa coopération porte sur l'agriculture et la formation technique (Sud-Sud).
\item \textbf{Vrai.} La coopération monétaire est un pilier historique des relations franco-camerounaises : la France garantit la convertibilité illimitée du FCFA en euros via un compte d'opérations au Trésor français, contre dépôt de réserves de change et respect de critères de convergence.
\item \textbf{Faux.} Le port de Kribi (phase 1) a été financé à \textbf{85\% par un prêt concessionnel de la Exim Bank of China} (remboursable sur 20 ans avec différé) et à 15\% par le budget national. Ce n'est pas un don, c'un prêt souverain, ce qui alourdit la dette extérieure.
\end{enumerate}
\textbf{Conclusion} : La maîtrise des exemples précis (noms d'agences, nature du financement, localisation) par pays permet de valider ce type de question sans hésitation.
\end{exoresolu}

\courssub{Méthode 4 : Traiter les données chiffrées d'un tableau de coopération (Calculs \& Interprétation)}

\begin{methode}[Exploiter un tableau statistique en ECM -- Compétence transversale]
Devant un tableau (échanges commerciaux, montants d'aide, nombre de bourses, flux d'IDE) :
\begin{enumerate}
\item \textbf{Lire le titre, les unités (millions USD, \%, nombre) et la source} (INS, BEAC, FMI, Douanes, Ministère de l'Économie) avant tout calcul.
\item \textbf{Repérer la valeur maximale et minimale} et nommer les pays concernés.
\item \textbf{Calculer un taux de variation} (évolution dans le temps) : $t = \dfrac{V_f - V_i}{V_i} \times 100$ (résultat en \%). $V_i$ = valeur initiale, $V_f$ = valeur finale.
\item \textbf{Calculer une part (structure)} : $p = \dfrac{\text{partie}}{\text{total}} \times 100$ (résultat en \%).
\item \textbf{Interpréter} : relier le chiffre à une notion du cours (dépendance, diversification des partenaires, poids de la Chine, érosion part européenne).
\item \textbf{Rédiger la conclusion} en une ou deux phrases, avec les unités et le sens de la variation.
\end{enumerate}
\textbf{Piège majeur} : Ne jamais confondre une \textbf{hausse en valeur absolue} (points de pourcentage) et une \textbf{hausse en pourcentage} (taux de variation). Ex : passer de 20\% à 25\% est une hausse de \textbf{5 points} mais de \textbf{25\%} en taux de variation.
\end{methode}

\begin{exoresolu}[Les exportations camerounaises par partenaire (Données simplifiées)]
Énoncé : Le tableau ci-dessous donne la part de quelques partenaires dans les exportations du Cameroun (année N, en \%).
\begin{center}
\begin{tabularx}{\linewidth}{|l|X|X|X|X|}
\hline
\rowcolor{popblueL} Partenaire & Chine & France & Pays-Bas & Brésil \\
\hline
Part (en \%) & 25 & 10 & 12 & 3 \\
\hline
\end{tabularx}
\end{center}
\begin{enumerate}
\item Calcule la part cumulée de la Chine et de la France dans les exportations camerounaises.
\item La part de la Chine passe de 20 \% (année N-5) à 25 \% (année N). Calcule le taux de variation de cette part.
\item Tire une conclusion sur l'évolution de la coopération bilatérale du Cameroun à la lumière de ces chiffres.
\end{enumerate}
\tcblower
\textbf{Solution détaillée.}
\begin{enumerate}
\item Part cumulée Chine + France : $25 + 10 = 35$~\%. Les deux pays absorbent donc \textbf{35 \%} des exportations camerounaises, soit plus du tiers du total affiché.
\item Taux de variation de la part chinoise :
\[ t = \dfrac{V_f - V_i}{V_i} \times 100 = \dfrac{25 - 20}{20} \times 100 = \dfrac{5}{20} \times 100 = 25~\%. \]
La part de la Chine a donc augmenté de \textbf{25 \%} en cinq ans (attention : il s'agit d'une hausse relative de 25 \%, soit 5 points de pourcentage, et non d'une hausse de 5 \%).
\item \textbf{Interprétation} : La Chine est devenue le premier client du Cameroun (25 \%), devançant nettement la France (10 \%), partenaire historique, et les Pays-Bas (12 \%). Cela traduit la \textbf{diversification réussie des partenaires commerciaux} du Cameroun et le poids croissant de la coopération sino-camerounaise (matières premières : bois, coton, minerais), tout en montrant que les partenaires traditionnels (France, UE via Pays-Bas) restent des débouchés significatifs. La part du Brésil (3 \%) reste marginale sur l'export mais croissante sur l'import (engrais, machines).
\end{enumerate}
\textbf{Conclusion} : Les chiffres confirment le \textbf{rééquilibrage géographique} de la coopération bilatérale camerounaise vers l'Asie (Chine), sans rupture brutale avec les partenaires occidentaux, conformément à la stratégie de diversification.
\end{exoresolu}

\courssub{Méthode 5 : Étudier un document (Extrait de discours, Article de presse, Texte d'accord)}

\begin{methode}[Analyser un document en ECM -- Démarche en 5 étapes]
\begin{enumerate}
\item \textbf{Présenter le document} : Nature (discours officiel, communiqué conjoint, article de presse, extrait de traité), Auteur/Émetteur, Date, Contexte (visite d'État, sommet, signature accord).
\item \textbf{Dégager l'idée principale} (thèse du document) en une phrase claire.
\item \textbf{Relever dans le texte} deux ou trois \textbf{éléments précis} (chiffres, noms de projets, engagements, mots forts) qui prouvent l'existence et la nature de la coopération bilatérale.
\item \textbf{Mettre le document en relation avec le cours} : Quel partenaire ? Quel domaine (économique, militaire, éducatif) ? Quels avantages / limites apparaissent ?
\item \textbf{Porter un jugement critique} fondé, sans parti pris : Le texte est-il un discours officiel (donc valorisant, langue de bois possible) ? Un article de presse (analyse externe) ? Un accord juridique (texte contraignant) ? Manque-t-il des informations (coût, contreparties, clause secrète) ?
\end{enumerate}
\textbf{Réflexe} : Citer brièvement le texte entre guillemets pour appuyer chaque réponse ; une affirmation sans preuve tirée du document est pénalisée.
\end{methode}

\begin{exoresolu}[Extrait de discours officiel -- Visite d'État]
Énoncé : Document : « Lors de sa visite à Yaoundé en juillet 2023, le Président de la République française a annoncé l'effacement d'une partie de la dette du Cameroun via le C2D (Contrat de Désendettement et de Développement) et la construction de trois nouveaux lycées techniques à Douala, Ngaoundéré et Bertoua, affirmant que "la France reste le premier partenaire du Cameroun". »
\begin{enumerate}
\item Présente ce document (nature, auteur, date, contexte).
\item Releve deux éléments du texte qui illustrent la coopération franco-camerounaise et précise pour chacun le domaine concerné.
\item Explique, en t'appuyant sur tes connaissances, pourquoi la France est qualifiée de « partenaire historique » du Cameroun.
\end{enumerate}
\tcblower
\textbf{Solution détaillée.}
\begin{enumerate}
\item Il s'agit d'un \textbf{extrait de discours officiel} (communication présidentielle), prononcé par le \textbf{Président de la République française} (Émetteur), en \textbf{juillet 2023} (Date), lors d'une \textbf{visite d'État à Yaoundé} (Contexte : relance relation bilatérale, sommet Afrique-France).
\item Deux éléments illustrant la coopération :
\begin{itemize}
\item « \emph{l'effacement d'une partie de la dette via le C2D} » : relève de la coopération \textbf{financière / budgétaire} (mécanisme dette-développement, conversion dette en investissements sociaux).
\item « \emph{la construction de trois nouveaux lycées techniques} » : relève de la coopération \textbf{éducative et technique} (formation professionnelle, insertion jeunes, appui à la stratégie "École pour tous" / PND 2030).
\end{itemize}
\item La France est le \textbf{partenaire historique} car : (1) Elle a administré la majeure partie du territoire (Cameroun oriental) sous \textbf{mandat SDN (1919) puis tutelle ONU (1945--1960)} ; (2) Le \textbf{français} est langue officielle et de travail (avec l'anglais) ; (3) Des \textbf{accords de coopération} (Défense, Monnaie, Éducation, Culture, Justice) ont été signés dès l'indépendance (1960) et régulièrement renouvelés ; (4) La \textbf{communauté d'affaires française} est la plus implantée (Bolloré, TotalEnergies, Orange, CFAO, Société Générale) ; (5) La France est un contributeur majeur via l'\cle{AFD} (Agence Française de Développement) et le \cle{C2D}.
\end{enumerate}
\textbf{Conclusion} : Ce discours illustre la \textbf{densité, la continuité et la multiformité} de la coopération franco-camerounaise (financière, éducative, linguistique, économique), tout en relevant d'une \textbf{communication officielle nécessairement valorisante}, qu'il faut lire avec \cle{esprit critique} (contreparties, conditionnalités, souveraineté).
\end{exoresolu}

\courssub{Méthode 6 : Comparer deux partenaires bilatéraux et en tirer un bilan synthétique}

\begin{methode}[Rédiger une comparaison équilibrée -- Tableau croisé]
\begin{enumerate}
\item \textbf{Choisir 3 à 4 critères communs} : Domaines d'intervention, Nature de l'appui (Don / Prêt / Investissement / Accès marché), Acteurs principaux, Avantages pour le Cameroun, Limites / Risques.
\item \textbf{Présenter chaque partenaire critère par critère} (et non l'un après l'autre sans lien) : utiliser un tableau comparatif.
\item \textbf{Utiliser le vocabulaire précis} : "Prêt concessionnel", "Don", "Accès préférentiel", "Transfert de technologie", "Conditionnalité", "Contenu local".
\item \textbf{Conclure par un bilan synthétique} : Complémentarité des partenaires, Diversification comme assurance, Risques communs (endettement, dépendance, gouvernance), Condition de réussite (appropriation nationale, transparence, formation des cadres camerounais).
\end{enumerate}
\textbf{Formule à retenir} : Un bon bilan = \cle{Avantages} + \cle{Limites} + \cle{Conditions de réussite} (Souveraineté, Transparence, Capital humain).
\end{methode}

\begin{exoresolu}[Comparer la coopération avec la Chine et avec les États-Unis]
Énoncé : Compare la coopération du Cameroun avec la Chine et celle avec les États-Unis selon trois critères pertinents, puis dégage un enjeu commun majeur.
\tcblower
\textbf{Solution détaillée.}
\begin{center}
\begin{tabularx}{\linewidth}{|l|X|X|}
\hline
\rowcolor{popblueL} \textbf{Critère de comparaison} & \textbf{Chine (RPC)} & \textbf{États-Unis (USA)} \\
\hline
\textbf{1. Domaines privilégiés} & Infrastructures lourdes (ports, barrages, routes, stades, bâtiments administratifs), Télécommunications / Numérique (backbone, 4G/5G), Santé (équipes médicales), Agriculture (riz). & Commerce (via \cle{AGOA}), Santé mondiale (\cle{PEPFAR} VIH/sida, \cle{PMI} Paludisme), Sécurité/Défense (formation BIR, renseignement), Éducation (Peace Corps, Fulbright), Gouvernance/Droits humains. \\
\hline
\textbf{2. Nature de l'appui financier / technique} & \textbf{Prêts concessionnels} (Exim Bank, taux bas, longs délais) + \textbf{Réalisation par entreprises chinoises} (EPC -- Engineering Procurement Construction). Peu de dons purs. & \textbf{Dons} (budgétaires, projets via USAID), \textbf{Accès de marché non réciproque} (AGOA = avantage comparatif), \textbf{Assistance technique} (experts, volontaires), \textbf{Formation militaire}. Pas de prêts directs massifs pour infrastructures. \\
\hline
\textbf{3. Principales limites / Risques pour le Cameroun} & \textbf{Endettement croissant} (risque soutenabilité), \textbf{Faible contenu local} (main-d'œuvre/encadrement chinois importés, matériaux importés), \textbf{Opacité} de certains contrats, \textbf{Dépendance technologique} (Huawei/ZTE). & \textbf{Conditionnalité politique} (droits humains, gouvernance, démocratie), \textbf{Volumes commerciaux limités} (peu de produits camerounais éligibles AGOA compétitifs), \textbf{Aide ciblée} (santé/sécurité) moins visible sur infrastructures lourdes. \\
\hline
\end{tabularx}
\end{center}

\textbf{Enjeu commun majeur} : Dans les deux cas, le Cameroun doit veiller à ce que la coopération serve \textbf{ses priorités nationales} (industrialisation, emploi des jeunes, transformation locale des matières premières, émergence 2035) et à \textbf{préserver sa souveraineté} dans la négociation, l'exécution et le contrôle des accords (transparence, clauses de contenu local, formation de contre-expertise nationale).

\textbf{Conclusion} : La coopération \cle{sino-camerounaise} est tournée vers les \textbf{infrastructures physiques financées par endettement}, tandis que la coopération \cle{améro-camerounaise} privilégie le \textbf{capital humain, la santé, la sécurité et l'accès au marché} ; les deux sont \textbf{complémentaires} (infrastructures + capital humain) mais exigent une \cle{gestion lucide et souveraine} de la dépendance (financière vis-à-vis de la Chine, politique/commerciale vis-à-vis des USA).
\end{exoresolu}

\courssub{Bilan : Avantages, limites et enjeux de la coopération bilatérale}

\begin{aretenir}[Bilan synthétique pour le Probatoire / Bac]
\begin{itemize}
\item \textbf{Avantages majeurs} : Diversification des sources de financement (dons, prêts, IDE), Accès à des technologies de pointe (numérique, agricole, médicale), Infrastructures structurantes (transport, énergie, numérique), Formation du capital humain (bourses, lycées techniques, centres de formation), Ouverture de marchés (AGOA, Chine), Renforcement de la souveraineté par le multipartenariat.
\item \textbf{Limites / Risques} : \cle{Endettement} extérieur (Chine, France), \cle{Dépendance} technologique et commerciale, \cle{Conditionnalité} politique (Occident) ou économique (clauses secrètes, ressources), \cle{Concurrence déloyale} aux PME locales (importations massives, main-d'œuvre étrangère), \cle{Gouvernance} (corruption, opacité contrats, faible évaluation d'impact).
\item \textbf{Enjeux pour le Cameroun (Vision 2035)} : \textbf{Négocier "gagnant-gagnant"} (contenu local, transfert tech., clauses de révision), \textbf{Renforcer l'expertise nationale} (juristes, économistes, ingénieurs pour auditer les accords), \textbf{Intégrer la coopération dans le PND} (pas de projets "cadeaux" hors priorités), \textbf{Promouvoir le "Made in Cameroon"} via les accords (ex. transformation bois/cacao avant export), \textbf{Évaluer ex post} chaque grand projet (audit citoyen, Cour des Comptes).
\end{itemize}
\end{aretenir}

\begin{exercice}[Entraînement Probatoire -- Sujet type]
\textbf{Durée conseillée : 30 min}
\begin{enumerate}
\item \textbf{Définis} la coopération bilatérale et cite les six pays partenaires prioritaires du Cameroun au programme. (3 pts)
\item \textbf{Complète} le tableau ci-dessous : (6 pts)
\begin{center}
\begin{tabularx}{\linewidth}{|l|X|X|}
\hline
\rowcolor{popblueL} \textbf{Pays} & \textbf{Agence / Instrument principal} & \textbf{Une réalisation emblématique (projet + localisation)} \\
\hline
Japon & \ldots & \ldots \\
\hline
USA & \ldots & \ldots \\
\hline
Brésil & \ldots & \ldots \\
\hline
\end{tabularx}
\end{center}
\item \textbf{Calcul} : En 2020, les importations camerounaises depuis la Chine s'élevaient à 800 milliards FCFA. En 2023, elles atteignent 1 200 milliards FCFA. Calcule le taux de variation de ces importations sur la période. Interprète brièvement ce résultat. (4 pts)
\item \textbf{Argumentation} : « La coopération avec la Chine est plus profitable au Cameroun que celle avec la France. » Discute cette affirmation en t'appuyant sur des exemples précis (infrastructures, dette, formation, souveraineté). (7 pts)
\end{enumerate}
\end{exercice}

\begin{corrige}[Exercice n°1 -- Entraînement Probatoire]
\begin{enumerate}
\item \textbf{Définition} : Relation d'échanges et d'assistance entre deux États souverains (2 pts). \textbf{Six pays} : France, Grande-Bretagne, Chine, Japon, USA, Brésil (1 pt).
\item \textbf{Tableau complété} :
\begin{center}
\begin{tabularx}{\linewidth}{|l|X|X|}
\hline
\rowcolor{popblueL} \textbf{Pays} & \textbf{Agence / Instrument principal} & \textbf{Une réalisation emblématique} \\
\hline
Japon & JICA (Japan International Cooperation Agency) & Second pont sur le Wouri (Douala) / Centre formation agricole Nkolbisson / Dons équipements médicaux. \\
\hline
USA & USAID / AGOA (Loi) / PEPFAR / Peace Corps & Programme PEPFAR (ARV gratuits, lutte VIH) / Accès marché US via AGOA (bois, cacao) / Formation officiers BIR. \\
\hline
Brésil & ABC (Agence Brésilienne Coopération) / Embrapa / SENAI & Projet Agropoles (transfert tech. agricole) / Centre formation pro Douala / Fourniture avions Super Tucano. \\
\hline
\end{tabularx}
\end{center}
(2 pts par ligne bien remplie = 6 pts).
\item \textbf{Calcul} : $t = \dfrac{1200 - 800}{800} \times 100 = \dfrac{400}{800} \times 100 = 50~\%$. (2 pts calcul + 1 pt formule/unité).
\textbf{Interprétation} : Hausse de 50\% en 3 ans, confirmant la Chine comme premier fournisseur du Cameroun (machines, électronique, textiles, produits manufacturés), illustrant l'approfondissement de l'interdépendance économique sino-camerounaise mais posant la question de la balance commerciale déficitaire et de la concurrence pour l'industrie locale naissante. (1 pt).
\item \textbf{Argumentation} (7 pts : Thèse 1 + 2 arguments Chine 2 + 2 arguments France 1 + Nuance 1 + Conclusion 1) :
\textit{Thèse} : L'affirmation est réductrice ; les deux coopérations sont de nature différente et complémentaires, la "profitabilité" dépend du secteur et de la gestion.
\textit{Arguments Chine} : + Infrastructures visibles (Kribi, barrages, stades, numérique) sans conditionnalité politique ; + Financements massifs rapides. -- Endettement lourd (risque soutenabilité), faible contenu local/emplois, opacité contrats.
\textit{Arguments France} : + Partenaire historique, confiance institutionnelle, monnaie stable (FCFA), formation technique/universitaire de qualité (lycées, ENS, bourses), C2D (dette convertie en investissements), entreprises créatrices d'emplois locaux (Bolloré, Total, Orange). -- Conditionnalité gouvernance, présence militaire parfois contestée, part de marché en érosion face à la Chine.
\textit{Nuance} : La Chine gagne sur le "hard" (infrastructures), la France sur le "soft" (institutions, formation, monnaie, droit). Le Cameroun a intérêt à les articuler.
\textit{Conclusion} : Ni l'une ni l'autre n'est "plus profitable" en absolu ; le profit dépend de la \cle{négociation souveraine}, de la \cle{transparence} et de l'\cle{alignement sur le PND 2030}.
\end{enumerate}
\end{corrige}

\begin{attention}[Les 5 erreurs qui coûtent des points au Bac / Probatoire]
\begin{enumerate}
\item \textbf{Confondre coopération bilatérale et multilatérale} : La coopération bilatérale lie le Cameroun à \emph{un seul} État (France, Chine, Japon...). L'ONU, la BAD, le FMI, l'Union européenne, la CEMAC, l'UA relèvent du \textbf{multilatéral}. Citer l'ONU ou la BAD dans une question sur la coopération bilatérale est \textbf{hors sujet}.
\item \textbf{Affirmer que la Grande-Bretagne a colonisé/administré la totalité du Cameroun} : Elle n'a administré que les régions occidentales (\cle{Southern Cameroons} et \cle{Northern Cameroons}), soit environ 1/5 du territoire, sous mandat SDN puis tutelle ONU. La \textbf{France} a administré les 4/5 restants (\cle{Cameroun oriental}). Dire "la Grande-Bretagne a colonisé le Cameroun" est \textbf{faux} et sanctionné.
\item \textbf{Attribuer l'AGOA ou la JICA au mauvais pays} : L'\cle{AGOA} est un dispositif des \textbf{États-Unis} (loi américaine), pas du Brésil, ni de la Chine, ni de l'UE. La \cle{JICA} est l'agence de coopération du \textbf{Japon}. La \cle{Exim Bank} finance les projets \textbf{Chine}. L'\cle{AFD} finance les projets \textbf{France}. Confondre ces acronymes = 0 point à la question.
\item \textbf{Se tromper dans les calculs de pourcentage (piège classique)} : Une part qui passe de 20 \% à 25 \% augmente de \textbf{5 points de pourcentage}, mais son \textbf{taux de variation} est de $\dfrac{5}{20}\times 100 = 25~\%$. Toujours diviser par la \textbf{valeur initiale} ($V_i$). Écrire "hausse de 5\%" au lieu de "25\%" ou "5 points" fait perdre les points de calcul et d'interprétation.
\item \textbf{Rédiger une argumentation sans exemple précis, sans nuance ni conclusion} : Une liste d'idées vagues (« la Chine aide beaucoup », « la France nous domine ») \textbf{sans projet nommé} (Kribi, Memve'ele, C2D, Lycées techniques, Pont Wouri, AGOA, PEPFAR), \textbf{sans nuance sur les limites} (dette, conditionnalité, contenu local) et \textbf{sans phrase de conclusion} répondant à la question fait perdre la moitié des points de synthèse (souvent notée sur 7 ou 10 pts).
\end{enumerate}
\end{attention}

\newpage

\coursec{Exercices}

\courssub{Je vérifie mes connaissances}

\begin{exercice}[QCM : Repères historiques et institutionnels]
Parmi les affirmations suivantes, une seule est exacte. La recopier en la complétant.
\begin{enumerate}
\item Le plébiscite du 11 février 1961 a concerné : 
\begin{itemize}
\item[A.] Le Cameroun français uniquement.
\item[B.] Le Cameroun britannique (Nord et Sud).
\item[C.] La République du Cameroun et le Cameroun méridional britannique.
\end{itemize}
\item Le Cameroun est membre du Commonwealth depuis :
\begin{itemize}
\item[A.] 1960 (indépendance).
\item[B.] 1961 (réunification).
\item[C.] 1995.
\end{itemize}
\item L'Agence Française de Développement (AFD) est l'opérateur principal de la coopération :
\begin{itemize}
\item[A.] Technique et financière avec la France.
\item[B.] Militaire avec les USA.
\item[C.] Commerciale avec la Chine.
\end{itemize}
\item L'AGOA (African Growth and Opportunity Act) est un dispositif d'accès au marché :
\begin{itemize}
\item[A.] Européen.
\item[B.] Américain.
\item[C.] Chinois.
\end{itemize}
\end{enumerate}
\end{exercice}

\begin{exercice}[Vrai / Faux justifié]
Dire si les affirmations sont vraies ou fausses. Justifier chaque réponse par une phrase précise.
\begin{enumerate}
\item La coopération bilatérale implique obligatoirement plus de deux États.
\item Le Japon intervient au Cameroun principalement à travers la JICA (Japan International Cooperation Agency) dans les secteurs de la santé, de l'éducation et des infrastructures.
\item La coopération Sud-Sud désigne exclusivement les relations entre le Cameroun et le Brésil.
\item Le mandat de la Société des Nations (SDN) confié à la France et à la Grande-Bretagne après 1919 a posé les bases linguistiques et administratives actuelles du Cameroun.
\end{enumerate}
\end{exercice}

\begin{exercice}[Définitions et concepts clés]
Définir brièvement (2 à 3 lignes) les termes suivants en les situant dans le contexte de la coopération bilatérale du Cameroun :
\begin{enumerate}
\item Coopération bilatérale.
\item Partenaire historique (citer les deux pays concernés pour le Cameroun).
\item Coopération décentralisée.
\item Accord de Cotonou (puis Accord de Samoa) : lien avec l'Union Européenne.
\end{enumerate}
\end{exercice}

\begin{exercice}[Calculs et données chiffrées]
On donne les données suivantes (estimations annuelles récentes) pour le Cameroun :
\begin{itemize}
\item Exportations totales : 3 500 milliards FCFA.
\item Exportations vers la Chine : 700 milliards FCFA.
\item Exportations vers les Pays-Bas : 525 milliards FCFA.
\item Exportations vers la France : 350 milliards FCFA.
\item Exportations vers l'Italie : 280 milliards FCFA.
\end{itemize}
\begin{enumerate}
\item Calculer la part (en \%) de la Chine dans les exportations camerounaises.
\item Calculer la part cumulée des quatre premiers clients cités ci-dessus.
\item En déduire la part des « autres pays ». Que cela suggère-t-il sur la diversification des partenariats ?
\end{enumerate}
\end{exercice}

\courssub{Je m'entraîne}

\begin{exercice}[Analyse de la coopération franco-camerounaise]
À partir de vos connaissances et des documents vus en cours (rapports AFD, accords de défense, zone franc), rédiger une réponse structurée (15 lignes maximum) aux questions suivantes :
\begin{enumerate}
\item Citer trois secteurs prioritaires de l'intervention de l'AFD au Cameroun.
\item Expliquer en quoi le « Franc CFA / Euro » constitue un pilier monétaire et financier de cette coopération.
\item Donner un exemple concret de coopération décentralisée entre une collectivité camerounaise et une collectivité française.
\end{enumerate}
\end{exercice}

\begin{exercice}[La coopération avec le Royaume-Uni et le Commonwealth]
\begin{enumerate}
\item Rappeler la particularité historique du Cameroun méridional britannique (plébiscite, réunification).
\item Citer deux avantages concrets de l'adhésion du Cameroun au Commonwealth (1995) pour la jeunesse et l'enseignement technique.
\item Identifier un domaine de coopération bilatérale actuelle entre Yaoundé et Londres (hors Commonwealth) : éducation, justice, ou sécurité.
\end{enumerate}
\end{exercice}

\begin{exercice}[Focus Asie : Chine et Japon — Modèles contrastés]
Compléter le tableau comparatif suivant en utilisant les mots-clés : \emph{Prêts concessionnels / Dons / Infrastructures lourdes / Transfert de technologie / Santé / Éducation / JICA / Exim Bank / FOCAC / TICAD}.

\begin{center}
\begin{tabularx}{\linewidth}{|l|X|X|}
\hline
\textbf{Critère} & \textbf{Coopération avec la Chine} & \textbf{Coopération avec le Japon} \\
\hline
Principal instrument financier & & \\
\hline
Forum multilatéral de référence & & \\
\hline
Agence d'exécution principale & & \\
\hline
Secteurs emblématiques & & \\
\hline
Type d'apport dominant & & \\
\hline
\end{tabularx}
\end{center}
\end{exercice}

\begin{exercice}[Amériques : USA et Brésil — Opportunités spécifiques]
\begin{enumerate}
\item L'AGOA permet l'entrée en franchise de droits de douane aux États-Unis pour certains produits. Citer trois produits camerounais éligibles (ex: bois, textile, produits agricoles transformés).
\item Le programme \emph{Power Africa} vise à améliorer l'accès à l'énergie. Donner un exemple de projet soutenu au Cameroun.
\item La coopération avec le Brésil (Sud-Sud) s'illustre notamment dans l'agriculture et la bioénergie. Citer l'institution brésilienne de référence pour la recherche agricole tropicale (Embrapa) et un projet de transfert de technologie (ex: coton, manioc, soja).
\end{enumerate}
\end{exercice}

\courssub{Je me prépare au Bac}

\begin{exercice}[Dissertation guidée : Chance ou dépendance ?]
\textbf{Sujet :} « La coopération bilatérale est-elle une chance ou une dépendance pour le Cameroun ? »

\textbf{Consignes :} Vous traiterez ce sujet sous la forme d'une dissertation structurée (Introduction, Développement en deux parties équilibrées, Conclusion).
\begin{enumerate}
\item \textbf{Introduction :} Accroche (actualité : sommet FOCAC, sommet France-Afrique, adhésion Commonwealth), définition de la coopération bilatérale, énoncé du paradoxe (ressources vs souveraineté), annonce du plan.
\item \textbf{Première partie : La coopération bilatérale, une chance pour le développement (arguments + exemples camerounais précis).} 
\begin{itemize}
\item Argument 1 : Apport en infrastructures structurantes (ex: barrages, stades, routes, ports — Chine, France).
\item Argument 2 : Appui au capital humain et services sociaux (santé, éducation, formation professionnelle — Japon, France, USA).
\item Argument 3 : Diversification des partenaires et pouvoir de négociation accru (leviers Chine / Brésil vs partenaires traditionnels).
\end{itemize}
\item \textbf{Deuxième partie : Les risques de dépendance et les limites structurelles (arguments + exemples).}
\begin{itemize}
\item Argument 1 : Endettement et conditionnalités (ratio dette/PIB, clauses des prêts Exim Bank, ajustements structurels FMI/BM appuyés par partenaires).
\item Argument 2 : Déséquilibre commercial et extraversion de l'économie (exportation matières premières brutes vs importation biens manufacturés, faible transfert de technologie effectif).
\item Argument 3 : Ingérence dans les choix de politique intérieure (droits de l'homme, gouvernance, alignement diplomatique votes ONU).
\end{itemize}
\item \textbf{Conclusion nuancée :} Bilan synthétique, condition de la « chance » (gouvernance, contrat gagnant-gagnant, vision stratégique \emph{Cameroun 2035}), ouverture sur la coopération régionale (CEMAC, UA, ZLECAf).
\end{enumerate}
\end{exercice}

\begin{exercice}[Étude de documents comparée : Chine vs France]
\textbf{Document 1 : Extrait de tableau — Principaux projets chinois au Cameroun (2010-2023)}
\begin{center}
\begin{tabularx}{\linewidth}{|l|X|l|}
\hline
\textbf{Projet} & \textbf{Secteur} & \textbf{Financement estimé} \\
\hline
Barrage de Memve'ele & Énergie / Hydroélectricité & Prêt Exim Bank \textasciitilde 250 Mrd FCFA \\
\hline
Autoroute Yaoundé-Nsimalen & Transport / Infrastructures & Prêt Exim Bank \textasciitilde 300 Mrd FCFA \\
\hline
Stade d'Olembé & Sport / Infrastructures & Don / Prêt \textasciitilde 35 Mrd FCFA \\
\hline
Hôpital de référence de Yaoundé & Santé & Don \textasciitilde 30 Mrd FCFA \\
\hline
Port en eau profonde de Kribi (Phase 1 \& 2) & Transport / Logistique & Prêt Exim Bank \textasciitilde 600 Mrd FCFA \\
\hline
\end{tabularx}
\end{center}

\textbf{Document 2 : Extrait — Coopération française (AFD) au Cameroun (Portefeuille 2020-2023)}
\begin{center}
\begin{tabularx}{\linewidth}{|l|X|l|}
\hline
\textbf{Secteur d'intervention} & \textbf{Projets emblématiques} & \textbf{Type d'instrument} \\
\hline
Eau \& Assainissement & Projet Eau Yaoundé, Assainissement Douala & Dons / Prêts très concessionnels \\
\hline
Éducation / Formation pro & PAREFOP, Lycées techniques, AFD Éducation & Dons (C2D) \\
\hline
Énergie / Climat & Projet Lom Pangar (appui), Solaire Nord & Prêts concessionnels / Dons \\
\hline
Appui budgétaire / Gouvernance & Contrats de désendettement-développement (C2D) & Dons (annulation dette) \\
\hline
Secteur privé / PME & Proparco, lignes de crédit banques locales & Prêts / Garanties / Fonds propres \\
\hline
\end{tabularx}
\end{center}

\textbf{Document 3 : Échanges commerciaux (Exportations Cameroun \textrightarrow Partenaire, 2022 est.)}
\begin{itemize}
\item Vers Chine : \textasciitilde 700 Mrd FCFA (Bois brut, Coton, Pétrole, Minerai).
\item Vers France : \textasciitilde 350 Mrd FCFA (Bois transformé, Cacao, Café, Aluminium).
\item Vers Pays-Bas / Italie / Espagne : Parts significatives (Pétrole, Bois, Cacao).
\end{itemize}

\textbf{Questions :}
\begin{enumerate}
\item \textbf{Caractériser la coopération sino-camerounaise} à partir du Document 1 : nature des projets, instrument financier dominant, secteurs ciblés. (5 lignes)
\item \textbf{Caractériser la coopération franco-camerounaise} à partir du Document 2 : nature des projets, instruments financiers, place du secteur social et de la gouvernance. (5 lignes)
\item \textbf{Comparer les deux modèles} au regard des Documents 1, 2 et 3 : 
\begin{itemize}
\item Structure des échanges (Document 3) : transformation locale vs export brut.
\item Type d'endettement généré.
\item Transfert de technologie et formation locale.
\end{itemize}
\item \textbf{Synthèse (10 lignes) :} En vous appuyant sur l'analyse précédente, discuter l'affirmation : « La Chine construit le hardware (infrastructures), la France finance le software (humain, gouvernance, institutions) ». Nuancer.
\end{enumerate}
\end{exercice}

\begin{exercice}[Sujet de débat argumenté : Choix stratégiques]
\textbf{Sujet :} « Le Cameroun doit-il privilégier ses partenaires historiques (France, Grande-Bretagne) ou ses nouveaux partenaires (Chine, Brésil, Turquie, Russie) pour son émergence à l'horizon 2035 ? »

\textbf{Consignes :} Vous rédigerez une production écrite argumentée (environ 300-400 mots) structurée en trois temps.
\begin{enumerate}
\item \textbf{Prise de position nuancée (Thèse) :} Affirmer clairement qu'il ne s'agit pas d'un choix binaire (ou l'un ou l'autre) mais d'une \emph{diversification stratégique} et d'une \emph{hiérarchisation par secteur}.
\item \textbf{Argumentation illustrée (Développement en 2 ou 3 paragraphes) :}
\begin{itemize}
\item \textbf{Paragraphe 1 : L'apport irremplaçable des partenaires historiques.} Cadre juridique (OHADA, zone Franc), appui budgétaire (C2D), formation d'excellence (Écoles d'ingénieurs, santé), coopération militaire/défense, langue et culture administrative partagée. Exemple : C2D, Lycées techniques, École de Santé Publique.
\item \textbf{Paragraphe 2 : L'apport massif et rapide des nouveaux partenaires.} Infrastructures lourdes « clé en main » (Chine), modèle agricole tropical (Brésil/Embrapa), coûts financiers souvent plus bas, absence de conditionnalités politiques explicites (principe non-ingérence). Exemple : Barrages, Autoroutes, Kribi, Coton/Soja.
\item \textbf{Paragraphe 3 (optionnel) : Les risques à gérer pour chacun.} Dette opaque / clauses secrètes / main-d'œuvre importée (Chine) ; Conditionnalités politiques / lenteurs procédurales / désengagement progressif (France/UE).
\end{itemize}
\item \textbf{Conclusion prospective :} Définir la doctrine camerounaise idéale : « Partenariat gagnant-gagnant », clause de contenu local (main-d'œuvre, transfert tech), alignement sur la \emph{Stratégie Nationale de Développement 2020-2030 (SND30)} et \emph{Cameroun 2035}, rôle de la diplomatie économique pour transformer l'essai (transformation locale, ZLECAf).
\end{enumerate}
\end{exercice}

\newpage

\coursec{Corrigés des exercices}

\coursec{Le Cameroun dans la coopération bilatérale}
\courssub{Exercices corrigés et analyses détaillées}

\begin{corrige}[Exercice 1 — QCM : Repères historiques et institutionnels]
\begin{enumerate}
\item \textbf{Réponse B.} Le plébiscite du 11 février 1961, organisé sous l'égide de l'ONU, a concerné le \cle{Cameroun britannique} (Nord et Sud) : le Nord a choisi de rejoindre le Nigeria, le Sud a choisi la réunification avec la République du Cameroun. Le Cameroun français était déjà indépendant depuis le 1\textsuperscript{er} janvier 1960.
\item \textbf{Réponse C.} Le Cameroun est membre du \cle{Commonwealth} depuis \textbf{1995} (sommet d'Auckland, Nouvelle-Zélande). C'est un cas particulier : un pays majoritairement francophone, sans passé entièrement britannique, admis en raison de son héritage anglophone.
\item \textbf{Réponse A.} L'\cle{AFD} (Agence Française de Développement) est l'opérateur principal de la coopération \textbf{technique et financière} avec la France (dons, prêts concessionnels, C2D).
\item \textbf{Réponse B.} L'\cle{AGOA} (African Growth and Opportunity Act, 2000) est un dispositif d'accès en franchise de droits au marché \textbf{américain} pour les produits des pays africains éligibles.
\end{enumerate}
\textbf{Points de vigilance :} ne pas confondre la date d'indépendance (1960), la réunification (1961) et l'adhésion au Commonwealth (1995) ; ne pas confondre AGOA (marché américain) avec les accords de Cotonou/Samoa (Union européenne).
\end{corrige}

\begin{corrige}[Exercice 2 — Vrai / Faux justifié]
\begin{enumerate}
\item \textbf{FAUX.} La coopération \cle{bilatérale} concerne par définition \textbf{deux} États (ou un État et une institution) ; c'est la coopération \emph{multilatérale} qui implique plus de deux acteurs (ONU, FMI, Union africaine).
\item \textbf{VRAI.} Le Japon intervient principalement à travers la \cle{JICA}, avec des dons et une assistance technique ciblés sur la santé (hôpitaux, vaccination), l'éducation (mathématiques et sciences) et les infrastructures (ponts, eau potable).
\item \textbf{FAUX.} La coopération \cle{Sud-Sud} désigne l'ensemble des relations entre pays en développement : le couple Cameroun-Brésil n'en est qu'un exemple, aux côtés du Cameroun-Chine, Cameroun-Inde ou Cameroun-Turquie.
\item \textbf{VRAI.} Après la défaite allemande de 1916, la SDN a confié en 1919-1922 des \cle{mandats} à la France (Cameroun oriental) et à la Grande-Bretagne (Cameroun occidental) : c'est l'origine du bilinguisme français-anglais et de la dualité des systèmes administratifs et éducatifs actuels.
\end{enumerate}
\textbf{Points de vigilance :} la justification est exigée pour chaque item ; une réponse « vrai/faux » seule ne vaut aucun point au Bac.
\end{corrige}

\begin{corrige}[Exercice 3 — Définitions et concepts clés]
\begin{enumerate}
\item \textbf{Coopération bilatérale :} ensemble des relations d'aide, d'échanges et de partenariats (économiques, financiers, techniques, culturels, militaires) établies entre \textbf{deux États} souverains, formalisées par des accords signés directement entre eux. Exemple : accords Cameroun-France, Cameroun-Chine.
\item \textbf{Partenaire historique :} État lié au Cameroun par un passé colonial ou mandataire et par des relations diplomatiques, économiques et culturelles anciennes et continues. Pour le Cameroun : la \textbf{France} (mandat puis tutelle sur le Cameroun oriental) et la \textbf{Grande-Bretagne} (mandat sur le Cameroun occidental).
\item \textbf{Coopération décentralisée :} partenariat conclu directement entre des \textbf{collectivités territoriales} (communes, régions) de deux pays, sans passer uniquement par les États centraux. Exemple : jumelage entre une commune camerounaise et une ville française (appui à l'eau potable, à l'assainissement, à la formation).
\item \textbf{Accord de Cotonou (2000), puis Accord de Samoa (2023) :} accords de partenariat entre l'\textbf{Union européenne} et les pays ACP (Afrique, Caraïbes, Pacifique), dont le Cameroun. Ils organisent l'aide au développement (FED) et l'accès préférentiel au marché européen ; l'Accord de Samoa succède à Cotonou et recentre le partenariat sur le développement durable, la paix et la gouvernance.
\end{enumerate}
\end{corrige}

\begin{corrige}[Exercice 4 — Calculs et données chiffrées]
\begin{enumerate}
\item \textbf{Part de la Chine :}
\[ \frac{700}{3500} \times 100 = 0{,}20 \times 100 = 20\ \% \]
La Chine absorbe \textbf{20 \%} des exportations camerounaises : c'est le premier client.
\item \textbf{Part cumulée des quatre premiers clients :}
\begin{align*}
\text{Total des 4 clients} &= 700 + 525 + 350 + 280 = 1855 \text{ milliards FCFA}\\
\text{Part cumulée} &= \frac{1855}{3500} \times 100 = 53\ \%
\end{align*}
Les quatre premiers clients (Chine, Pays-Bas, France, Italie) représentent \textbf{53 \%} des exportations.
\item \textbf{Part des autres pays :}
\[ 100\ \% - 53\ \% = 47\ \% \quad\text{soit}\quad 3500 - 1855 = 1645 \text{ milliards FCFA} \]
\textbf{Interprétation :} plus de la moitié des exportations camerounaises dépend de seulement quatre pays. Cela montre une \textbf{concentration forte} des débouchés commerciaux et donc une diversification encore insuffisante : une crise ou une baisse de la demande chez l'un de ces partenaires affecterait fortement les recettes d'exportation du Cameroun. D'où l'enjeu de la diversification des partenariats et de la transformation locale (ZLECAf).
\end{enumerate}
\textbf{Points de vigilance :} vérifier que la somme des parts fait bien 100 \% ; présenter les calculs avec la formule $\frac{\text{partie}}{\text{total}} \times 100$.
\end{corrige}

\begin{corrige}[Exercice 5 — Analyse de la coopération franco-camerounaise]
\begin{enumerate}
\item \textbf{Trois secteurs prioritaires de l'AFD au Cameroun :} (a) l'\textbf{eau et l'assainissement} (projets Eau Yaoundé, assainissement de Douala) ; (b) l'\textbf{éducation et la formation professionnelle} (appui aux lycées techniques, programme PAREFOP) ; (c) l'\textbf{énergie et le climat} (appui au barrage de Lom Pangar, solaire dans le Nord). On peut ajouter l'appui au secteur privé via Proparco.
\item \textbf{Le Franc CFA, pilier monétaire :} le FCFA est \textbf{arrimé à l'euro} (parité fixe : 1 \euro{} = 655,957 FCFA) et sa convertibilité est garantie par le Trésor français. Cela assure la \textbf{stabilité monétaire}, facilite les échanges et les investissements français au Cameroun (membre de la CEMAC, zone franc) et sécurise les prêts et dons de l'AFD. En contrepartie, il lie la politique monétaire camerounaise à celle de la zone euro, ce qui alimente le débat sur la souveraineté monétaire.
\item \textbf{Exemple de coopération décentralisée :} le jumelage entre la \textbf{Communauté urbaine de Douala} et des collectivités françaises (région Hauts-de-France / métropole de Lille) pour l'assainissement et la gestion des déchets ; ou la coopération entre la ville de \textbf{Yaoundé} et des communes françaises pour l'eau potable et la formation des agents municipaux.
\end{enumerate}
\textbf{Points de vigilance :} citer des exemples précis et datés ; distinguer coopération d'État à État (AFD) et coopération décentralisée (collectivités).
\end{corrige}

\begin{corrige}[Exercice 6 — La coopération avec le Royaume-Uni et le Commonwealth]
\begin{enumerate}
\item \textbf{Particularité historique :} le Cameroun méridional britannique était sous \textbf{mandat SDN} puis \textbf{tutelle ONU} britannique depuis 1919-1922. Lors du \textbf{plébiscite du 11 février 1961}, sa population a choisi la réunification avec la République du Cameroun (alors que le Nord a rejoint le Nigeria). Cette réunification a donné naissance, le 1\textsuperscript{er} octobre 1961, à la \textbf{République fédérale du Cameroun}, État bilingue et biculturel.
\item \textbf{Deux avantages du Commonwealth pour la jeunesse et l'enseignement technique :} (a) l'accès aux \textbf{bourses du Commonwealth} (Commonwealth Scholarships) pour étudier dans les universités britanniques et du Commonwealth ; (b) les \textbf{programmes d'échange et de formation professionnelle et technique} (Commonwealth of Learning, formations à distance, renforcement des filières techniques en langue anglaise), ainsi que la participation aux Jeux du Commonwealth, vecteur de rayonnement sportif des jeunes.
\item \textbf{Domaine de coopération bilatérale actuelle :} l'\textbf{éducation} (programme British Council, enseignement de l'anglais, certification des enseignants) ; on peut aussi citer la \textbf{justice} (appui à la lutte contre la corruption et à la formation des magistrats) et la \textbf{sécurité} (coopération dans la lutte contre le terrorisme de Boko Haram dans l'Extrême-Nord).
\end{enumerate}
\end{corrige}

\begin{corrige}[Exercice 7 — Focus Asie : Chine et Japon — Modèles contrastés]
\begin{center}
\begin{tabularx}{\linewidth}{|l|X|X|}
\hline
\rowcolor{popblueL}
\textbf{Critère} & \textbf{Coopération avec la Chine} & \textbf{Coopération avec le Japon} \\
\hline
Principal instrument financier & Prêts concessionnels & Dons \\
\hline
Forum multilatéral de référence & FOCAC (Forum sur la coopération sino-africaine) & TICAD (Conférence internationale de Tokyo sur le développement de l'Afrique) \\
\hline
Agence d'exécution principale & Exim Bank of China & JICA \\
\hline
Secteurs emblématiques & Infrastructures lourdes (barrages, routes, ports, stades) & Santé, Éducation (sciences, mathématiques), eau, ponts \\
\hline
Type d'apport dominant & Infrastructures lourdes « clé en main » & Transfert de technologie et formation (assistance technique) \\
\hline
\end{tabularx}
\end{center}
\textbf{Commentaire :} le modèle chinois privilégie le \emph{hardware} (ouvrages financés par prêts à rembourser), le modèle japonais privilégie le \emph{software} (capital humain financé par dons, avec forte implication d'experts et de volontaires japonais).
\end{corrige}

\begin{corrige}[Exercice 8 — Amériques : USA et Brésil — Opportunités spécifiques]
\begin{enumerate}
\item \textbf{Trois produits camerounais éligibles à l'AGOA :} le \textbf{bois et ses dérivés}, les \textbf{produits agricoles transformés} (cacao transformé, café, huile de palme, fruits transformés) et le \textbf{textile-habillement}. (Le pétrole brut constitue l'essentiel des exportations vers les USA, mais l'objectif de l'AGOA est de diversifier vers les produits manufacturés.)
\item \textbf{Exemple de projet Power Africa au Cameroun :} l'appui au développement des \textbf{centrales solaires} (projets solaires de Guider et Maroua dans le Nord) et l'assistance technique à la \textbf{régulation du secteur électrique} (ARSEL) et à ENEO pour améliorer l'accès à l'électricité ; appui aux mini-réseaux ruraux.
\item \textbf{Coopération avec le Brésil :} l'institution de référence est l'\textbf{Embrapa} (Empresa Brasileira de Pesquisa Agropecuária), leader mondial de la recherche agricole tropicale. Exemples de transfert de technologie : amélioration des variétés de \textbf{coton} et de \textbf{manioc} (variétés à haut rendement, résistantes aux maladies), culture du \textbf{soja} en zone soudanienne, et coopération en \textbf{bioénergie} (éthanol à partir de canne à sucre).
\end{enumerate}
\textbf{Points de vigilance :} l'AGOA est un dispositif unilatéral américain révisé périodiquement (le Cameroun en a été partiellement suspendu en 2019 pour des questions de gouvernance) ; la coopération avec le Brésil relève de la coopération \cle{Sud-Sud}.
\end{corrige}

\begin{corrige}[Exercice 9 — Dissertation guidée : Chance ou dépendance ?]
\textbf{Éléments de corrigé (plan détaillé et contenu attendu).}

\textbf{Introduction (attendus) :} accroche d'actualité (sommets FOCAC, France-Afrique, adhésion du Cameroun au Commonwealth en 1995) ; définition de la coopération bilatérale (relations d'aide et de partenariat entre deux États) ; paradoxe : le Cameroun, riche en ressources, reste dépendant des financements extérieurs — la coopération est-elle un levier de développement ou un facteur de soumission ? Annonce du plan dialectique en deux parties.

\textbf{Première partie — Une chance pour le développement :}
\begin{itemize}
\item \emph{Argument 1 :} infrastructures structurantes impossibles à financer seul : barrage de Memve'ele et port de Kribi (Chine), autoroute Yaoundé-Nsimalen, appui français à Lom Pangar.
\item \emph{Argument 2 :} capital humain et services sociaux : hôpitaux et vaccination (Japon/JICA), lycées techniques et formation professionnelle (France/AFD), santé et éducation (USA/PEPFAR, bourses).
\item \emph{Argument 3 :} diversification des partenaires (France, Grande-Bretagne, Chine, Japon, USA, Brésil) qui accroît le pouvoir de négociation du Cameroun et réduit la dépendance à un seul partenaire.
\end{itemize}

\textbf{Deuxième partie — Les risques de dépendance :}
\begin{itemize}
\item \emph{Argument 1 :} endettement croissant (prêts Exim Bank liés aux grands ouvrages, clauses parfois opaques) et conditionnalités des bailleurs (programmes d'ajustement FMI/BM).
\item \emph{Argument 2 :} déséquilibre commercial : exportation de matières premières brutes (bois, pétrole, coton) contre importation de biens manufacturés ; faible transfert effectif de technologie ; main-d'œuvre parfois importée.
\item \emph{Argument 3 :} pressions sur les choix intérieurs : conditionnalités de gouvernance et de droits humains (suspension partielle de l'AGOA en 2019), alignements diplomatiques attendus aux enceintes internationales.
\end{itemize}

\textbf{Conclusion (attendus) :} bilan nuancé — la coopération est une chance \emph{à condition} d'une bonne gouvernance, de contrats gagnant-gagnant, de clauses de contenu local et d'une vision stratégique (SND30, Cameroun 2035) ; ouverture sur la coopération régionale (CEMAC, Union africaine, ZLECAf) comme voie de réduction de la dépendance.

\textbf{Barème indicatif (sur 20) :} Introduction (accroche, définitions, problématique, plan) : 4 pts ; Première partie (3 arguments étayés d'exemples camerounais précis) : 6 pts ; Deuxième partie (3 arguments étayés) : 6 pts ; Conclusion nuancée avec ouverture : 2 pts ; Qualité de la langue et de l'organisation : 2 pts.

\textbf{Points de vigilance :} exiger des exemples \emph{nommés} (Memve'ele, Kribi, C2D, JICA, AGOA) ; éviter le hors-sujet « coopération multilatérale » ; les deux parties doivent être équilibrées ; la conclusion doit trancher avec nuance, pas rester indécise.
\end{corrige}

\begin{corrige}[Exercice 10 — Étude de documents comparée : Chine vs France]
\begin{enumerate}
\item \textbf{Coopération sino-camerounaise (Document 1) :} elle se caractérise par de \textbf{grands ouvrages d'infrastructure} : énergie (barrage de Memve'ele), transports (autoroute Yaoundé-Nsimalen, port de Kribi), sport (stade d'Olembé) et santé (hôpital de référence). L'instrument financier dominant est le \textbf{prêt de l'Exim Bank of China} (environ 250 à 600 milliards FCFA par projet), donc un financement \emph{remboursable} qui accroît la dette, complété par quelques dons. Ce sont des projets « clé en main », visibles et rapides.
\item \textbf{Coopération franco-camerounaise (Document 2) :} elle cible les \textbf{services essentiels et le social} : eau et assainissement, éducation et formation professionnelle (PAREFOP, lycées techniques), énergie-climat. Les instruments sont surtout des \textbf{dons et prêts très concessionnels}, dont les \textbf{C2D} (contrats de désendettement-développement : annulation de dette convertie en projets). Elle accorde une place centrale à la \textbf{gouvernance} (appui budgétaire) et au \textbf{secteur privé} (Proparco, lignes de crédit aux PME).
\item \textbf{Comparaison des deux modèles :}
\begin{itemize}
\item \emph{Structure des échanges (Document 3) :} vers la Chine, le Cameroun exporte des produits \textbf{bruts} (bois brut, coton, pétrole, minerai : 700 Mrd FCFA) ; vers la France, des produits davantage \textbf{transformés} (bois transformé, cacao, café, aluminium : 350 Mrd FCFA). La relation sino-camerounaise renforce donc l'extraversion et l'exportation de matières premières.
\item \emph{Type d'endettement :} la Chine génère une \textbf{dette nouvelle} (prêts Exim Bank) ; la France privilégie la \textbf{réduction de la dette} (C2D) et les dons.
\item \emph{Transfert de technologie :} limité côté chinois (entreprises et main-d'œuvre souvent chinoises) ; plus présent côté français (formation professionnelle, appui institutionnel), mais avec des procédures plus lentes.
\end{itemize}
\item \textbf{Synthèse (éléments attendus) :} l'affirmation « la Chine construit le hardware, la France finance le software » est \textbf{largement vérifiée} : la Chine livre des infrastructures physiques massives (barrages, ports, routes) financées par prêts, tandis que la France investit dans l'humain, les institutions et la gouvernance (éducation, C2D, appui budgétaire). \textbf{Nuance :} la Chine finance aussi du « software » (hôpital de référence, bourses, formation) et la France du « hardware » (Lom Pangar, solaire, eau) ; par ailleurs, le hardware chinois crée une dette qui peut hypothéquer l'avenir, tandis que le software français entretient parfois une dépendance institutionnelle et monétaire (zone franc). L'idéal pour le Cameroun est d'articuler les deux : utiliser les infrastructures chinoises tout en consolidant le capital humain, en exigeant dans tous les cas des clauses de contenu local et de transfert de technologie.
\end{enumerate}
\textbf{Barème indicatif (sur 20) :} Q1 : 4 pts ; Q2 : 4 pts ; Q3 : 6 pts (2 pts par aspect) ; Q4 : 6 pts (thèse 2 pts, nuances 3 pts, expression 1 pt).
\textbf{Points de vigilance :} s'appuyer \emph{explicitement} sur les documents (citer les chiffres) ; ne pas opposer les partenaires de façon caricaturale ; la synthèse doit nuancer, c'est le cœur de l'épreuve.
\end{corrige}

\begin{corrige}[Exercice 11 — Sujet de débat argumenté : Choix stratégiques]
\textbf{Corrigé type (production attendue, 300-400 mots).}

\emph{Prise de position :} Pour son émergence à l'horizon 2035, le Cameroun ne doit pas choisir entre partenaires historiques et nouveaux partenaires : il doit pratiquer une \textbf{diversification stratégique} et une \textbf{hiérarchisation par secteur}, en tirant le meilleur de chaque relation.

\emph{Paragraphe 1 — L'apport irremplaçable des partenaires historiques :} la France et la Grande-Bretagne offrent un cadre juridique et monétaire stabilisé (zone franc, OHADA), un appui budgétaire décisif (C2D : annulation de dette convertie en projets sociaux), une formation d'excellence (lycées techniques, écoles d'ingénieurs et de santé publique) et une coopération de défense éprouvée. La langue et la culture administrative partagées réduisent les coûts de transaction. Se priver de ces appuis fragiliserait les services publics essentiels.

\emph{Paragraphe 2 — L'apport massif et rapide des nouveaux partenaires :} la Chine réalise en quelques années des infrastructures lourdes « clé en main » (barrage de Memve'ele, port de Kribi, autoroute Yaoundé-Nsimalen) à des coûts compétitifs et sans conditionnalités politiques explicites (principe de non-ingérence) ; le Brésil apporte son expertise agricole tropicale unique (Embrapa : coton, manioc, soja, bioénergie) dans le cadre de la coopération Sud-Sud. Ces partenaires répondent aux besoins urgents d'équipement que les partenaires traditionnels financent trop lentement.

\emph{Paragraphe 3 — Les risques à gérer :} côté nouveaux partenaires : dette parfois opaque, clauses secrètes, main-d'œuvre importée limitant l'emploi local ; côté partenaires historiques : conditionnalités politiques, lenteurs procédurales, désengagement progressif. La diversification impose donc une négociation rigoureuse de chaque contrat.

\emph{Conclusion prospective :} la doctrine idéale est le \textbf{partenariat gagnant-gagnant} : clauses de contenu local (emploi et formation des Camerounais), transfert effectif de technologie, alignement de tout accord sur la SND30 et la vision Cameroun 2035, et diplomatie économique active pour transformer localement les matières premières et conquérir le marché africain de la ZLECAf. Ainsi, les partenaires historiques et nouveaux deviennent complémentaires au service d'une souveraineté renforcée.

\textbf{Barème indicatif (sur 20) :} prise de position claire et nuancée : 4 pts ; qualité et exactitude des arguments et exemples (au moins 4 exemples nommés) : 10 pts ; conclusion prospective articulée à SND30/Cameroun 2035 : 4 pts ; langue et organisation : 2 pts.

\textbf{Points de vigilance :} éviter le choix binaire (c'est le piège du sujet) ; chaque argument doit être illustré par un exemple camerounais précis ; citer correctement les sigles (C2D, Embrapa, SND30, ZLECAf) ; rester dans le volume demandé.
\end{corrige}

\newpage

\coursec{Fiche bilan}

\begin{popfigure}
\begin{tikzpicture}[
  mindmap,
  grow cyclic,
  every node/.style={concept, execute at begin node=\hskip0pt},
  concept color=popblue,
  level 1/.append style={level distance=4.5cm, sibling angle=360/6},
  level 2/.append style={level distance=3cm, sibling angle=360/4},
  text=white,
  font=\small\bfseries
]
\node[concept, scale=1.2, align=center]{Coopération\\ Bilatérale\\ du Cameroun}
  child[concept color=poporange] {node[concept, align=center]{Généralités\\ (Définition,\\ Cadre juridique,\\ Principes)}}
  child[concept color=popgreen] {node[concept, align=center]{France\\ (Histoire,\\ Accords,\\ Domaines prioritaires)}}
  child[concept color=poppurple] {node[concept, align=center]{Grande-Bretagne\\ (Commonwealth,\\ Échanges,\\ Partenariats)}}
  child[concept color=poppink] {node[concept, align=center]{Chine\\ (FOCAC,\\ Infrastructures,\\ Investissements)}}
  child[concept color=popgold] {node[concept, align=center]{Japon\\ (TICAD,\\ Agriculture,\\ Formation)}}
  child[concept color=popdark] {node[concept, align=center]{USA \& Brésil\\ (AGOA, USAID,\\ Coopération Sud-Sud)}}
;
\end{tikzpicture}
\legende{Carte mentale : Les piliers de la coopération bilatérale du Cameroun (Terminale Technique, ECM).}
\end{popfigure}

\begin{bilan}
\courssub{Définitions et repères clés}
\begin{definition}[Coopération bilatérale]
Ensemble des relations officielles, juridiques, techniques, économiques et culturelles établies \textbf{entre deux États souverains} sur la base d'accords signés (traité, convention, protocole d'accord). Elle se distingue de la coopération multilatérale (ONU, UE, CEMAC, etc.).
\end{definition}

\begin{propriete}[Principes directeurs (Charte de l'ONU, Droit international)]
\begin{itemize}
\item Respect de la \cle{souveraineté} et de l'égalité souveraine.
\item \cle{Non-ingérence} dans les affaires intérieures.
\item \cle{Consentement mutuel} (pacta sunt servanda).
\item Recherche de l'\cle{intérêt mutuel} (gagnant-gagnant).
\end{itemize}
\end{propriete}

\courssub{Synthèse par partenaire (Essentiel à connaître)}

\begin{center}
\begin{tabularx}{\linewidth}{|l|X|X|}\hline
\rowcolor{popblueL}
\textbf{Partenaire} & \textbf{Cadre / Instruments principaux} & \textbf{Domaines phares (Exemples Bac)} \\ \hline
\textbf{France} & Accords de coopération (1960, 1974, 1994), C2D (Contrat de Désendettement-Développement), Sommets Afrique-France. & Éducation/Enseignement supérieur, Défense/Sécurité, Santé, Infrastructures, Francophonie, Appui budgétaire. \\ \hline
\textbf{Grande-Bretagne} & Commonwealth (adhésion 1995), Accords de commerce (UK-ESA EPA), British Council, Chevening. & Anglais/Éducation, Gouvernance/Droits de l'homme, Commerce/Investissements, Climat/Environnement. \\ \hline
\textbf{Chine} & FOCAC (Forum sur la Coopération Sino-Africaine), Initiative « Ceinture et Route » (BRI), Accords économiques et techniques. & Infrastructures lourdes (barrages, routes, stades), Santé (hôpitaux, équipes médicales), Agriculture/Foresterie, Numérique, Formation professionnelle. \\ \hline
\textbf{Japon} & TICAD (Conférence internationale de Tokyo sur le développement de l'Afrique), JICA, Volontaires (JOCV). & Développement rural/Agriculture (riziculture), Santé (lutte contre paludisme), Éducation technique/Formation, Infrastructures de base, Paix/Sécurité. \\ \hline
\textbf{USA} & AGOA (Loi sur la croissance et les opportunités en Afrique), USAID, MCC (Millennium Challenge Corp.), Peace Corps. & Santé (VIH/SIDA, Palu, Santé maternelle), Gouvernance/Démocratie, Sécurité régionale (Bassin du Lac Tchad), Éducation/Entrepreneuriat (YALI). \\ \hline
\textbf{Brésil} & Coopération Sud-Sud, ABC (Agence brésilienne de coopération), Accords techniques. & Agriculture tropicale (Embrapa : coton, soja, manioc), Énergies renouvelables (éthanol, solaire), Santé (médicaments génériques), Formation technique professionnelle. \\ \hline
\end{tabularx}
\end{center}

\courssub{Méthode express : Analyser un accord de coopération (Type Bac)}
\begin{methode}[Démarche en 4 étapes]
\begin{enumerate}
\item \textbf{Identifier} : Nature de l'acte (Traité, Accord-cadre, Protocole), Parties signataires, Date, Durée.
\item \textbf{Qualifier} : Domaine(s) concerné(s) (éducation, santé, infrastructure, défense, commerce...), Type d'appui (don, prêt concessionnel, expertise, formation, transfert de technologie).
\item \textbf{Évaluer} : Retombées concrètes pour le Cameroun (emplois, infrastructures réalisées, capacités renforcées, balance commerciale) ET contreparties (alignement diplomatique, ouverture marchés, clauses conditionnelles).
\item \textbf{Conclure} : Apport au \cle{développement national} (PND 2020-2030, ODD) et respect de la \cle{souveraineté}.
\end{enumerate}
\end{methode}

\courssub{Bilan critique : Avantages, Limites, Enjeux}
\begin{itemize}
\item \textbf{Avantages} : Ressources financières/expertises externes ; Transfert de compétences/technologies ; Appui aux secteurs sociaux (santé, éducation) ; Diversification des partenariats (réduction dépendance unique) ; Rayonnement diplomatique.
\item \textbf{Limites} : Conditionnalités politiques/économiques parfois lourdes ; Risque d'\cle{endettement} (prêts non concessionnels) ; Asymétrie de pouvoir (négociation inégale) ; Projets « clés en main » peu intégrés au tissu local ; Dilution des priorités nationales au profit des agendas donateurs.
\item \textbf{Enjeux pour le Cameroun} : Maîtriser la \cle{planification} (alignement sur PND/ODD) ; Renforcer la \cle{négociation} (clauses de transfert de technologie, main-d'œuvre locale) ; Améliorer l'\cle{absorption} des fonds (taux d'exécution) ; Promouvoir la \cle{coopération triangulaire} et \cle{Sud-Sud} (Brésil, Chine, Inde, Maroc) pour plus d'autonomie.
\end{itemize}
\end{bilan}

\begin{aretenir}[Checklist avant le Bac]
\begin{itemize}
\item $\square$ Définir la coopération bilatérale et la distinguer de la multilatérale.
\item $\square$ Citer 3 principes juridiques qui la régissent (souveraineté, non-ingérence, consentement, intérêt mutuel).
\item $\square$ Pour la \textbf{France} : Citer le C2D et 2 domaines prioritaires (ex: éducation, défense).
\item $\square$ Pour la \textbf{Grande-Bretagne} : Citer le Commonwealth et l'APE (UK-ESA EPA) + 1 domaine (ex: gouvernance).
\item $\square$ Pour la \textbf{Chine} : Citer le FOCAC et la « Ceinture et Route » + 1 réalisation type (infrastructure lourde).
\item $\square$ Pour le \textbf{Japon} : Citer la TICAD et la JICA + 1 spécificité (ex: agriculture rizicole, volontariat).
\item $\square$ Pour les \textbf{USA} : Citer l'AGOA et l'USAID + 1 domaine (santé, gouvernance, sécurité).
\item $\square$ Pour le \textbf{Brésil} : Identifier la « Coopération Sud-Sud » et l'expertise agricole (Embrapa).
\item $\square$ Savoir lister 2 avantages et 2 limites de la coopération bilatérale pour le Cameroun.
\item $\square$ Savoir expliquer l'enjeu de l'\textbf{appropriation nationale} (alignement PND/ODD).
\item $\square$ Être capable d'analyser un extrait d'accord (identifier parties, objet, type d'appui, conditionnalité).
\end{itemize}
\end{aretenir}

\findecours

\end{document}
